% concatenated from asymptotics_of_hLc_250916_arxiv_version.tex
%\documentclass{article}
%\usepackage{graphicx} % Required for inserting images
\documentclass[12pt]%{article}
{amsart}
\usepackage{fancyhdr,url}

\usepackage{amsmath, amsthm, amssymb, tikz, hyperref}
\usepackage{amsfonts}
% %\usepackage{fullpage}
% \usepackage[margin=2cm]{geometry}
% %\topmargin=-0.2cm
% \headheight=15pt
% \linespread{1.3}
\usepackage[a4paper,
%top=2cm,bottom=2cm,left=2.5cm,right=2.5cm,
margin=2.5cm,
marginparwidth=1.75cm]{geometry}

\newcommand{\bbn}{\mathbb{N}}
\newcommand{\bbz}{\mathbb{Z}}
\newcommand{\bbr}{\mathbb{R}}
\newcommand{\bbq}{\mathbb{Q}}
\newcommand{\FF}{\mathbb{F}}
\renewcommand\S{{\mathcal S}}
\newcommand\T{{\mathcal T}}
\newcommand\X{{\mathcal X}}
\newcommand{\EEE}{{\mathcal{E}}}

\DeclareMathOperator\Des{Des}
\DeclareMathOperator\des{des}
\DeclareMathOperator\CDes{CDes}
\DeclareMathOperator\cdes{cdes}
\DeclareMathOperator\maj{maj} \DeclareMathOperator\Inv{Inv}
\DeclareMathOperator\inv{inv} \DeclareMathOperator\stat{stat}
\DeclareMathOperator\Reg{Reg}
\DeclareMathOperator\Conj{Conj}

\newcommand{\cDes}{{\operatorname{cDes}}}

\newcommand{\Hilb}{{\operatorname{Hilb}}}

\newcommand{\Q}{{\mathcal Q}}
\newcommand{\TTT}{{\mathcal{T}}}

\newcommand{\la}{{\lambda}}
\newcommand{\al}{{\alpha}}
\newcommand{\Si}{{\Sigma}}
\newcommand{\si}{{\sigma}}
\newcommand{\om}{{\omega}}
\newcommand{\Lie}{\operatorname{Lie}}
\newcommand{\Ind}{\operatorname{Ind}}
\newcommand{\Pl}{M_{\rm Planch}}
\newcommand{\Prob}{{\rm Prob}}
\newcommand{\conj}{M_{\rm conj}}
\newcommand{\PD}{\operatorname{PD}}
\def\red#1{{\color{red}{#1}}}
\def\magenta#1{{\color{magenta}{#1}}}

\theoremstyle{plain}
\newtheorem{theorem}{Theorem}[section]
\newtheorem{proposition}[theorem]{Proposition}
\newtheorem{lemma}[theorem]{Lemma}
\newtheorem{corollary}[theorem]{Corollary}
\newtheorem{conjecture}[theorem]{Conjecture}
\newtheorem{question}[theorem]{Question}
\newtheorem{problem}[theorem]{Problem}

\theoremstyle{definition}
\newtheorem{defn}[theorem]{Definition}
\newtheorem{definition}[theorem]{Definition}
%\newtheorem{example}[theorem]{Example}
\newtheorem*{ex}{Example}
%\newenvironment{ex}[1][Example.]{\begin{trivlist}
%\item[\hskip \labelsep {\bfseries
%#1}]}{\end{trivlist}}

\theoremstyle{remark}
\newtheorem{notation}[theorem]{Notation}
\newtheorem{remark}[theorem]{Remark}
\newtheorem{observation}[theorem]{Observation}



\newcommand\coolrightbrace[2]{%
	\left.\vphantom{\begin{matrix} #2 \end{matrix}}\right\}#1}

%\newcommand{\Q}{\mathcal Q}
%\newcommand{\f}{\mathcal F}
\newcommand{\symm}{S}
\newcommand{\C}{C}%{\mathcal C}
\newcommand{\A}{\mathcal A}

%\newcommand{\sign}{{\rm sign}}
\DeclareMathOperator\sign{sign}
\newcommand{\Irr}{{\rm Irr}}

\def\n{\mathcal{N}}
\def\m#1{\mathcal{N}_{#1}}
\def\p{(\textnormal{mod}\ p)}
\def\x#1{(\textnormal{mod}\ x^{#1})}
\def\F{\mathbb{F}}%_{p}}
\def\b{\bar{n}}

\DeclareMathOperator\lcm{lcm}
\newcommand{\id}{{\operatorname{id}}}

\DeclareMathOperator\sgn{sgn}
\def\CC{{\mathbb C}}
\def\KK{{\mathbb K}}
\def\LL{{\mathbb L}}
\def\NN{{\mathbb N}}
\def\PP{{\mathbb P}}
\def\QQ{{\mathbb Q}}
\def\ZZ{{\mathbb Z}}
\def\RR{{\mathbb R}}
\def\kk{{\mathbf k}}

\newcommand{\blambda}{{\underline\lambda}}%{{\mathbf{\lambda}}}
\newcommand{\bmu}{{\underline\mu}}
\newcommand{\br}{{\underline r}}
\newcommand{\bs}{{\underline s}}
\newcommand{\bt}{{\underline t}}
\newcommand{\bc}{{\underline c}}


\newcommand{\rank}{{\operatorname{rank}}}

\newcommand{\Par}{{\operatorname{Par}}}
\newcommand{\ve}{\varepsilon}

\newcommand{\cyc}{{\operatorname{cyc}}}

\newcommand\comp[2]{\alpha{(#2,#1)}}
\newcommand\ccomp[2]{\alpha^\cyc{(#2,#1)}}

\newcommand{\cs}{\tilde s}

\newcommand{\ch}{{\operatorname{ch}}}
\newcommand{\fix}{{\operatorname{fix}}}
%\newcommand{\sign}{\operatorname{\sign}}

\newcommand{\roots}{\varphi} 
\newcommand{\sroots}{{\overline{\varphi}}} 
\newcommand{\spsi}{{\overline{\psi}}} 
\newcommand{\sk}{{k,-}}%{{\overline{k}}} 

\newcommand{\SYT}{{\operatorname{SYT}}}

\newcommand{\todo}[1]{\vspace{5 mm}\par \noindent
	\marginpar{\textsc{ToDo}} \framebox{\begin{minipage}[c]{0.95
				\textwidth}
			#1 \end{minipage}}\vspace{5 mm}\par}

\title{Asymptotics of higher Lie characters}

% \author{Ron M.\ Adin, Yuval Roichman, and Natalia Tsilevich (not for distribution)}

\author[R.\ M.\ Adin]{Ron M.\ Adin}
\address{Department of Mathematics, Bar-Ilan University, Ramat-Gan 52900, Israel}
\email{radin@math.biu.ac.il}

\author[Y.\ Roichman]{Yuval Roichman}
\address{Department of Mathematics, Bar-Ilan University, Ramat-Gan 52900, Israel}
\email{yuvalr@math.biu.ac.il}

\author[N.\ Tsilevich]{Natalia Tsilevich}
\address{Department of Applied Mathematics, Braude College of Engineering, Karmiel 2161002, Israel} 
\email{natalia.tsilevich@gmail.com}


\date{September 16, 2025}

\thanks{Part of this research was done while N.~Tsilevich was hosted by  Bar-Ilan University in the framework of the Shapira program.}

\keywords{Asymptotic representation theory, symmetric group, higher Lie characters, conjugacy representation,  Thrall's problem}

%\subjclass{20C30}

\begin{document}

\begin{abstract} 
    Higher Lie characters form a distinguished family of symmetric group characters, which appear in many areas of algebra and combinatorics. An old open problem of Thrall is to decompose them into irreducibles. 
    We propose a novel asymptotic approach to this problem, showing that many families of 
    higher Lie characters 
    tend to be proportional to the regular character. 
    In particular, a random higher Lie character 
    tends, in probability, to be proportional to the regular character.
\end{abstract}

\maketitle

\tableofcontents

\section{Introduction}

\subsection{Higher Lie characters and Thrall's problem}

Higher Lie characters form a distinguished family of symmetric group characters. 
Their study can be traced back to Schur~\cite{Schur}, and then to  papers in Lie theory by Thrall~\cite{Thrall} 
%, Brandt~\cite{Brandt}, Wever~\cite{Wever}
and others. 
The important special case of Lie characters was thoroughly studied by Klyachko~\cite{Klyachko}, Kra\'skiewicz and Weyman~\cite{KraskiewiczWeyman}, and Garsia~\cite{Garsia}, who pointed out 
their significant role in combinatorics. 
A new era in the study of higher Lie characters began when Gessel and Reutenauer~\cite{GR93} presented them by means of quasisymmetric functions, leading to surprising applications in various areas of algebra and combinatorics: Lyndon words~\cite[Ch.~7]{Reutenauer}, \cite[Exercise 7.89]{Stanley_ECII}, root enumeration~\cite[Ex.~7.69]{Stanley_ECII}, construction of Gelfand models~\cite{Inglis}, permutation statistics 
%(interplay of fixed points, and more generally cycle structure with descent statistics), 
and card shuffling~\cite{Saliola}. 
%Intimate relations
Strong connections were found to the Whitney homology of the partition lattice~\cite{Stanley1982},  
the cohomology of configuration spaces~\cite{Hersh}, 
and Varchenko--Gelfand and Orlik--Solomon algebras~\cite{Almousa}. 
%by Stanley~\cite{Stanley?}, Hersh and Reiner~\cite{Hersh}, Reiner and Sundaram~\cite{Almousa}. 
%{ReinerSundaram}.

% \todo{Short pa regarding the significance of higher Lie characters in various fields:  
% Lyndon words, Lie theory (Klyachko, ~\cite{KraskiewiczWeyman}, Garsia), 
% construction of Gelfand models, root enumeration in finite groups, theory of quasisymmetric functions and permutation statistics~\cite{GR93}, 
% hyperplane arrangements (Orlik-Solomon
% Gelfand-Varchenko algebras), (graph colorings ??).

% NT: Also mention the construction of higher Lie's from the action of $\S_n$ in the free Lie algebra as an explanation for the term.
% }

%\subsection{Thrall's problem (Old version, to be updated)}




%~\cite{schur}. % and Thrall~\cite{thrall}. 
% An old problem of Thrall~\cite{Thrall} 
% asks for an explicit combinatorial interpretation of the multiplicities of irreducible characters in each higher Lie character;
% see also~\cite[Exercise 7.89(i)]{Stanley_ECII}. 
For a partition $\lambda\vdash n$,
denote by $\chi^\lambda$ and by $\psi^\lambda$ the irreducible character and the higher Lie character indexed by $\lambda$, respectively. 

\begin{problem} {\rm (Thrall~\cite{Thrall})}
Find a combinatorial interpretation of the inner products 
$c_\nu^{\lambda}:=\langle \psi^\lambda,\chi^\nu\rangle$.
\end{problem}

%See also~\cite[Exercise 7.89(i)]{Stanley_ECII}. 
This 83 years old open problem of Thrall 
%(see also~\cite[Exercise 7.89(i)]{Stanley_ECII})
%is still open, 
%\red{
%Thrall’s problem remains open in general, and 
has received significant attention over the years; 
%in the literature since its formulation; 
see, e.g., \cite{Brandt, Klyachko, KraskiewiczWeyman, Garsia, GR93, Reutenauer, Sundaram94, Schocker}. 
%}   
%[Bra44, Kly74, BBG90, Gar90, GR93, Sun94, KW01, Sch03, AS18].} 
%For this problem, %Only partial results are known. 
%For more recent discussions %of Thrall's problem 
%For 
%%a detailed introduction see~\cite{Reutenauer, Garsia}; for 
%a lucid survey see~\cite{Reiner-Banff}.} 
Solutions are known for distinguished families of partitions.  
The partitions $\lambda = (2^k 1^{n-2k})$ 
%, corresponding to involutions with a given number of fixed points, 
yield $c_\nu^\lambda \in \{0,1\}$, an important multiplicity-free property that was applied to construct Gelfand models~\cite{Inglis}. 
The case $\lambda = (n)$ was solved by Kra\'{s}kiewicz and Weyman~\cite{KraskiewiczWeyman}, using the major index statistic; see Theorem~\ref{thm:KW} below. 
The best general result so far is Schocker's expansion~\cite[Theorem~3.1]{Schocker}, which 
however involves signs and rational coefficients.
%\red{
Thrall's problem was approached from different viewpoints in the last decade; 
see, e.g., \cite{Hersh, Reiner-Banff, AS_Thrall, Sundaram18, 
AHR, Lim, Pak, Armon, Amendola, Commins}.  
%}

%They showed that
%%, for $\lambda =(n)$,
%$Q(A_{(n)}) = \sum_{T} s_{\lambda(T)}$,
%where $\lambda(T)$ is the shape of $T$ and the summation is over all standard Young tableaux $T$ with $n$ cells having {\em major index} $\maj(T) :=
%\sum_{i \in \Des(T)}i$ congruent to $1$ modulo $n$. 
%One can then fairly easily show that, for a general partition $\lambda=1^{m_1}	2^{m_2} 3^{m_3} \cdots$,
%	\[
%	Q(A_\lambda)= \prod_{i \geq 1} h_{m_i}[Q_{A_{(i)}}],
%	\]
%	where $h_m$ is a complete homogeneous symmetric function and
%	$f[g]$ denotes plethystic composition of the symmetric function
%	$g$ into the symmetric function $f$. In particular, using the
%	Littlewood-Richardson rule for the product $s_\mu s_\nu$, one can
%	reduce Thrall's problem to the special case of a rectangle
%	$\lambda=b^a$, where $Q(A_{b^a}) = h_{a}[Q_{A_{(b)}}]$; but due to
%	the plethysm, this doesn't help so much. 
%A bijective proof of this classical result was %recently 
%given by Ahlbach and Swanson~\cite{AS_Thrall}. 
%Their proof applies a refined cyclic sieving phenomenon (CSP) on words, combined with the RSK correspondence. A possible application of this approach to the original Thrall's problem could be a great breakthrough.

%\cite{Sundaram} 
%{\bf update references ??}.  

%\todo{RA: Write the most relevant references here.}
% \red{
In this paper we propose a novel asymptotic approach to Thrall's problem. 
%\red{%Namely, 
Instead of focusing on the multiplicities $\langle \psi^\la,\chi^\nu \rangle$ for fixed finite partitions $\la$ and $\nu$, we study their asymptotic behavior as both partitions grow. 
% infinitely. 
It turns out that, for many families of partitions $\la$,  
%for $n$ large enough, 
the higher Lie character $\psi^\la$ 
% is 
% close 
%\red{nearly proportional}
tends to be proportional 
to the regular character.
% ; see Subsection~\ref{sec:main_results} for details. 
%The only two previously known results in this direction are described in the next subsection.
% } 
%}

%\todo{RA: Use stronger words.}

% \todo{Old version:

% \begin{conjecture}
%     For most partitions $\lambda \vdash n$, the higher Lie character $\psi^\lambda$ is asymptotically proportional to the regular character.\\
%     More precisely, for every sequence of partitions  $\mu^n \vdash n$  with no bounded parts, the higher Lie characters $\psi^{\mu^n}$ are asymptotically norm equal to the character of the regular representation.
% \end{conjecture}

% This conjecture has been verified for several families of partitions. 
% }


\subsection{Asymptotic proportionality of characters}\label{sec:tendtoregular}

% Our main interest in this paper is the following question: {\it for which $\la$ is the higher Lie character $\psi^{\la}$ 
% % asymptotically proportional 
% asymptotically close
% to the regular character $\Reg$?}  

There are various ways to measure how asymptotically close two character sequences $\rho = (\rho_n)$ and $\rho' = (\rho'_n)$ are.
% We choose 
Our choice is
to compare multiplicities 
%in $\psi^{\la}$ and $\Reg$ of random 
of irreducible characters,
$\langle \rho_n,\chi^{\nu^{(n)}} \rangle$ and $\langle \rho'_n,\chi^{\nu^{(n)}} \rangle$,
where the sequence of diagrams $\nu = (\nu^{(n)})$ is typical with respect to the Plancherel measure; see Subsection~\ref{sec:Plancherel} for more details about this measure. 
% As explained in Subsection~\ref{sec:Plancherel}, the most natural choice for the distribution of a random character is the Plancherel distribution. 
%
%Let us make the question precise.
%
The choice of Plancherel-typical ``test characters'' $\chi^\nu$ is based on the following consideration.
% (for more details, see Subsection~\ref{sec:Plancherel}).
The Plancherel probability of a 
% Young diagram 
diagram $\nu$ is proportional to the dimension of the isotypic component of the regular representation $\Reg_n$ corresponding to $\chi^\nu$, which is $(\dim \chi^\nu)^2$. Thus, from a representation-theoretic point of view, it is natural to assume that a ``randomly chosen'' irreducible character $\chi^{\nu}$ is distributed according to the Plancherel measure.
We therefore introduce the following definition.

\begin{definition}
    Two character sequences $\rho = (\rho_n)$ and $\rho' = (\rho'_n)$, where $\rho_n$ and $\rho'_n$ are characters of $\S_n$, are said to be {\em asymptotically proportional} if 
    \begin{equation*}
        \frac{\langle \rho_n,\chi^{\nu^{(n)}} \rangle}{\dim\rho_n}
        \sim \frac{\langle \rho'_n,\chi^{\nu^{(n)}} \rangle}{\dim\rho'_n}
        \quad \text{as} \quad n\to\infty
    \end{equation*}
    for a Plancherel-typical sequence of Young diagrams $\nu=(\nu^{(1)},\nu^{(2)},\ldots)$.
\end{definition}

% \todo{RA: 
%     Why regular?
% }
In particular, if $\rho'_n = \Reg_n$, the regular character of $\S_n$, then $\dim\Reg_n = n!$ and $\langle \Reg_n,\chi^{\nu^{(n)}}\rangle = f^{\nu^{(n)}}$.

\begin{definition}\label{def:regularity}
    A sequence of characters $\rho = (\rho_n)$ {\em tends to be regular} if it satisfies
    \begin{equation}\label{tends_to_regular}
        \frac{\langle \rho_n,\chi^{\nu^{(n)}} \rangle}{\dim\rho_n}
        \sim \frac{f^{\nu^{(n)}}}{n!}
        \quad \text{as} \quad n\to\infty
    \end{equation}
    for Plancherel-typical $\nu = (\nu^{(n)})$.   
  
\end{definition}



%Namely,

%\begin{definition}
%A {\em random higher Lie character of size $n$} $\psi^\la$ is indexed 
%by the cycle type of a uniformly chosen random permutation in $\S_n$. 
%\end{definition}


%\begin{definition}\label{def:regularity}
%%(Sequences of) diagrams $\la$ satisfying~\eqref{claim}, or equivalently  \eqref{symmfunc}, for typical $\nu$ are said to be {\em regular}.   
%A sequence of higher Lie characters $\psi^\la$ {\em tends to be regular} if it 
%satisfies~\eqref{claim} 
%%, or equivalently  \eqref{symmfunc}, 
%for typical $\nu$.   
%\end{definition}

% \todo{RA: Should we add the following?

% Note that,
% since $\nu_1 \cdot \nu'_1 \ge n$, 
% if $\max\{\nu_1, \nu'_1\} \le C \cdot \sqrt n$ then 
% $\min\{\nu_1, \nu'_1\} \ge C^{-1} \cdot \sqrt n$.
% }


\subsection{Existing asymptotic results}

% \todo{Motivating subsection (existing results regarding asymptotics of sum of hLc)\\
% To be rewritten, to fit as a part of introduction}

%\todo{
%It is well known that the conjugacy character of $\S_n$ is asymptotically the same as the regular character~\cite{AF, R97}. } 

%\section{Desired applications}

%\section{Old Section: Derangements}
  
%\subsection{The result}

%\todo{YR: It is desired to have a stronger version of Theorem~\ref{thm:main}
%which holds, for sufficiently large finite $n$.
%We would like the strengthened version to imply the below analogous results.}

%There is a number of papers dealing with the asymptotic properties of conjugacy and higher Lie characters. 

 
%It is well known that 
The sum of all higher Lie characters of $\S_n$, each with multiplicity $1$, is the regular character $\Reg_n$; see Theorem~\ref{t:sum_of_hLc} for a formulation in terms of symmetric functions. 
% ; see Theorem~\ref{t:sum_of_hLc} below for a formulation in terms of symmetric functions. 
Coarse versions of Thrall's problem (see~\cite{Reiner-Banff}) aim to describe the decomposition into irreducibles of other (large) sums of higher Lie characters.
For example, the {\em derangement character} ${\mathcal D}_n$ of the symmetric group $\S_n$  is the sum of higher Lie characters over all partitions with no parts of size~$1$. Thus, the dimension of ${\mathcal D}_n$ is equal to the number $d_n$ of derangements (fixed-point-free permutations) in $\S_n$. For more about the derangement character see, e.g., \cite{DesarmenienWachs, ReinerWebb, AAER}.
Following a combinatorial interpretation by D\'esarm\'enien and Wachs~\cite{DesarmenienWachs}, 
Regev~\cite{Regev99} presented the following remarkable asymptotic result. 
%solution of coarse Thrall's. 

% \todo{
% D\'esarm\'enien and Wachs~\cite{DesarmenienWachs} resolved a coarse version of Thrall's problem, considering 
% %the sum of higher Lie characters over all partitions without a singleton part (corresponding to derangements); see also~\cite{ReinerWebb}. 
% the {\em derangement character} $D_n$ of the symmetric group $\S_n$, which is the direct sum of higher Lie characters over all partitions with no parts of size~$1$. Thus, the dimension of $D_n$ is equal to the number $d_n$ of derangements (fixed-point-free permutations) in $\S_n$.
% %(see~\cite{DesarmenienWachs, ReinerWebb}). 
% Combining \cite[Theorem~8.1]{GR93}, \cite[Propositions~2.1 and~2.3]{ReinerWebb}, and \cite[Theorem~7.2]{AAER} (which, in turn, is based on \cite{DesarmenienWachs} and \cite{ReinerWebb}), one can see that that the multiplicity $\langle D_n,\chi^\nu\rangle$ of an irreducible character $\chi^\nu$ in $D_n$ is exactly the coefficient $r_\nu$ studied by Regev~\cite{Regev99}. Thus, since for a typical diagram $\nu$ its rows and columns grow sublinearly, Theorem~0.3 in \cite{Regev99} says, in particular, the following.
% }

\begin{theorem}\label{thm:Regev}{\rm (Regev~\cite{Regev99})}
    For Plancherel-typical $\nu$,  
    \[
        \langle {\mathcal D}_n,\chi^{\nu} \rangle
        \sim \frac{f^{\nu}}{e} 
        \quad \text{as} \quad n\to \infty.
    \]
\end{theorem}

% For the meaning of Plancherel-typical shapes see Subsection~\ref{sec:Plancherel} below.
Note that $\lim_{n \to \infty} d_n / n! = 1/e$, so that Regev's result can be written as
\[
    \frac{\langle {\mathcal D}_n,\chi^{\nu} \rangle}{\dim {\mathcal D}_n}
    \sim \frac{f^{\nu}}{n!}
    \quad \text{as} \quad n\to\infty
\]
for Plancherel-typical $\nu$.
In other words, ${\mathcal D}_n$ tends to be regular in the sense of Definition~\ref{def:regularity}.

% and more??.

% \todo{
% It is well known that 
% $$\frac1e=\lim\limits_{n\to\infty}\frac{d_n}{n!}
% $$ 
% is the asymptotic fraction of derangements among all permutations in $\S_n$, in the same way as $1/z_\la$ is the fraction of permutations with cycle structure $\la$ among all permutations. Thus, Theorem~\ref{thm:Regev} is an exact analog of our Theorem~\ref{thm:main} for the derangement character (which is a sum of higher Lie characters) instead of a random individual higher Lie character. Also observe that the sum of the higher Lie characters $\psi^\la$ over all $\la\vdash n$ is exactly the regular character of $\S_n$. 
% }

\medskip

The action of $\S_n$ on itself by conjugation is closely related to higher Lie characters, as explained in Subsection~\ref{sec:conjugacy} below; denote by ${\mathcal C}_n$ the correspoding character.
%be the conjugacy character of $\S_n$ on itself. 
The asymptotics of ${\mathcal C}_n$
%the conjugacy character on the whole group $\S_n$ 
was first studied in~\cite{AF}, where it was shown that it tends, in a certain sense, to the regular character. 
A well known problem (see, e.g., \cite{Heide}) is to decompose 
%the conjugacy character 
${\mathcal C}_n$ into irreducible characters. 
The following result follows from~\cite[Theorem 2.1]{R97}. 

\begin{theorem}\label{thm:Roichman}
    For Plancherel-typical $\nu$, 
\[
    \frac{\langle {\mathcal C}_n,\chi^{\nu} \rangle}{\dim {\mathcal C}_n}
    \sim \frac{f^{\nu}}{n!}
    \quad \text{as} \quad n\to\infty.
\]
    % \[
    %     \langle C_n,\chi^{\nu} \rangle
    %     \sim \frac{f^{\nu}}{1} 
    %     \quad \text{as} \quad n\to \infty.
    % \]
\end{theorem}
Thus, ${\mathcal C}_n$ also tends to be regular in the sense of Definition~\ref{def:regularity}.

% \todo{
% On the other hand, it is well known (\cite{AF, R97}) that the conjugacy character  of $\S_n$ (which is the sum of $\phi^\lambda$ over all $\lambda\vdash n$) is asymptotically close to the regular character. 
% %We would like the strengthened version to imply these results.
% In particular, it follows from the proofs in~\cite{R97}, that the normalized multiplicity of a typical  irreducible character $\chi^\nu$ in the conjugacy character   tends to $1$. Thus, this result is again an analog of our Theorem~\ref{thm:main} for the conjugacy character. 
% }


% \todo{To be moved to final section??

% Looking at the collection of these results suggests the following question.

% \begin{question}
%     Consider the sum of $\phi^\la$ over all partitions $\la\vdash n$ with no rows of length~$1$.
%     Does the normalized multiplicity of a typical Plancherel irreducible character in it tend to $\frac{1}{e}$?   
% \end{question}
% }

%\begin{question}
% \todo{Is the condition that the rows and columns grow sublinearly a necessary condition?}
%\end{question}


\subsection{Main results}\label{sec:main_results}

In this paper we consider
%/confront
Thrall's original problem (rather than its coarse versions) in an asymptotic setting; namely, we study the multiplicities of irreducibles in single higher Lie characters (rather than sums of characters), indexed by diagrams of growing sizes.  

Our first main result deals with sequences of higher Lie characters indexed by rectangular diagrams. 
Note that Thrall's problem may be reduced, via the Littlewood--Richardson rule, to the special case of a rectangular diagram; see, e.g., \cite{Reiner-Banff} and references therein.  

\begin{theorem}\label{main:rectangular}
Any sequence of higher Lie characters, indexed by rectangular diagrams $\la = (m^k)\vdash n$ with
$m \to \infty$ and $k = o(m)$, tends to be regular. 
% Namely,
% In particular,  
% for Plancherel-typical $\nu$, 
% \[
%     \frac{\langle \psi^{(m^k)},\chi^{\nu} \rangle}{\dim \psi^{(m^k)}}
%     \sim \frac{f^{\nu}}{n!}
%     \quad \text{as} \quad n\to\infty. 
% \]
\end{theorem}

For a precise definition of the notion {\em tends to be regular} see Definition~\ref{def:regularity} below. For  
a stronger version of this result see Theorem~\ref{thm:rectangular} below.

%\medskip

% \todo{
% Our main result is that a random higher Lie character tends in probability to be regular. More precisely, 
% let $\la(\pi)$ be the cycle structure of a permutation $\pi\in \S_n$. 
% }

The second main result of this paper deals with the asymptotics of a random higher Lie character. 
% For the exact distribution, see Definition~\ref{def:random_hLc}. 
%Subsection~\ref{sec:???}.
%In Section~\ref{sec:asymp} we consider the asymptotic behavior of a randomly chosen higher Lie character.
%Let us now turn to a ``randomly chosen'' higher Lie character. 
%The sum of all higher Lie characters of $\S_n$, each with multiplicity $1$, is the regular representation $\Reg_n$ (see Theorem~\ref{t:sum_of_hLc} for a formulation in terms of symmetric functions), 
As mentioned above,
the sum of higher Lie characters $\psi^\la$ over all partitions $\la \vdash n$ is the regular character $\Reg_n$. Also, $\dim\psi^\la=|C_\la|$, where $C_\la$ is the conjugacy class consisting of all permutations in $\S_n$ with cycle type $\la$. Thus, exactly the same logic as in the choice of Plancherel measure for irreducible characters $\chi^\nu$ suggests that 
%the following: 
% \begin{definition}\label{def:random_hLc}
%that we should take 
a random higher Lie character $\psi^\la$ ($\la \vdash n$) be chosen with probability  $|C_\la|/n!$. 
%to be proportional to $|C_\la|$. 
% \end{definition}

%Let $\la(\si)$ be the cycle structure of a permutation $\si \in \S_n$. 

% \todo{Very important: Explain the dichotomy: why do we use Plancherel measure for irreducible constituents, versus cycle structure of a random permutation for a ``typical" higher Lie characters}

\begin{theorem}\label{thm:main} 
     % If $\la_n$ is a random Lie character,  then $\psi^{\la_n}$ 
     A random higher Lie character $\psi^{\la_n}$, where $\la_n \vdash n \to \infty$, 
    tends in probability to be regular.
\end{theorem}

For stronger quantitative versions, see Theorems~\ref{thm:main1} and~\ref{thm:main_general} below. 

%\medskip

To prove Theorem~\ref{thm:main} we analyze, first, diagrams with distinct row lengths, and show that under mild conditions they tend to be regular (Theorem~\ref{thm:distinct} below).  
Another important tool in the proof of Theorem~\ref{thm:main} is the Gluing Lemma (Lemma~\ref{l:gluing} below),
which allows us to glue relatively small diagrams to 
large diagrams 
and preserve the tending-to-regularity property. 
An immediate consequence of the Gluing Lemma is a sufficient condition for higher Lie characters of 
hook shapes to tend to be regular (Corollary~\ref{cor:hooks}).   


\begin{remark}
    Theorems~\ref{main:rectangular} and~\ref{thm:main}, as well as all other results in this paper, hold also when higher Lie characters are replaced by any other characters induced from linear characters of centralizers, in particular for the characters of the conjugacy action of $\S_n$ on its conjugacy classes;
    see Theorem~\ref{thm:conjugacy} and Corollary~\ref{cor:conjugacy} below. 
\end{remark}

%\todo{RA:
For the benefit of the reader, here is a list of other results in this paper: 
Theorem~\ref{prop:Lie} (one-row diagrams), 
Theorem~\ref{thm:distinct} (distinct row lengths),
Lemma~\ref{l:gluing} (Gluing Lemma),
Corollary~\ref{cor:hooks} (hook shapes),
Corollary~\ref{thm:main_cor1} (good bounds on the error term or, alternatively, the error probability in Theorem~\ref{thm:main_general}),
Proposition~\ref{lem:virtual} 
and Conjecture~\ref{conj:virtual} (virtual permutations),
Theorem~\ref{thm:conjugacy} (conjugacy characters),
and Corollary~\ref{cor:conjugacy} ($\S_n$-actions on cohomology).
%}

The rest of the paper is organized as follows. 
In Section~\ref{sec:preliminaries} we 
%introduce notation, and cite basic facts and known results used in the paper.
give some necessary background,
and introduce the concept of characters tending to be regular. 
In Section~\ref{sec:1row} we refine a result of Swanson regarding the classical (one-row) Lie character.
In Section~\ref{sec:rectangular} we prove our first main result, Theorem~\ref{thm:rectangular} (a strong version of Theorem~\ref{main:rectangular}), concerning rectangular diagrams.
Sections~\ref{sec:distinct} and~\ref{sec:gluing} contain two useful results, regarding diagrams with distinct parts and the Gluing Lemma. 
%which will also be used later,
These results are used in Section~\ref{sec:asymp} for the proof of Theorem~\ref{thm:main_general}, which implies Theorem~\ref{thm:main}.
Variants of the main results appear in Section~\ref{sec:variations}, and Section~\ref{sec:questions} concludes with some further questions.


\section{Preliminaries}\label{sec:preliminaries}


\subsection{Notation}

We identify partitions with the corresponding Young (or Ferrers) diagrams.
For a 
% Young diagram 
partition
$\la = (\la_1, \la_2, \dots) = (1^{m_1} 2^{m_2} \cdots) \vdash n$,
\begin{itemize}
    \item 
    $\la_1$ and $\la'_1$ are the lengths of the first row and the first column of $\la$, respectively;
    \item 
    $\chi^\la$ is the irreducible character of $\S_n$ corresponding to $\la$;
    \item 
    $f^\la=\dim\chi^\la$;
    \item 
    $C_\la$ is the conjugacy class in $\S_n$ consisting of all the permutations with cycle type $\la$;
    \item 
    $z_\la:=\prod_ii^{m_i}m_i!$, so that $|C_\la| = n! / z_\la$.
\end{itemize}
 
By $e_\la, h_\la, p_\la, s_\la$ we denote the elementary, complete homogeneous, power sum, and Schur symmetric functions, respectively.

%In what follows, we will deal with homogeneous symmetric functions, for which we make the following definition.

%\begin{definition}\label{def:schur_le}
%    Let $f$ and $g$ be symmetric functions, homogeneous of degree $n$. Write $f \le_s g$ if 
%   \[
%        \langle f,s_\la \rangle \le \langle g,s_\la \rangle 
%        \qquad (\forall \la \vdash n).
%    \]
%\end{definition}

\subsection{Higher Lie characters}\label{sec:hLc}

For $\la=(i^{m_i})\vdash n$, the higher Lie character $\psi^\la$ can be constructed as follows. 
Denote by $Z_\la$ the centralizer in $\S_n$ of a permutation of cycle type $\la$. Then
\[
    Z_\la
    \simeq \prod_{i=1}^n\S_{m_i}[\bbz_i],
\]
where $\bbz_i$ is the cyclic subgroup of $\S_n$ generated by a cycle of length $i$, and $\S_{m_i}[\bbz_i]$ is the wreath product of $\bbz_i$ by $\S_{m_i}$. For each $i$, let $\zeta_i$ be a primitive irreducible character of $\bbz_i$. Then
\begin{equation}\label{eq:wreathproduct}
     \psi^\la
     = \Ind_{Z_\la}^{\S_n} \left(\id_{m_1}[\zeta_1] \otimes \id_{m_2}[\zeta_2] \otimes \dots \otimes \id_{m_n}[\zeta_n] \right),
\end{equation}
where $\id_{m_i}[\zeta_i]$ is the wreath product of $\zeta_i$ by the trivial character $\id_{m_i}$ of $\S_{m_i}$.

In particular, the character $\psi^{(i)}$ corresponding to the one-row 
%partition
diagram 
$\la=(i)$ is the much-studied ordinary Lie character corresponding to the action of $\S_n$ on the multilinear component of the free Lie algebra with $n$ generators.

In what follows, we also use the following formula for higher Lie characters (\cite{Thrall}):%, see also \cite{CalHanSun}):
\begin{equation}\label{eq:hlcformula}
\psi^{\la}=\Ind_{\S_{m_1}[\S_1]\times\S_{m_2}[\S_2]\times\ldots}^{\S_n}
\left(\id_{m_1}[\psi^{(1)}]\otimes\id_{m_2}[\psi^{(2)}]\otimes\ldots\right),
\end{equation}
where $\S_{m_i}[\S_i]$ is the wreath product of $\S_i$ by $\S_{m_i}$ and $\id_{m_i}[\psi^{(i)}]$ is the wreath product of the character $\psi^{(i)}$ of $\S_i$ by the trivial character $\id_{m_i}$ of $\S_{m_i}$.



Let $L_\la$ be the Frobenius image of the higher Lie character $\psi^\la$ in the algebra of symmetric functions. 
These functions are sometimes called {\em Lyndon symmetric functions}.  
For convenience, denote
$\Lie_n := L_{(n)}$, the Frobenius image of the ordinary Lie character.

% \todo{YR: Quasisymmetric functions. 

% The quasisymmetric approach will not be used in this paper. 
% However, to have a better understanding we briefly sketch it here. 
% Exceptionally, Theorem~\ref{t:sum_of_hLc} below has a minor use in the paper; note that this theorem has more complicated proofs using other techniques . 

% {\bf Define: QSF, fundamental QSF} 

% Using the Gessel--Reutenauer necklace bijection~\cite{GR93}, Thrall's problem can be recast as follows.


% For a partition $\lambda \vdash n$, let $C_\lambda$ be the conjugacy class consisting of all permutations of cycle type $\lambda$. This set is Schur-positive, and thus there exist nonnegative integers $c_\mu^\lambda$ such that
% %\begin{equation}\label{eq_Alambda_1}
% \[
%     Q(C_\lambda)
%     := \sum_{\pi \in C_\lambda} F_{n,\Des(\pi)}
%     = \sum_\nu c_\nu^\lambda s_\nu.
% \]
% %\end{equation}
% By~\cite{GR93}, 
% the Frobenius preimage of $\Q(C_\lambda)$ is the higher Lie character
% $\psi^\lambda$. 
% %It follows that higher Lie characters can be used to prove the existence of cyclic descent extensions, as explained in~\cite{AHR_conj}.
% }
% \todo{
The following theorem is well known; see, e.g., \cite[Theorem VII]{Thrall}.

\begin{theorem}\label{t:sum_of_hLc}
    For every $n\ge 1$
    \[
        \sum_{\la \vdash n} L_\la 
        = h_1^n 
        = \sum_{\la\vdash n} f^\la s_\la .
    \]
\end{theorem}

% \begin{proof}
% \[
% \sum\limits_{\la \vdash n} L_\la =  \sum\limits_{\la \vdash n} \sum\limits_{\pi\in C_\la} F_{\Des(\pi)} = \sum\limits_{\pi\in S_n}
% F_{\Des(\pi)} = \sum\limits_{\la\vdash n} f^\la  \sum\limits_{T\in \SYT(\la)} F_{\Des(T)} =   \sum\limits_{\la\vdash n} f^\la s_\la. 
% \]
% \end{proof}    

% }


\subsection{The Plancherel measure}\label{sec:Plancherel}

% \todo{RA: 
% In the description below of $\nu=(\nu^{(1)},\nu^{(2)},\ldots)$, is it required that $\nu^{(n)} \subset \nu^{(n+1)}$ ?
% }
The (infinite-dimensional) {\em Plancherel measure} is a distinguished measure on the space of infinite Young tableaux, 
%that is, of sequences $\nu=(\nu^{(1)},\nu^{(2)},\ldots)$ of Young diagrams $\nu^{(n)}\vdash n$, where $\nu^{(n)} \subset \nu^{(n+1)}$.
% of Young diagrams 
which plays an important role in the Vershik--Kerov asymptotic representation theory \cite{VershikKerovPlanch, VershikKerovChar}. In particular, it corresponds to the character of the regular representation of the infinite symmetric group. 
%The finite-dimensional projections of this measure coincide with the finite Plancherel measures:
% on Young diagrams:
%\[
%    \Pl^n(\nu^{(n)}) 
%    = \frac{( f^{\nu^{(n)}})^2}{n!}
%    \qquad (\forall \nu \vdash n).
%\]

This measure is defined as follows (see~\cite{VershikKerovPlanch}). 
Denote by ${\mathcal T}_n$ the space of of all sequences 
$\nu=(\nu^{(1)},\ldots, \nu^{(n)})$ of Young diagrams $\nu^{(i)}\vdash i$, where $\nu^{(i)} \subset \nu^{(i+1)}$; this is the same as the space of all standard Young tableaux with $n$ cells.
Consider the probability measure on ${\mathcal T}_n$ defined by
\[
    M_n(\nu)
    = \frac{f^{\nu^{(n)}}}{n!} 
    \qquad (\forall\, \nu\in{\mathcal T}_n).
\]
We have the natural projections $p_n:{\mathcal T}_n\to{\mathcal T}_{n-1}$, forgetting the last element of a sequence, and ${\mathcal T}:=\varprojlim{\mathcal T}_n$ is the space of infinite sequences  $\nu =(\nu^{(1)},\nu^{(2)},\ldots)$ of Young diagrams $\nu^{(n)}\vdash n$, where $\nu^{(n)} \subset \nu^{(n+1)}$.
It is easy to check that the pushforward of $M_n$ by $p_n$ coincides with $M_{n-1}$.

\begin{definition}
The measure $M:=\varprojlim M_n$ on the
space ${\mathcal T}$ of infinite sequences of Young diagrams is called the (infinite-dimensional) {\em Plancherel measure}.  
\end{definition}

The finite-dimensional marginals of $M$ coincide with the finite Plancherel measures: for every $\al \vdash n$,
\[
M \left( \left\{\nu =(\nu^{(1)},\nu^{(2)},\ldots) \in {\mathcal T} \,:\, \nu^{(n)} = \al \right\} \right) = \frac{(f^\al)^2}{n!}.
\]

\begin{remark}\label{rem:Plancherel_is_balanced}
    The famous Vershik--Kerov limit shape theorem~\cite{VershikKerovPlanch} implies that, for 
    % almost every (with respect to the Plancherel measure)
    a Plancherel-typical
    sequence of Young diagrams $\nu$, the first row and the first column of $\nu^{(n)}$ grow like $2\sqrt n$  as $n\to\infty$.
    %, so $\nu$ is typical according to Definition~\ref{def:typical}.     
\end{remark}

%\begin{remark}\label{rem:Planch}
%    It follows that all theorems proved below remain valid if we replace ``for typical $\nu$'' by ``for almost every $\nu$ with respect to the Plancherel measure.''    
%\end{remark}


\subsection{Tending to be regular}

Recall that, by Definition~\ref{def:regularity},
a sequence of characters $\rho = (\rho_n)$ tends to be regular if 
    \begin{equation}\label{eq:tends_to_regular2}
        \frac{\langle \rho_n,\chi^{\nu^{(n)}} \rangle}{\dim\rho_n}
        \sim \frac{f^{\nu^{(n)}}}{n!}
        \quad \text{as} \quad n\to\infty
    \end{equation}
for Plancherel-typical $\nu = (\nu^{(n)})$.  
However, all of the results in this paper hold under a weaker assumption on the test sequence $\nu$, namely for balanced sequences in the sense of the following definition.

\begin{definition}\label{def:balanced}\cite{DousseFeray}
    A sequence of diagrams $\nu =(\nu^{(1)},\nu^{(2)},\ldots)$ with $\nu^{(n)} \vdash n$ is said to be {\em balanced} if the lengths of the first row and the first column of $\nu^{(n)}$ are both $O(\sqrt n)$.
\end{definition}

By 
%the Vershik--Kerov limit shape theorem~\cite{VershikKerovPlanch}, 
Remark~\ref{rem:Plancherel_is_balanced},
a Plancherel-typical sequence of diagrams is balanced, so if~\eqref{eq:tends_to_regular2} holds for balanced sequences $\nu$, then $\rho$ tends to be regular. 


% \todo{RA: 
% In the description below of $\nu=(\nu^{(1)},\nu^{(2)},\ldots)$, is it required that $\nu^{(n)} \subset \nu^{(n+1)}$ ?
% }

Further, if $\rho_n = \psi^{\la^{(n)}}$ is %and $\beta_n = \Reg_n$ 
the higher Lie character corresponding to a diagram $\la^{(n)} \vdash n$,
%and the regular character of $\S_n$, respectively.
then $\dim\psi^{\la^{(n)}} = |C_{\la^{(n)}}| = n!/z_{\la^{(n)}}$.
%and $\dim\Reg_n = n!$; 
%also, $\langle \Reg_n,\chi^{\nu^{(n)}}\rangle = f^{\nu^{(n)}}$.
%So, for a sequence of Young diagrams $\la^{(n)} \vdash n$,
In this case
\eqref{eq:tends_to_regular2}~takes the form
\[
    \langle \psi^{\la^{(n)}},\chi^{\nu^{(n)}} \rangle
    \sim
    \frac{f^{\nu^{(n)}}}{z_{\la^{(n)}}} 
    \quad \text{as} \quad n\to \infty.
\]
For convenience, we will abbreviate this to 
\begin{equation}\label{claim}
    \langle \psi^\la,\chi^\nu\rangle\sim\frac{f^\nu}{z_\la}
    \qquad\text{or, equivalently,}\qquad
    \langle L_\la,s_\nu\rangle\sim\frac{f^\nu}{z_\la},
\end{equation}
where the second asymptotic equality is a reformulation of the first one 
in terms of (Lyndon and Schur) symmetric functions.

\subsection%{Useful results}
{Standard Young tableaux: enumeration and bounds}
\label{sec:useful_results}

% \todo{This phrase is no longer relevant.

% For convenience, we cite two theorems which, combined, essentially provide the asymptotics of the ordinary Lie characters $\psi^{(n)}$ (see Theorem~\ref{prop:Lie} below).
% }

In this subsection we quote results regarding 
%the enumeration and asymptotics of 
standard Young tableaux, which will be used in the paper. 
We begin with two theorems which, combined, essentially provide the asymptotics of the ordinary Lie characters $\psi^{(n)}$; see Theorem~\ref{prop:Lie} below.

For a skew shape $\la/\mu$, 
denote by $\SYT(\lambda/\mu)$ the set of standard Young tableaux of shape $\la/\mu$ and let $f^{\la/\mu}:=|\SYT(\la/\mu)|$. 
For $T\in \SYT(\la)$, define the {\em descent set}
\[
    \Des(T)
    := \{i \,:\, i+1 \text{ is in a lower row of } T \text{ than } i\} 
%\quad \text{and} \quad \maj(T):=\sum\limits_{i\in \Des(T)} i.
\]
and the {\em major index} 
%of $T$ is  
\[
    \maj(T) := \sum_{i \in \Des(T)} i.
\]

%\todo{Notation and definitions of SYT and $f$ of SYT of (skew) shapes; major index. }

\begin{theorem}{\rm (Kra\'skiewicz--Weyman~\cite{KraskiewiczWeyman}, see also~\cite[Theorem 8.4]{Garsia})}%
\label{thm:KW}
    For every partition $\nu\vdash n$, the multiplicity $\langle \psi^{(n)},\chi^{\nu} \rangle$ is equal to the cardinality of the set
    \[
        %A_{k,n,1}=
        \{T \in \SYT(\nu) \,:\, \maj(T) \equiv 1 \pmod n\}.
    \]	
%where $\SYT(\la)$ is the set of standard Young tableaux of shape $\la$ and $\maj(T)$ is the major index of $T$.
\end{theorem}


\begin{theorem}{\rm (Swanson~\cite[Theorem 1.8]{Swanson})}%
\label{thm:Swanson}
    For every partition $\nu\vdash n\ge 1$ and all $r$,
    \[
        \left| \frac{|\{T\in \SYT(\nu) \,:\, \maj(T) \equiv r \pmod n\} |}{f^\nu}-\frac{1}{n}
        \right|
        \le \frac{2n^{3/2}}{\sqrt {f^\nu}}.
    \]	
\end{theorem}

For further results in this direction see~\cite{
%BKS_SLC, BKS_AdvApplMath, 
Billey} and references therein. 
% \todo{NT: Do we really need all these references?\\
%  YR: References here and all over should be updated.}


The following result is due to Dousse and F\'eray~\cite{DousseFeray}.
%It is the the special case $\al(n) = D \sqrt n$, $\beta(m) = m$ of Theorem 8 there.

\begin{theorem}\label{t:DousseFeray}\rm{\cite[Theorem 8]{DousseFeray}} 
    There exist constants $C'_1$ and $C'_2$ such that the following holds: if $\la \vdash n$ with $\max\{\la_1, \la'_1\} \le \al(n)$, and $\mu \vdash m$ with $\max\{\mu_1, \mu'_1\} \le \beta(m)$, such that $m < C'_1 n / \al(n)$, then
    \[
        m! \cdot \frac{f^{\la/\mu}}{f^\la f^\mu} 
        \le e^{C'_2 m \beta(m) n^{-1} \al(n)}. 
    \]
\end{theorem}

Letting  $\al(n) = D \sqrt n$ and $\beta(m) = m$, and denoting $d := C'_1/D$ and $b := C'_2 D$, we deduce the following.

\begin{corollary}\label{lem:sub_generic}
    For any $D > 0$ there exist constants $d > 0$ and $b > 0$ such that, 
    %for any typical 
    if $\la \vdash n$ with $\max\{\la_1, \la'_1\} \le D \sqrt n$ and $\mu \vdash m$ with $m < d \sqrt n$, then
    \[
        \frac{f^{\la/\mu}}{f^\la} 
        \le e^{b m^{2} n^{-1/2}} \cdot \frac{f^\mu}{m!}.
    \]
\end{corollary}

% \todo{RA:
% Decide about terminology: Young diagram/shape/partition ?
% }
In several places in this paper we need a strengthening of the following result of Vershik and Kerov~\cite{VershikKerovChar} about dimensions of characters: 
if $m$ is fixed then, for every $\beta\vdash m$,
\begin{equation}\label{eq:skewdim}
    \frac{f^{\nu/\beta}}{f^\nu} 
    = \frac{f^\beta}{m!}(1+o(1))
    \quad \text{as} \quad n \to \infty
\end{equation}
for Plancherel-typical $\nu \vdash n$.
However, we need~\eqref{eq:skewdim} in the situation where $m$ is not fixed but grows not very fast, and $\nu$ is balanced but not necessarily Plancherel-typical. 
% The case we need is covered by 
% %%Theorem~\ref{t:DousseFeray}
% Corollary~\ref{lem:sub_generic}.
%
% \begin{lemma}\label{lem:strong_Vershik_Kerov}
%     If $m = o(n^{1/4})$, $\nu \vdash n$ is balanced, and $\beta \vdash m$ is arbitrary, then~\eqref{eq:skewdim} holds
%     uniformly on $\beta$ as $n \to \infty$,
%     with error term $O(m^2n^{-1/2})$.
% \end{lemma}
The case we need is covered by~\cite[Theorem 7]{DousseFeray}, with $r = 0$, $\al(n) = D \sqrt n$ and $\beta(m) = m$; it follows easily from its proof that the bound is uniform on $\beta$. 

\begin{lemma}\label{lem:strong_Vershik_Kerov}
    If $m = o(n^{1/4})$, $\nu \vdash n$ is balanced, and $\beta \vdash m$ is arbitrary, then~\eqref{eq:skewdim} holds
    uniformly on $\beta$ as $n \to \infty$,
    with error term $O(m^{2}n^{-1/2})$.
\end{lemma}

% \todo{RA: 
% Corollary~\ref{lem:sub_generic} gives an upper bound.
% What about a lower bound?
% }

% \begin{proof}
% Since $\nu$ is typical, $\max\{\nu_1,\nu_1'\}\le c\sqrt n$ for some constant $c$.
% Take $\al(n)=c\sqrt n$ and $\beta(k)=k$. Then, for $n$ large enough, $k<Cn/\al(n)$ for any constant $C>0$. Therefore, by Theorem~\ref{t:DousseFeray},
% \[
%     \frac{k!f^{\nu/\beta}}{f^\nu f^\beta}
%     \le e^{C' k \beta(k) n^{-1} \al(n)} 
%     = e^{C'' k^2  n^{-1/2}},
% \]
% which implies the lemma.
% \end{proof}




%\begin{theorem}\label{lem:sub_generic}\rm{\cite[Theorem 8]{DousseFeray}} 
%Assume that $\nu\vdash n$ is a  typical partition.  Then
%%\begin{itemize}
%%On the other hand, 
%%\item[(1)]
%%\item[(2)] 
%(Theorem 8 of~\cite{DousseFeray} 
%with $\al(n) = D \sqrt n$, $\beta(m) = m$)
%There exist constants $C$ and $c_4$ such that if $\mu \vdash n$ is typical (namely $\max\{\mu_1, \mu'_1\} \le D \sqrt n$) and $\al \vdash m$ is arbitrary, and $m < \frac{C}{D} \sqrt n$ (old: For every $m\le c_0 \sqrt n$ and $\al\vdash m$ there exist a constant $c_4>0$ such that),
%\[
%    \frac{f^{\mu/ \al}}{f^\mu} 
%    \le e^{c_4 D m^{2} n^{-1/2}} \cdot \frac{f^\al}{m!}. 
%\]
%%\end{itemize}
%\end{theorem}

%\todo{
%In other words, there exist constants  $\epsilon>0$ and $c_4>0$, such that 
%for any typical $\mu\vdash n$ (namely,  $\max\{\mu_1, \mu'_1\} \le D \sqrt n$) and any partition $\al\vdash m< \epsilon \sqrt n$
%\[
%    \frac{f^{\mu/ \al}}{f^\mu} 
%    \le e^{c_4 D m^{2} n^{-1/2}} \cdot \frac{f^\al}{m!}. 
%\]
%}

% \todo{
% \begin{theorem}
% (Theorem 7 of~\cite{DousseFeray})
% For any $m= o(n^{1/4})$ and $\al\vdash m$, 
% the probability  that the first $m$ letters in a random SYT of the typical limit shape $\nu$ form a SYT of shape $\al\vdash m$ is %asymptotically equal to
% \[
%     \frac{f^{\mu/ \al}}{f^\mu} 
%     = \left( 1+ O(m^{3/2}n^{-1/2}) \right) \cdot \frac{f^{\al}}{m!}.
% \]
% \end{theorem}
% }

% \todo{\rm NT: Did you mean \cite[Corollary~2]{DousseFeray}? Then this is not correct because this corollary is valid only for balanced $\alpha$. 
% On the other hand, \cite[Theorem~8]{DousseFeray}, which is cited below, does indeed imply our Lemma~\ref{lem:strong_Vershik_Kerov}.
% }



\section{The classical one-row Lie character}\label{sec:1row}

%\todo{In this section...}

%\todo{Align use of Lyndon symmetric functions versus higher Lie characters}

We begin with the proof of~\eqref{claim} for 
%the simplest case of a higher Lie character -- 
the ordinary Lie character 
$\psi^{(n)}$,
corresponding to the Lyndon symmetric function $\Lie_n = L_{(n)}$. 
In this case \eqref{claim} immediately 
follows from Theorems~\ref{thm:KW} and~\ref{thm:Swanson},  
but %for applications 
%in rest of the paper 
%further purpose %we also 
we want a good estimate of the error term.
%Recall Definition~\ref{def:regularity}.

% \todo{NT: Maybe promote Proposition~\ref{prop:Lie} to Theorem?}

\begin{theorem}\label{prop:Lie}
    The sequence of 
    %one-row diagrams $(n)$ 
    ordinary Lie characters $\psi^{(n)}$ 
    tends to be regular. 
    Furthermore, as $n \to \infty$,
    \[
        \langle \Lie_n,s_\nu \rangle
        \sim \frac{f^\nu}{n} 
    \]
    for balanced (and hence for Plancherel-typical) $\nu \vdash n$. 
    %More precisely,
    In fact,
    \[
        \langle \Lie_n,s_\nu \rangle
        = \frac{f^\nu}n(1+r_\nu)
    \]
    with the following two bounds on the error term $r_\nu$.
    \begin{enumerate}
    \item[\rm(a)] 
        For every $s > 0$ there exist constants $N > 0$ and $c > 0$ (both depending only on $s$) such that, for every $n \ge N$, if $\nu \vdash n$ with $\nu_1, \nu'_1 \le n-c$ then 
        \begin{equation}\label{swanson}
            |r_\nu| \le n^{-s}.
        \end{equation}
        %for sufficiently large $n$;
    
    \item[\rm(b)] 
        For balanced $\nu\vdash n$, there exists a constant $C>0$ (depending only on the constant in $O(\sqrt n)$ in Definition~\ref{def:balanced}) such that
        \begin{equation}
            |r_\nu| 
            \le \left(\frac{C}{n}\right)^\frac{n}{4}.
        \end{equation}
    
    \end{enumerate}
\end{theorem}

\begin{remark}
    The number of partitions $\nu \vdash n$ not satisfying the condition in Theorem~\ref{prop:Lie}(a) is bounded by a constant depending on $s$ but not on $n$.    
\end{remark}

\begin{proof}
    (a) %Write $\Lie_n = \sum_{\nu \vdash n} a_\nu s_\nu$. 
%\todo{YR suggestion: 
%By [Rasala , ??], for every $s\in\bbn$ there exists a constant $c$ such that if $\nu_1,\nu'_1<n-c$, then $f^{\nu} \ge n^{2s+5}$ for sufficiently large~$n$. Thus, by    Theorems~\ref{thm:KW} and~\ref{thm:Swanson} 
%%and the proof of \cite[Theorem~1.9]{Swanson}, 
%since for typical $\nu$ we have $\nu_1,\nu'_1 = o(n)$, so $\nu_1,\nu'_1 < n-c$ for sufficiently large $n$, which completes the proof.
%}
    By Theorems~\ref{thm:KW} and~\ref{thm:Swanson},
    \[
        \left| \frac{\langle \psi^{(n)},\chi^{\nu} \rangle}{f^\nu}-\frac{1}{n} \right|
        \le \frac{2n^{3/2}}{\sqrt {f^\nu}}.
    \]	
    In other words,
    \[
        |r_\nu| \le \frac{2n^{5/2}}{\sqrt {f^\nu}}.
    \]
    Given a Young diagram $\nu$ and a cell $(i,j)\in\nu$, denote by $h_{i,j}^{\rm op}$ the opposite hook length at $(i,j)$, i.e., $h_{i,j}^{\rm op}:=i+j-1$. Fix $c>0$ to be chosen later, and let $n\ge 2c$.
    Observe that if $\nu_1,\nu'_1\le n-c$, then $i,j\le n-c$ while $ij\le n$, whence
    \[
        i+j-1\le n-c+\frac n{n-c}-1\le n-c+1,
    \] 
    and for the diagonal excess $N(\nu)$ of $\nu$ we obtain 
    \[  
        N(\nu):=|\nu|-\max\limits_{(i,j)\in\nu}h_{i,j}^{\rm op}\ge c-1.
    \]

    Now fix $s>0$ and take $c=2s+7$.
    By~\cite[Corollary 4.13]{Swanson}, 
    \[
        f^\nu
        \ge \frac{1}{c} \binom{n}{c-1}
        \ge 4 n^{2s+5}
    \]
    for sufficiently large~$n$ depending only on $s$. Thus, 
    \[
        |r_\nu| \le \frac{2n^{5/2}}{\sqrt {f^\nu}}
        \le \frac1{n^s}.
    \]
    
    %But for typical $\nu$ we have $\nu_1,\nu'_1 = o(n)$, so $\nu_1,\nu'_1 < n-c$ for sufficiently large $n$, which completes the proof.

    (b) %By Theorem~\ref{thm:Swanson}, we have 
    As mentioned above, we always have $|r_\nu|\le 2n^{5/2}(f^\nu)^{-1/2}$. Let us bound $f^\nu$ from below for balanced~$\nu$. By Definition~\ref{def:balanced}, $\nu_1,\nu'_1\le c\sqrt n$ for some constant $c$. Thus, by the hook length formula and Stirling's formula,
    \[
        f^\nu\ge\frac{n!}{(2c\sqrt n)^n}
        \sim \sqrt{2\pi n}\left(\frac{\sqrt n}{2ce}\right)^n
        \ge \left( \frac{\sqrt n}{2ce} \right)^n, 
    \]
    whence
    \[
        |r_\nu|
        \le 2n^{5/2} \left( \frac{2ce}{\sqrt n} \right)^{\frac n2}
        \le \left( \frac{C}{\sqrt n} \right)^\frac n2.
     \qedhere
    \]
\end{proof}


\section{Rectangular diagrams}\label{sec:rectangular}

%\todo{NT: Explain why this case is important.}

Our next goal is to prove~\eqref{claim}, with a good error estimate, for 
%a Young diagram $\la \vdash n = mk$ of 
rectangular diagrams $\la = (m^k)$. 
Here $z_{(m^k)} = m^k k!$.

\begin{theorem}\label{thm:rectangular}
    Any sequence of higher Lie characters indexed by rectangular diagrams $\la=(m^k)$ with $m\to\infty$ and $k = o(m)$ tends to be regular. Furthermore,
    \begin{equation*}%\label{rectangular}
        \langle L_{(m^k)},s_\nu \rangle
        \sim \frac{f^\nu}{m^kk!} 
        \quad \text{as} \quad m \to \infty
    \end{equation*}
    for balanced (and hence for Plancherel-typical) $\nu \vdash mk$. 
    
%    More precisely, fix $C, c, \delta >0$. For any $\al > 0$ there exists $M > 0$, depending only on $\alpha$, $C$, $c$ and $\delta$,
%    such that
%    \begin{equation}\label{rectangularerror}
%        \langle L_{(m^k)},s_\nu \rangle
%        = \frac{f^\nu}{m^kk!} (1+R_\nu)
%        \quad \text{ with } \quad |R_\nu|\le m^{-\al}
%    \end{equation}
%    for any $m \ge M$, $k \le c m^{1 - \delta}$ and $\nu \vdash km$ satisfying $\nu_1, \nu'_1 \le C \sqrt{km}$.
\end{theorem}

%\todo{YR: Improve the bound on the error term.

%NT: Currently impossible, because we must use Theorem~\ref{prop:Lie} for almost all $\nu$ in order to get Theorem~\ref{thm:rectangular} for typical $\nu$.}

% \todo{RA:
% Natalia will go over the proof and see what bound is actually obtained, without any additional assumptions.

% NT: No good estimate for the error term can be obtained without a stronger assumption on $k$ than $k=o(m)$, so I removed the second part of the statement.
% }

\begin{proof}
    The case $k=1$ is essentially Theorem~\ref{prop:Lie}(a).

    Assume that $k \ge 2$, and let $s > 0$.
    Fix $c$ as in the proof of Theorem~\ref{prop:Lie}(a). We will say that 
    $\la\vdash m$ is $c$-\textit{good} if $\max\{\la_1,\la_1'\}\le m-c$ and 
    $c$-\textit{bad} otherwise.
    Then, for sufficiently large $m$,
    \begin{equation}\label{aux}
        \Lie_m
        = \sum_{\la\vdash m}\frac{f^\la}m(1+r_\la) s_\la
    \end{equation}
    where for $c$-good $\la$, from~\eqref{swanson} we have $|r_\la|\le m^{-s}$.

By an easy generalization of~\cite[(8.8)]{Macdonald}, for any symmetric functions $f_i$ we have
\[
    h_k \left[ \sum f_i \right]
    = s_{(k)} \left[ \sum f_i \right]
    = \sum_{\substack{j_i \ge 0 \\ \sum j_i=k}}\prod_i h_{j_i}[f_i].
\]
For a collection of indices $\{j_\la\}_{\la\vdash m}$, denote 
\[
    r_{\{j_\la\}}
    := \prod_\la(1+r_\la)^{j_\la}-1.
\]
By \eqref{eq:hlcformula} and \eqref{aux},
\begin{eqnarray}
    L_{(m^k)}
    &=& h_k[\Lie_m] 
    = \sum_{\substack{j_\la \ge 0 \\ \sum j_\la=k}} \prod_{\la \vdash m}h_{j_\la} \left[ \frac{f^\la}m(1+r_\la) s_\la \right] 
    = \frac{1}{m^k} \sum_{\substack{j_\la \ge 0 \\ \sum j_\la=k}} \prod_\la(1+r_\la)^{j_\la} \prod_{\la\vdash m} h_{j_\la} \left[ f^\la s_\la \right] \nonumber \\
    &=& \frac{1}{m^k}\sum_{\substack{j_\la \ge 0 \\ \sum j_\la=k}} (1+r_{\{j_\la\}}) \prod_{\la \vdash m} h_{j_\la} \left[ f^\la s_\la \right]
    = \frac{1}{m^k}\sum_{\substack{j_\la \ge 0 \\ \sum j_\la=k}} \prod_{\la \vdash m} h_{j_\la} \left[f^\la s_\la \right] + R \nonumber \\
    &=& \frac{1}{m^k}h_k \left[ \sum_{\la \vdash m} f^\la s_\la \right] + R 
    = \frac{1}{m^k}h_k \left[h_1^m \right] + R, 
\label{eq:R}
\end{eqnarray}
where
\[
    R = \frac{1}{m^k}\sum_{\substack{j_\la \ge 0 \\ \sum j_\la=k}} r_{\{j_\la\}} \prod_{\la\vdash m} h_{j_\la}\left[f^\la s_\la\right].
\]
Now let $R=R_1+R_2$,
where $R_1$ includes all summands for which $j_\la \ne 0$ only for $c$-good $\la$, and $R_2$ includes all summands for which $j_\la\ne0$ for at least one $c$-bad $\la$.

In order to bound $R_1$, note that for $0 \le x \le 0.5$,
\[
    -2x \le \ln(1-x) \le 0 \le \ln(1+x) \le x ,
\]
and therefore, for $-0.5 \le x \le 0.5$,
\[
    |\ln(1+x)| \le 2|x|.
\]
Similarly, for $0 \le x \le \ln 2$,
\[
    -x \le e^{-x} - 1 \le 0 \le e^x - 1 \le 2x,
\]
and therefore, for $-\ln 2 \le x \le \ln 2$,
\[
    |e^x - 1| \le 2|x|.
\]
For each summand of $R_1$ we have $|r_\la| \le m^{-s} \le 0.5$ (for $m$ large enough), so that
\[
    \left| \ln \left( \prod_\la(1+r_\la)^{j_\la} \right) \right|
    \le \sum_\la j_\la |\ln(1+r_\la)|
    \le 2 \sum_\la j_\la |r_\la|
    \le 2k m^{-s}.
    %= o(m^{-s+1}).
\]
%since $|r_\la| \le m^{-s}$ and $\sum_\la j_\la = k = o(m)$. 
Then, as $m \to \infty$, $2k m^{-s} \le \ln 2$ and therefore
\begin{align*}
    |r_{\{j_\la\}}|
    &= \left| \prod_\la(1+r_\la)^{j_\la} - 1 \right|
    = \left| \exp \left( \ln \left( \prod_\la(1+r_\la)^{j_\la} \right) \right) -1 \right|
    \le 2 \left| \ln \left( \prod_\la(1+r_\la)^{j_\la} \right) \right| \\
    &\le 4k m^{-s}
    = o(m^{-s+1}).
\end{align*}

Further, it is well known that the plethysm of two Schur functions is Schur-positive.
Therefore,
\begin{eqnarray}
    |\langle R_1,s_\nu \rangle|
    &\le& o(m^{-s+1})\left\langle \frac{1}{m^k} \sum_{\substack{\sum j_\la=k \\ j_\la=0\text{ for $c$-bad }\la}} \prod_{\la\vdash m} h_{j_\la} \left[f^\la s_\la \right], s_\nu \right\rangle \\
    &\le& o(m^{-s+1}) \left\langle \frac{1}{m^k} \sum_{\sum j_\la=k} \prod_{\la\vdash m} h_{j_\la} \left[f^\la s_\la \right], s_\nu \right\rangle \nonumber \\
    &=& o(m^{-s+1}) \left\langle \frac{1}{m^k}h_k \left[ \sum_{\la\vdash m} f^\la s_\la \right], s_\nu \right\rangle
    = o(m^{-s+1}) \left\langle \frac{1}{m^k} h_k \left[h_1^m \right], s_\nu \right\rangle.
\label{sigma1}
\end{eqnarray}

Now, each summand from $R_2$ contains a factor of the form $h_{j}[s_\la]$ with $c$-bad $\la$, for which either $\la_1\ge m-c$ or $\la'_1\ge m-c$. Then, by \cite[Theorem~5.1]{GutierrezRosas}, the decomposition of $h_j[s_\la]$ in the Schur basis contains only $s_\tau$ with $\tau\supset\la$, hence with either $\tau_1\ge m-c$ or $\tau'_1\ge m-c$.
But then the decomposition of the whole product $\prod_\la h_{j_\la}[s_\la]$ contains only $s_\tau$ with either $\tau_1\ge m-c$ or $\tau'_1\ge m-c$.
However, for balanced $\nu\vdash km$ we have $\nu_1,\nu'_1=O(\sqrt{km})=o(m)$, which means that for sufficiently large $m$ the corresponding inner product vanishes. Thus, $\langle R_2,s_\nu\rangle=0$ for sufficienly large~$m$. 

Combining this conclusion with~\eqref{eq:R} and~\eqref{sigma1}, we obtain
\begin{equation}\label{firststep}
    \langle L_{(m^k)},s_\nu\rangle
    %=\left(1+o(m^{-s+1})\right)\left\langle h_k\left[\frac{h_1^m}m\right],s_\nu\right\rangle
    = \left(1+o(m^{-s+1})\right)\frac1{m^k}\langle h_k[h_1^m],s_\nu\rangle.
\end{equation}    
In order to deal with $\langle h_k[h_1^m],s_\nu\rangle$, we
expand $h_k$ in the basis of power sums:        
\[
    h_k[h_1^m]
    = \sum_{\mu\vdash k} \frac{p_\mu[h_1^m]}{z_\mu}
    = \frac{h_1^n}{k!} 
    + \sum_{\mu \ne (1^k)} \frac{p_{m\mu}}{z_\mu},
\]
where $m \mu \vdash mk$ is obtained from $\mu \vdash k$ by repeating each row $m$ times. 
Thus,  recalling that $\langle h_1^n,s_\nu \rangle = f^\nu$
and $\langle p_\tau,s_\nu\rangle = \chi^\nu(\tau)$ is the value of the irreducible character $\chi^\nu$ at permutations with cycle type $\tau$, we obtain
\begin{equation}\label{eq:A}
    \langle h_k[h_1^m],s_\nu \rangle
    = \frac{f^\nu}{k!} 
    + \sum_{\mu \ne (1^k)} \frac{\langle p_{m\mu},s_\nu \rangle}{z_\mu}
    = \frac{f^\nu}{k!}
    + \sum_{\mu \ne (1^k)}
    \frac{\chi^\nu(m\mu)}{z_\mu}
    = \frac{f^\nu}{k!}(1+A),
\end{equation}
with
\[
    A
    = \sum_{\mu \ne (1^k)}
    \frac{k!}{z_\mu}
    \frac{\chi^\nu(m\mu)}{f^\nu}=\sum_{e\ne g\in\S_k}\frac{\chi^\nu(mg)}{f^\nu},
\]
where, for $g \in \S_k$, the cycle type of the permutation $mg\in\S_{mk}$ is obtained from the cycle type of $g$ by repeating each part $m$ times. Now we use Roichman's~\cite{R96} estimate: there exist constants $0 < a < 1$ and $b > 0$ such that for every $\nu\vdash n$ and every $g\in\S_n$,
\[
    \frac{|\chi^\nu(g)|}{f^\nu}
    \le \left[ \max \left( \frac{\nu_1}{n}, \frac{\nu'_1}{n}, a \right) \right]^{b(n-\fix(g))},
\]
where $\fix(g)$ is the number of fixed points of $g$. In our case, since $\nu$ is balanced, $\frac{\nu_1}{mk},\frac{\nu'_1}{mk} = o(1)$; besides, $\fix(mg)=m\,\fix(g)$. Thus, denoting $q:=a^b$ (obviously, $0<q<1$), for sufficiently large $m$, the left-hand side does not exceed
$q^{m(k-\fix(g))}$, so 
\[
    |A|
    \le \sum_{e \ne g \in \S_k} q^{m(k-\fix(g))}
    = \sum_{i=0}^{k-1} D_{k,i} q^{m(k-i)},
\]
where the so-called ``rencontres number'' $D_{k,i}$ is the number of permutations $g\in\S_k$ with $\fix(g)=i$.
Clearly, $D_{k,i} = \binom{k}{i} D_{k-i,0}$, where $D_{k-i,0} = d_{k-i}$ are the usual derangement numbers. Thus
\[
    |A| 
    \le \sum_{i=0}^{k-1} \binom{k}{i} d_{k-i} q^{m(k-i)}
    = \sum_{j=1}^{k} \binom{k}{j} d_{j} q^{mj}
    \le \sum_{j=1}^{k} \binom{k}{j} j! q^{mj}
    \le \sum_{j=1}^{k} k^j q^{mj}
    =\sum_{j=1}^{k} (kq^m)^j.
\]
Observe that $kq^m=o(mq^m)=o(m^{-t})$ for any $t>0$, so $A = o(m^{-t})$. Plugging this into~\eqref{eq:A} and~\eqref{firststep}, we obtain
$$
\langle L_{(m^k)},s_\nu\rangle=\left(1+o(m^{-s+1})\right)\left(1+o(m^{-t})\right)
\frac{f^\nu}{m^kk!},
$$
which completes the proof.
\end{proof}




%\todo{YR: Are there counterexamples beyond $m=2$?} 

%\begin{conjecture}\label{conj1}
%The assumption $k=o(m)$ 
%%(or, which is the same, $k=o(\sqrt n)$) 
%in Theorem~\ref{thm:rectangular} 
%is necessary and sufficient.   
%\end{conjecture}



\section{%Shapes
Diagrams with distinct row lengths}\label{sec:distinct}


In this section we prove~\eqref{claim}, with a good error estimate, for 
%a Young diagram $\la \vdash n = mk$ of 
diagrams with distinct row lengths. The main result of this section, Theorem~\ref{thm:distinct}, plays a crucial role in the proof of Theorem~\ref{thm:main}. 
% will be used in Section~\ref{}
 %$\la = (m^k)$. Here $z_{(m^k)} = m^k k!$.

In the following theorem we actually consider a sequence $\mu^{(1)}, \mu^{(2)}, \dots$ of diagrams, but we write only a single $\mu$ for readability.
%\todo{
%$g(n) = C' (\ln n)^\beta$ and $f(n) \ge C (\ln n)^{\beta + 1}$.
%}

\begin{theorem}\label{thm:distinct}
    Let $\mu =(m_1, \dots, m_t) \vdash n$ with $m_1 > \dots > m_t$. 
    If $t\le g(n)$ and $m_t\ge f(n)$ for some functions $f,g$ such that $f(n) \to \infty$, $g(n) = O(f(n) / \ln n)$, and $g(n) = o(\sqrt n)$ (as $n \to \infty$), then $\psi^\mu$ tends to be regular. Furthermore,
    \[
        \langle L_{\mu},s_\nu \rangle 
        \sim \frac{f^\nu}{z_{\mu}} 
        \quad \text{as} \quad n\to \infty,
    \]
    %where $\nu$ is Plancherel-typical. 
    for balanced (and hence for Plancherel-typical) $\nu$.
    Moreover, the error term is uniform over $\nu$:
    \[
        \langle L_{\mu}, s_\nu \rangle 
        = \frac{f^\nu}{z_{\mu}} (1+R_\nu)
        \quad \text{where} \quad |R_\nu| = o(f(n)^{-q})
    \]
    for any fixed (arbitrarily large) $q$.
%    \[
%        \langle L_{\mu}, s_\nu \rangle 
%        = \frac{f^\nu}{z_{\mu}} (1+R_\nu)
%        \quad \text{where} \quad |R_\nu| < \omega(n)
%    \]
%    for some function $\omega(n) = o(1)$.
\end{theorem}

\begin{proof}
%    Let us first deal with $L_\mu$. 
%The Littlewood--Richardson coefficient $c^\nu_{\al\beta}$ is nonzero only for $\al\subset\nu$. Since $\nu$ is typical,
%the rows and columns of $\nu$ are $O(\sqrt{n})=O(\sqrt{n-k})$, hence the same is true for $\al\subset\nu$, i.e., $\al\vdash n-k$ is typical. 
%Denote by $t$ the number of different blocks $(m_i^{j_i})$ in $L_\mu$. 


Since $\mu$ has distinct parts, we have 
\[
    L_\mu
    = \prod_{i=1}^t L_{(m_i)}
    = \prod_{i=1}^t \Lie_{m_i}.
\]
%We will use the following well-known formula \red{(reference?)}: for arbitrary symmetric functions $f$ and $g$,
%\[
%    \langle fg,s_\nu \rangle
%    = \sum_{\alpha\subset\nu} \langle f,s_\alpha\rangle \langle g, s_{\nu/\alpha} \rangle 
%    =\sum_{\al,\beta\subset\nu} c^\nu_{\al\beta} \langle f,s_\al \rangle \langle g,s_\beta \rangle ,
%\]
%where $c^\nu_{\al\beta}$ are Littlewood-Richardson coefficients.
    Using the Hopf algebra structure on $\Lambda$, for any two symmetric functions $f$ and $g$, and Schur function $s_\nu$ we have
    \begin{eqnarray}\label{eq:Hopf}
        \langle fg, s_\nu \rangle
        &=& \langle f \otimes g, \Delta s_\nu \rangle
        = \langle f\otimes g, \sum_{\la \subseteq \nu} s_\la \otimes s_{\nu/\la} \rangle
        = \sum_{\la \subseteq \nu} \langle f, s_\la \rangle \langle g, s_{\nu/\la} \rangle \nonumber 
        \\
        &=& \sum_{\la \subseteq \nu} \langle f, s_\la \rangle \langle g,\sum_{\mu \subseteq \nu} c^\nu_{\la \mu} s_\mu \rangle 
        = \sum_{\la,\mu \subseteq \nu} c^\nu_{\la \mu} \langle f, s_\la \rangle \langle g, s_\mu\rangle.
    \end{eqnarray}
where $c^\nu_{\la \mu}$ are the Littlewood--Richardson coefficients.
By an obvious generalization of this identity
to an arbitrary number of factors,
%\begin{equation}\label{eq:decompmu2}
\[
    \langle L_\mu,s_\nu \rangle
    = \langle \prod_{i=1}^t L_{(m_i)}, s_\nu \rangle
    = \sum_{\substack{\nu_1,\ldots,\nu_t\\ \nu_i \vdash m_i}} c^\nu_{\nu_1,\ldots,\nu_t} \prod_{i=1}^t \langle L_{(m_i)},s_{\nu_i} \rangle,
\]
%\end{equation}
where  
$c^\nu_{\nu_1,\ldots,\nu_t}$ is a ``multi-Littlewood--Richardson coefficient'' equal to the coefficient of $s_\nu$ in the product $s_{\nu_1} \cdots s_{\nu_t}$.

For a partition $\nu \vdash m$, define an error term $r_\nu$ by
\begin{equation}\label{eq:def_error}
    \langle L_{(m)},s_\nu \rangle
    = \frac{f^\nu}{m} \cdot (1 + r_\nu) .
\end{equation}
%Equation~\eqref{eq:decompmu2} can then be rewritten as 
We can then write
\[
    \langle L_\mu,s_\nu \rangle
    = \sum_{\substack{\nu_1,\ldots,\nu_t \\ \nu_i \vdash m_i}} c^\nu_{\nu_1,\ldots,\nu_t} \prod_{i=1}^t \frac{f^{\nu_i}}{m_i} (1 + r_{\nu_i}) .
\]
Defining $r_{\nu_1, \dots, \nu_t}$ by
\[
    1 + r_{\nu_1, \dots, \nu_t} 
    := \prod_{i=1}^{t} (1 + r_{\nu_i}) 
\]
and recalling that $z_\mu  = m_1 \cdots m_t$,
this becomes
\[
    \langle L_\mu,s_\nu \rangle
    = \frac{1}{z_\mu} \sum_{\substack{\nu_1,\ldots,\nu_t \\ \nu_i \vdash m_i}} c^\nu_{\nu_1,\ldots,\nu_t} (1 + r_{\nu_1, \dots, \nu_t}) \prod_{i=1}^t f^{\nu_i} .
\]
Recalling (e.g., from~\cite[Corollary 7.12.5]{Stanley_ECII}) that
\[
    h_1^m = \sum_{\nu \vdash m} f^\nu s_\nu \,,
\]
we have
\begin{align*}
    \sum_{\substack{\nu_1,\ldots,\nu_t \\ \nu_i \vdash m_i}} c^\nu_{\nu_1,\ldots,\nu_t} \prod_{i=1}^t f^{\nu_i} 
    &= \sum_{\substack{\nu_1,\ldots,\nu_t \\ \nu_i \vdash m_i}} \langle s_{\nu_1} \cdots s_{\nu_t}, s_\nu \rangle \prod_{i=1}^t f^{\nu_i} 
    = \left\langle \prod_{i=1}^t \sum_{\nu_i \vdash m_i} f^{\nu_i} s_{\nu_i}, s_\nu \right\rangle 
    = \left\langle \prod_{i=1}^t h_1^{m_i}, s_\nu \right\rangle \\
    &= \langle h_1^n, s_\nu \rangle
    = f^\nu .
\end{align*}
Therefore
\[
    \langle L_\mu,s_\nu \rangle
    = \frac{f^\nu}{z_\mu} \cdot (1 + R) ,
\]
where
\begin{equation}\label{eq:total_error}
    R
    := \frac{1}{f^\nu} \sum_{\substack{\nu_1,\ldots,\nu_t \\ \nu_i \vdash m_i}} c^\nu_{\nu_1,\ldots,\nu_t} r_{\nu_1, \dots, \nu_t} \prod_{i=1}^t f^{\nu_i} .
\end{equation}
Our goal is to prove that, 
for $\mu \vdash n$ 
%(namely $m_1, \dots, m_t$) 
as in the statement of the theorem and $\nu \vdash n$ balanced, 
$R$ tends to zero (as $n$ tends to infinity)
uniformly over $\nu$.
By the definition of $\nu$ being balanced, there exist a constant $a>0$ such that
\begin{equation}\label{eq:balanced}
    \max\{\nu_1,\nu_1'\}\le a\sqrt n\,.   
\end{equation}

%\red{It follows from the Vershik--Kerov limit shape theorem~\cite{VershikKerovPlanch} %that for a.e.\ $\nu$ with respect to Plancherel measure we have, for $n$ large enough, 
%\begin{equation}\label{eq:planch}
%    a_1\sqrt n<\nu_1,\nu_1'<a_2\sqrt n \,,
%\end{equation}   
%where the constants $0<a_1<2<a_2$ can be chosen arbitrarily close to $2$.}

%In the rest of the proof we apply Theorem~\ref{prop:Lie}(a) with an arbitrary fixed $s>1$. 

Fix an arbitrary real number $s > 1$. %, to be chosen later.

By Theorem~\ref{prop:Lie}(a), there exist constants $M > 0$ and $c > 0$ (both depending only on~$s$) such that, for every $m \ge M$ and $\nu \vdash m$ with $\max\{\nu_1, \nu'_1\} \le m-c$, the error term~$r_\nu$ in Equation~\eqref{eq:def_error} satisfies
\[
    |r_\nu| \le m^{-s}.
\]
Call a partition $\nu \vdash m$ {\em 
$c$-good} if $\max\{\nu_1,\nu'_1\} \le m-c$, and {\em $c$-bad} otherwise.
Write 
\[
    R = R_1 + R_2 + R_3,
\]
where
\begin{enumerate}

    \item 
    $R_1$ includes the summands in Equation~\eqref{eq:total_error} for which all $\nu_i$ are $c$-good; 

    \item 
    $R_2$ includes the summands for which at least one $\nu_i$ is $c$-bad and has size $|\nu_i| = m_i < d \sqrt n$, with $d$ to be determined later; and

    \item 
    $R_3$ includes all other summands,
    namely: those for which there exists at least one $c$-bad $\nu_i$, and all the $c$-bad $\nu_i$ have size at least $d \sqrt n$.
        
\end{enumerate}

Let us first deal with $R_1$:
\[
    R_1
    = \frac{1}{f^\nu} \sum_{\substack{\nu_1,\ldots,\nu_t \\ \forall i\, \nu_i \vdash m_i \text{ is $c$-good}}} c^\nu_{\nu_1,\ldots,\nu_t} r_{\nu_1, \dots, \nu_t} \prod_{i=1}^t f^{\nu_i} .
\]
As observed in the proof of Theorem~\ref{thm:rectangular}, 
%For $0 \le x \le 0.5$,
%\[
%    -2x \le \ln(1-x) \le 0 \le \ln(1+x) \le x ,
%\]
%and therefore, 
%for $-0.5 \le x \le 0.5$,
\[
    |x| \le 0.5 
    \quad \Longrightarrow \quad 
    |\ln(1+x)| \le 2|x|
\]
%Similarly, for $0 \le x \le \ln 2$,
%\[
%    -x \le e^{-x} - 1 \le 0 \le e^x - 1 \le 2x,
%\]
and %therefore, 
%for $-\ln 2 \le x \le \ln 2$,
\[
    |x| \le \ln 2 
    \quad \Longrightarrow \quad
    |e^x - 1| \le 2|x|.
\]
As mentioned above, by Theorem~\ref{prop:Lie}(a), 
if $m \ge M$ and $\nu \vdash m$ is $c$-good, then $|r_\nu| \le m^{-s}$.
We can assume that $M \ge 2^{1/s}$. It follows that, if $m_t \ge M$ and $\nu_1, \dots, \nu_t$ are $c$-good, then (for each $1 \le i \le t$) $|r_{\nu_i}| \le m_i^{-s} \le M^{-s} \le 0.5$ and therefore
\[
    \left| \ln \left( \prod_{i = 1}^t (1 + r_{\nu_i}) \right) \right|
    \le \sum_{i = 1}^t |\ln(1+r_{\nu_i})|
    \le 2 \sum_{i = 1}^t |r_{\nu_i}|
    \le 2 \sum_{i = 1}^t m_i^{-s}
    \le 2t m_t^{-s} .
\]
Since $s>1$, $2t m_t^{-s} \le  2g(n)f^{-s}(n) \to 0$. Hence, for $n$ large enough, $2t m_t^{-s} \le \ln 2$ and 
\[
    |r_{\nu_1, \dots, \nu_t}|
    = \left| \exp \left( \ln \left( \prod_{i = 1}^t (1+r_{\nu_i}) \right) \right) - 1 \right|
    \le 2 \left| \ln \left( \prod_{i = 1}^t (1+r_{\nu_i}) \right) \right|
    \le 4t m_t^{-s} .
\]
Thus, as $n \rightarrow \infty$,
\begin{equation}\label{eq:R1}   
    |R_1| 
    = \left| \frac{1}{f^\nu} \sum_{\substack{\nu_1,\ldots,\nu_t \\ \forall i\, \nu_i \vdash m_i \text{ is $c$-good}}} c^\nu_{\nu_1,\ldots,\nu_t} r_{\nu_1, \dots, \nu_t} \prod_{i=1}^t f^{\nu_i} \right|
    \le 4t m_t^{-s} 
    %\le \red{4 c_1 c_2^{-s} (\ln n)^{1-2s}}
    \le 4g(n)f^{-s}(n)
    \rightarrow 0 \,.
\end{equation}
%Observe that, again, $c^\al_{\al_1,\ldots,\al_t}$ is nonzero only if $\al_i\subset\al$.
%We have seen that the rows and columns of $\al$ are $O(\sqrt{n})$, which is the same thing as $O(\sqrt{m_ij_i})$ in view of assumption~(a), and hence the same it true for $\al_i$. Thus, each $\al_i\vdash m_ij_i$ is typical.

%\todo{Theorem 4.1:

%    The sequence of one-row diagrams $(n)$ is regular:
%    as $n \to \infty$,
%    \[
%        \langle \Lie_n,s_\nu \rangle
%        \sim \frac{f^\nu}{n} 
%    \]
%    for  typical $\nu \vdash n$. 
%    More precisely, 
%    \[
%        \langle \Lie_n,s_\nu \rangle
%        = \frac{f^\nu}n(1+r_\nu)
%    \]
%    where
    %\begin{enumerate}
    %\item[\rm(a)] 
%        for every $s > 0$ there exist constants $c > 0$ and $N > 0$ (both depending only on $s$) such that, for every $n \ge N$, if $\nu \vdash n$ with $\nu_1, \nu'_1 \le n-c$ then 
        %\begin{equation}\label{swanson2}
%            $|r_\nu| \le n^{-s}$.
        %\end{equation}
        %for sufficiently large $n$;
    
    %\item[\rm(b)] 
    %    for typical $\nu\vdash n$, there exists a constant $C>0$ (depending only on the constant in $O(\sqrt n)$ in Definition~\ref{def:typical}) such that
    %    \begin{equation}
    %        |r_\nu|\le \left(\frac{C}{\sqrt n}\right)^\frac n2.
    %    \end{equation}
    
   %\end{enumerate}
%}
%Since $m_1,\ldots,m_t> c_2 n^b$, for sufficiently large $n$ we can apply Theorem~\ref{prop:Lie}(a) to obtain


Let us now deal with $R_2$:
\[
    R_2 
    = \frac{1}{f^\nu} \sum_{\substack{\nu_1,\ldots,\nu_t \\ \forall i\, \nu_i \vdash m_i \\
    \exists i\, \nu_i \text{ is $c$-bad and } m_i < d \sqrt n}} c^\nu_{\nu_1,\ldots,\nu_t} r_{\nu_1, \dots, \nu_t} \prod_{i=1}^t f^{\nu_i} .
\]
Recall (see Theorem~\ref{t:sum_of_hLc})
%and~\cite{GR93})
that the sum of all higher Lie characters is the regular character.
%, in which the multiplicity of each $s_\nu$ is $f^\nu$.
Equation~\eqref{eq:def_error} thus implies that, for any $\nu_i \vdash m_i$, 
 \[
    \frac{f^{\nu_i}}{m_i} (1+r_{\nu_i}) 
    = \langle L_{(m_i)}, s_{\nu_i} \rangle 
    \le \sum_{\mu \vdash m_i} \langle L_\mu, s_{\nu_i} \rangle
    = \langle h_1^{m_i}, s_{\nu_i} \rangle
    = f^{\nu_i},
\]
so that
\[
    0 \le 1 + r_{\nu_i} \le m_i .
\]
Then
\[
    0
    \le 1+ r_{\nu_1, \ldots, \nu_t} 
    = \prod_{i=1}^{t} \left( 1+r_{\nu_i} \right)
    \le \prod_{i=1}^{t} m_i
    \le n^{t}.
\]
Therefore
\begin{align*}
    |R_2| 
    &\le %e^{c_1 (\ln n)^2} 
    n^{t} \cdot
    \frac{1}{f^\nu} \sum_{\substack{\nu_1,\ldots,\nu_t \\ \forall i\, \nu_i \vdash m_i \\
    \exists i\, \nu_i \text{ is $c$-bad and } m_i < d \sqrt n}} c^\nu_{\nu_1,\ldots,\nu_t} \prod_{i=1}^t f^{\nu_i} \\
    &\le n^{t} \cdot
    \sum_{\substack{i_0 = 1 \\ m_{i_0} < d \sqrt n}}^{t} \,\,\sum_{\substack{\al \vdash m_{i_0} \\ \al \text{ is $c$-bad}}}
    \frac{1}{f^\nu} \sum_{\substack{\nu_1,\ldots,\nu_t \\ \forall i\, \nu_i \vdash m_i \\
    \nu_{i_0} = \al}} c^\nu_{\nu_1,\ldots,\nu_t} \prod_{i=1}^t f^{\nu_i} \\
    &= n^{t} \cdot
    \sum_{\substack{i_0 = 1 \\ m_{i_0} < d \sqrt n}}^{t} \,\,\sum_{\substack{\al \vdash m_{i_0} \\ \al \text{ is $c$-bad}}}
    \frac{f^{\nu/\al} f^\al}{f^\nu} .
\end{align*}
%Fix $D > 0$.
Recall that $\max\{\nu_1,\nu_1'\}\le a\sqrt n$ by~\eqref{eq:balanced}. Then,
by Corollary~\ref{lem:sub_generic},      
%for typical $\nu\vdash n$ and $c$-bad $\al=\al_{i_o} \vdash m_{i_0}< \epsilon \sqrt n$ %with $m_i = o(n^{1/4})$ 
there exist constants $d > 0$ and $b > 0$ such that
%, for any typical 
%if $\nu \vdash n$ with $\max\{\nu_1, \nu'_1\} \le D \sqrt n$, 
if $\al \vdash m$ with $m < d \sqrt n$, then
\[
    \frac{f^{\nu/\al}}{f^\nu} 
    \le e^{b m^{2} n^{-1/2}} \cdot \frac{f^\al}{m!}
    < e^{b d m} \cdot \frac{f^\al}{m!} \, .
\]  
It is easy to see that, for any $c$-bad $\al \vdash m$, $f^\al = O(m^c)$. 
It follows that, for a suitable constant $b' > 0$,
\[
    \frac{f^{\nu/\al} f^\al}{f^\nu} 
    \le \frac{e^{b d m}}{m!} \cdot (f^\al)^2
    \le \frac{e^{b d m}\cdot b' m^{2c}}{m!}
    \le\frac{b'e^{(b d+1)m}m^{2c}} {m^m}.
\]
The function
\[
    \ln\left(\frac{e^{(b d+1)m}m^{2c}} {m^m}\right)
    =(bd+1)m+2c\ln m-m\ln m
\]
is monotone decreasing for $m$ large enough.
A $c$-bad $\al$ has at most $c$ cells outside its first row (or column), which means that the total number of $c$-bad $\al \vdash m$ is at most a constant $c'$  independent of $m$. Denote $C=c'b'$.
By our assumptions $t \le g(n)$, $m_t \ge f(n) \to \infty$, and $g(n) = O(f(n) / \ln n)$, we thus obtain
\begin{align}\label{eq:fromR2}
    |R_2| 
    &\le n^{t} \cdot
    \sum_{\substack{i_0 = 1 \\ m_{i_0} < d \sqrt n}}^{t} \,\,\sum_{\substack{\al \vdash m_{i_0} \\ \al \text{ is $c$-bad}}}
    \frac{f^{\nu/\al} f^\al}{f^\nu} \nonumber\\
    &\le n^t \cdot t \cdot C \cdot 
   e^{(b d + 1) f(n) + 2c\ln f(n) - f(n) \ln f(n) } \\
%\begin{align*}
%    |R_2| 
%    &\le e^{c_1 (\ln n)^2} \cdot
%    \sum_{\substack{i_0 = 1 \\ m_{i_0} < \ve \sqrt n}}^{t} \,\,\sum_{\substack{\al \vdash m_{i_0} \\ \al \text{ is $c$-bad}}}
%    \frac{f^{\nu/\al} f^\al}{f^\nu} 
%    \le e^{c_1 (\ln n)^2} \cdot t \cdot C \cdot e^{(c' \ve + 1) m_t - m_t \ln m_t + O(\ln n)} \\
%    &\le e^{\red{- 2 c_2 (\ln n)^2 \ln \ln n + O((\ln n)^2)}} \to 0
%\end{align*}
    &\le C \cdot e^{g(n)\ln n+\ln g(n)
    +(b d + 1) f(n) + 2c\ln f(n) - f(n) \ln f(n) } ,\nonumber
\end{align}
so that, as $n \to \infty$,
\begin{equation}\label{eq:R2}
    |R_2| \le e^{-f(n)\ln f(n)(1+o(1))} \to 0 \,.  
\end{equation}


%if, this ratio is asymptotically equal to 

%  Now apply Corollary~\ref{lem:sub_generic} to show that the second sum is small.
%complete the proof of the upper bound on the error term. 

\medskip

%Now sum over $c$-bad $\al_i\vdash m_i$ with $O(n^{1/3}) \le m_i= O(n^{1/2})$. 

%\medskip

%It remains to sum over $c$-bad $\al_i\vdash m_i$ with $ m_i= O(n^{1/2})$.

%The key observation here is that their number is "small".
%}

%If $\nu$ is typical and $\al_i$ is $c$-bad then $m_i=o(n)$.
%This follows from the assumption that $\nu$ is typical and 
%$c^\nu_{\al_1,\dots,\al_t}\ne 0$ that 
%for every $i$, $m_i=o(n)$. Otherwise, 
%%there exists $C>0$ such that $\al_1\vdash m_1\ge Cn$ then for large enough $n$, 
%for some $i$, the first row of $\al_i$ is of order $n$, since the LR coefficient is nonzero $\al_i\subset \nu$, thus  
%$\nu_1>> \sqrt n$ and $\nu$ is not typical. 

%By Observation~\ref{obs:inclusion} and the fact that $\al_i$ is $c$-bad,  $m_i\le c_o \sqrt n+c$. We consider separately $c$-bad $\al_i\vdash m_i$, by the range of $m_i$.

%\medskip



%    Hence, for typical $\nu\vdash n$, and   with all $c$-bad $\al \vdash m_{i_0}$ with      
%    $m_i = o(n^{1/4})$ 

%    \smallskip

%We will consider separately $c$-bad $\al_i$-s with: 
%$m_i = o(n^{1/4})$; 
%$O(n^{1/4})\le m_i < O(\sqrt n)$; 
%and $m_{i_0} = O(\sqrt n)$. 

%\medskip
%By Corollary~\ref{lem:sub_generic}, summing the ratio in Equality~\eqref{eq:ratio} 

%    \[
%    \le  c_1 \log n\cdot C(c)\cdot n^{2c}  \cdot  n^{c_1 \log n} \cdot \left(\frac{e}{n}\right)^{bn^b}   \rightarrow 0.
%\]
%To explain the factor $n^{c_1 \log n}$ in the upper bound note that 
% for every $c$-bad $\al_i\le m_i$, 


%\todo{Thus, it suffices to bound the error on the first sum. Since for all $i$, $\al_i$ is $c$-good, by Theorem 4.1.a the error satisfies
%\[
%\prod_i (1+r_{\al_i}) -1 \le \prod_i (1+m_i^{-s}) -1 \sim \sum_i m_i^{-s}\le c_1 \log n \cdot c_2^{-s} n^{-bs} \rightarrow 0.
%\]
%
%\qed
%}

%\todo{

%First, recall that
%\[
%\sum_{\al_i\vdash m_i}c^\nu_{\al_1,\ldots\al_t}\prod f^{\al_i}=f^\nu,
%\]
%Hence, the first sum in \eqref{eq:1111} is  $\le (1+\sum_i  m_i^{-s})f^\nu \le (1+ \log n\cdot n^{bs}) f^\nu$. 


%$c$-bad is exponentillay decreasing with $m_1$.   
%By symmetry of LR coefficients, this holds for every $1\le i\le \ell$. 

%{\bf Proof of Observation.} First, observe that the shape $\nu_1$ formed by the prefix of 
%a SYT of a typical shape $n$ is also typical as long as $m_1=|\nu_1|$ tends to infinity. Now, 
%the number of SYT of a typical shape $\nu_1$ is asymptotically  
%equal to 
%\[
%\frac{\sqrt {m_1!}}{e^{\sqrt m_1}}.
%\]
%The number of SYT of shape $\nu_1$ that are $\epsilon$-bad is 
%\[
%\le \binom{m_1}{\epsilon m_1} \binom{\epsilon m_1}{\epsilon m_1/2} {\sqrt {(1-\epsilon)m_1!}\over e^{\sqrt{(1-\epsilon m_1}} }. 
%\]
%The quotient is exponentially decreasing with $m_1$. One concludes that the second sum\\ $\le  \cdot f^\nu \prod\limits_i e^{\sqrt m_i}\gamma^{-m_i}(1+m_i^{-1})$. The proof of lower bound is similar. 
%}




%\todo{

%    For $c$ and $N$ as in the proof of Theorem~\ref{prop:Lie}(a), let 
%    \[
%        B_m 
%        := \{\la\vdash m\colon f^\la< m^{2s+5}\}
%        \subset\{\la\vdash m\colon \la_1\ge m-c\text{ or } \la'_1\ge m-c\}
%    \] 
%    be the set of ``bad'' $\la$ for which \eqref{swanson} may fail. 

%}
%\todo{Now consider the sum}

%\todo{Old:\\
%For $s>1$ we have 
%\[
%    \left|\ln\prod_{i=1}^t(1+R_{m_i})\right|
%    = \left|\sum_{i=1}^t\ln(1+R_{m_i})\right|\le\sum_{i=1}^t|R_{m_i}|
%    \le \sum_{i=1}^t\frac1{{m_i}^s}
%    \le \sum_{j=1}^\infty\frac1{j^s}
%    \le\int_1^\infty\frac1{x^s}dx
%    =\frac{1}{s-1},
%\]
%hence
%\[
%    \left|\prod_{m_i>a}(1+R_{m_i})-1\right|
%    \le e^{1/(s-1)}-1
%    \le \varepsilon
%\]
%for sufficiently large $s$, which we now choose to satisfy this inequality. Then, in view of~\eqref{eq:decompmu},
%\[
%    \langle L_\nu,s_\al \rangle
%    = (1+R)\frac1{z_\mu}\sum_{\substack{\al_1,\ldots,\al_t\\ \al_i \vdash j_im_i}} c^\al_{\al_1,\ldots,\al_t} \prod_{i=1}^t f^{\al_i}
%    =(1+R)\frac{\langle h_1^{n-k},s_\al\rangle}{z_\mu}
%    =(1+R)\frac{f^\al}{z_\mu} 
%    \qquad\text{where}\qquad |R|<\varepsilon,
%\]
%since
%\[
%    \sum_{\substack{\al_1,\ldots,\al_t\\ \al_i \vdash j_im_i}} c^\al_{\al_1,\ldots,\al_t} \prod_{i=1}^t f^{\al_i}
%    =\sum_{\substack{\al_1,\ldots,\al_t\\ \al_i \vdash j_im_i}} c^\al_{\al_1,\ldots,\al_t} \prod_{i=1}^t \langle h_1^{j_im_i}, s_{\al_i}\rangle
%    =\langle h_1^{n-k},s_\al\rangle,
%\]
%again by the generalization of~\eqref{eq:product} to an arbitrary number of factors.
%}

\medskip

It remains to deal with $R_3$:
%which consists of the summands in~\eqref{eq:total_error} for which there exists at least one $c$-bad $\al_i$, and all the $c$-bad $\al_i$ have size $|\al_i| \ge \ve \sqrt n$.

\[
    R_3 
    = \frac{1}{f^\nu} \sum_{\substack{\nu_1,\ldots,\nu_t \\ \forall i\, \nu_i \vdash m_i \\
    \exists i\, \nu_i \text{ is $c$-bad} \\ \forall i\,\, \nu_i \text{ is $c$-bad } \Rightarrow\, m_i \ge d \sqrt n}} c^\nu_{\nu_1,\ldots,\nu_t} r_{\nu_1, \dots, \nu_t} \prod_{i=1}^t f^{\nu_i} .
\]
As in the discussion of $R_2$, we have
% \[
%     0 \le 1 + r_{\nu_1, \ldots, \nu_t} 
%     %< \prod_i \left| 1+r_{\nu_i} \right|
%     %\le \prod_i m_i\le n^{t}
%     %\le n^{c_1 \ln n} 
%     %= e^{c_1 (\ln n)^2} 
%     \le n^t
%     %=e^{g(n)\ln n}
% \]
% and therefore
\begin{align}\label{eq:expansionR3}
    |R_3| 
    %&\le e^{c_1 (\ln n)^2} \cdot
%    &\le n^t\cdot
%    \frac{1}{f^\nu} \sum_{\substack{\nu_1,\ldots,\nu_t \\ \forall i\, \nu_i \vdash m_i \\
%    \exists i\, \nu_i \text{ is $c$-bad} \\ \forall i\,\, \nu_i \text{ is $c$-bad } \Rightarrow\, m_i \ge d \sqrt n}} c^\nu_{\nu_1,\ldots,\nu_t} \prod_{i=1}^t f^{\nu_i} \nonumber\\
%    &\le n^t\cdot
%    \sum_{\substack{i_0 = 1 \\ m_{i_0} \ge d \sqrt n}}^{t} \,\,\sum_{\substack{\al \vdash m_{i_0} \\ \al \text{ is $c$-bad}}}
%    \frac{1}{f^\nu} \sum_{\substack{\nu_1,\ldots,\nu_t \\ \forall i\, \nu_i \vdash m_i \\
%    \nu_{i_0} = \al}} c^\nu_{\nu_1,\ldots,\nu_t} \prod_{i=1}^t f^{\nu_i} 
%    \nonumber\\
%    &= 
    &\le n^t \cdot
    \sum_{\substack{i_0 = 1 \\ m_{i_0} \ge d \sqrt n}}^{t} \,\,\sum_{\substack{\al \vdash m_{i_0} \\ \al \text{ is $c$-bad}}}
    \frac{f^{\nu/\al} f^\al}{f^\nu} .
\end{align}
Here we cannot use Corollary~\ref{lem:sub_generic}, and need a more delicate argument to bound $f^{\nu/\al}/f^\nu$.

Consider a $c$-bad $\al\vdash m_{i_0}$. For simplicity, write $m := m_{i_0}$.
By the $R_3$ assumption, $m \ge d \sqrt n$. 
Assume that $\al_1 = \max \{\al_1, \al'_1\}$; otherwise, exchange rows and columns in the following discussion.
Thus $\al_1 > m - c$.
For sufficiently large $n$, we have $d \sqrt n - c\ge \delta\sqrt n$ for a suitable constant $\delta>0$.
%Assume first that $\al = (\al_1) \vdash m_{i_0}$ has a single row. 
Since $\al$ is contained in $\nu$, we have by~\eqref{eq:balanced}
\begin{equation}\label{eq:alpah_nu_bounds}
    \delta\sqrt n
    \le d \sqrt n - c
    \le m - c
    < \al_1
    \le \nu_1 
    < a \sqrt n. 
\end{equation}

Given a diagram $\nu$, let $\nu^-:=\nu/(\nu_1)$ denote the (ordinary) diagram obtained from $\nu$ by removing the first row.
We can build a standard Young tableau of shape $\nu/\al$ by first choosing the entries in its first row (of length $\nu_1 - \al_1$), and then placing the other entries into  
$\nu^-/\al^-$ (if possible). Hence
\begin{eqnarray}\label{eq:skew}
   \frac{f^{\nu/\al} f^\al}{f^\nu}
   \le \binom{n-m}{\nu_1-\al_1} \frac{f^{\nu^-/\al^-} f^\al}{f^\nu}
   = \binom{n-m}{\nu_1-\al_1} \cdot \frac{f^{\nu^-/\al^-}}{f^{\nu^-}} \cdot \frac{f^{\nu^-}}{f^\nu} \cdot f^\al. 
\end{eqnarray}

We will bound, separately, each of the four factors in the RHS of~\eqref{eq:skew}.
To that end, we will use the following asymptotic formula (see, e.g., 
\cite[(5.12)]{Spencer}): for $k=o(n)$,
\begin{equation}\label{eq:binomial}
    \binom{n}{k}
    \sim \left( \frac{ne}{k} \right)^k (2\pi k)^{-1/2} e^{-\frac{k^2}{2n}(1+o(1))}
    \sim \frac{n^k}{k!} e^{-\frac{k^2}{2n}(1+o(1))}.
\end{equation}

For sufficiently large $n$, we have $\nu_1-\al_1 \le (a-\delta) \sqrt n$. 
%\todo{Check the case $\nu_1 - \al_1 = 0$.}
%Denoting $a := a-2\delta$ and 
%Using the unimodality of binomial coefficients and
Using the first asymptotic equality in~\eqref{eq:binomial} for $k = \nu_1 - \al_1$,
we obtain, for a suitable constant $d_1 > 0$,
%\begin{align*}
\begin{equation}\label{eq:factor1}
   \binom{n-m}{\nu_1-\al_1}
   \le \binom{n}{\nu_1-\al_1}
   \le d_1^{\sqrt n}n^{\nu_1-\al_1}(\nu_1-\al_1)^{-(\nu_1-\al_1)}\, 
   %\le {n\choose a\sqrt n}
   %\le \left(\frac{e\sqrt n}{a}\right)^{a\sqrt n}
   %= d_1^{\sqrt n}n^{(a_2/2-\delta)\sqrt{n}}.
   %\le C\frac{n^ne^{a\sqrt{n}}e^{n-a\sqrt{n}}}{e^n(a\sqrt{n})^{a\sqrt{n}}(n-a\sqrt{n})^{n-a\sqrt{n}}} \\
   %&= C \left( \frac{n}{n-a\sqrt{n}} \right)^{n-a\sqrt{n}}\left( \frac{n}{a\sqrt n} \right)^{a\sqrt n} 
   %= C \left(1+\frac{a\sqrt n}{n-a\sqrt{n}}\right)^{n-a\sqrt{n}} (a\sqrt n)^{a\sqrt n}.
%\end{align*}
\end{equation}
(with the convention that if $\nu_1-\al_1=0$, then $(\nu_1-\al_1)^{-(\nu_1-\al_1)}=1$).
%As in the proof of Lemma~\ref{lem:removingfirstrow},
%\[
%    \ln\left(1+\frac{a\sqrt n}{n-a\sqrt{n}}\right)^{n-a\sqrt{n}}
%    \le a\sqrt n + O(1).
%\]
%So, using Lemma~\ref{lem:removingfirstrow} and recalling that $a=a_2-2\delta$, we obtain
%\[
%    {n-m\choose \nu_1-\al_1}\le C'e^{a\sqrt n}(a\sqrt n)^{a\sqrt n}
%    = d_3\cdot d_4^{\sqrt n}n^{(a_2/2-\delta)\sqrt{n}}.
%\]

Observe that $\al^-\vdash c'$ with $c' < c$, and $\nu^-$ satisfies
$\nu^-_1,(\nu^-)'_1\le a\sqrt n$.
Hence, by Lemma~\ref{lem:strong_Vershik_Kerov},
\begin{equation}\label{eq:factor2}
    \frac{f^{\nu^-/\al^-}}{f^{\nu^-}}
    \sim \frac{f^{\al^-}}{c'!}
    = O(1).
\end{equation}

By the hook length formula, denoting by $h_{ij}$ the hook length at a cell $(i,j)$ of $\nu$ and observing that $h_{ij} \le \nu_1 + \nu_1' \le 2a \sqrt{n}$, we have
\[
    \frac{f^{\nu^-}}{f^\nu}
    = \frac{(n-\nu_1)!}{\prod_{(i,j)\in\nu^-}h_{ij}}
    \frac{\prod_{(i,j)\in\nu}h_{ij}}{n!}
    =\frac{(n-\nu_1)!\prod_{j=1}^{\nu_1}h_{1j}}{n!}
    \le\frac{(n-\nu_1)!}{n!}(2a\sqrt{n})^{\nu_1}.
\]
%Denote $k=a_1\sqrt n$. 
By the assumptions on~$\nu$ and the second asymptotic equality in~\eqref{eq:binomial} for $k = \nu_1$, 
\[
    \frac{(n-\nu_1)!}{n!} 
%    \le \frac{(n-a_1\sqrt n)!}{n!}
    = \left( \nu_1! \binom{n}{\nu_1} \right)^{-1}
    \sim  n^{-\nu_1} e^{\frac{\nu_1^2}{2n}(1 + o(1))} = n^{-\nu_1} \cdot O(1).
\]
Therefore,
\begin{equation}\label{eq:factor3}
   \frac{f^{\nu^-}}{f^\nu}
   \le O(1) \cdot n^{-\nu_1} \cdot
   (2a\sqrt{n})^{\nu_1}
   \le  d_2^{\sqrt n} n^{-\nu_1/2}.
   %\le d_2^{\sqrt n} n^{-a_1\sqrt n/2}.
\end{equation}

Finally, for any $c$-bad $\al \vdash m$, 
\begin{equation}\label{eq:factor4}
    f^\al = O(m^c) = O(n^c).
\end{equation} 

%Fix $0<\ve<\delta$ and
%choose $a_1,a_2$ in~\eqref{eq:planch} such that $a_2-a_1<\ve$. 
Plugging inequalities~\eqref{eq:factor1}, \eqref{eq:factor2}, \eqref{eq:factor3}, and~\eqref{eq:factor4} into~\eqref{eq:skew}, we obtain
\[
    \frac{f^{\nu/\al} f^\al}{f^\nu}
    \le d_1^{\sqrt n}n^{\nu_1-\al_1}(\nu_1-\al_1)^{-(\nu_1-\al_1)}
    %\le d_1^{\sqrt n}n^{(a_2/2-\delta)\sqrt{n}} 
    \cdot O(1) 
    \cdot d_2^{\sqrt n} n^{-\nu_1/2}
    %d_2^{\sqrt n} n^{-a_1\sqrt n/2} 
    \cdot O(n^c)
    \le d_3^{\sqrt n} n^{(\nu_1/2-\al_1)} (\nu_1-\al_1)^{-(\nu_1-\al_1)}\, .
%    \le d_3^{\sqrt n}n^{(\ve-\delta)\sqrt n}\,.
\]
Consider two cases. 
% \red{If $\al_1=\nu_1$, then $(\nu_1-\al_1)^{-(\nu_1-\al_1)}=1$, so that
% \[
%     \frac{f^{\nu/\al} f^\al}{f^\nu}
%     \le d_3^{\sqrt n} n^{(\nu_1/2-\al_1)}
%     = d_3^{\sqrt n} n^{-\nu_1/2}
%     \le d_3^{\sqrt n} n^{-\delta\sqrt n/2}.
% \]}
If $2\nu_1/3\le\al_1 \le \nu_1$, then 
$(\nu_1-\al_1)^{-(\nu_1-\al_1)} \le 1$ (by convention, also if $\nu_1 - \al_1 = 0$).
Also $\nu_1/2-\al_1 \le -\nu_1/6 \le - \delta \sqrt n/6$, so that
\[
    \frac{f^{\nu/\al} f^\al}{f^\nu}
    \le d_3^{\sqrt n} n^{-\delta \sqrt n/6}.
\]
On the other hand, if $\al_1<2\nu_1/3$, then $\nu_1-\al_1 > \nu_1/3 \ge \delta \sqrt n/3$, so that
\[
    (\nu_1-\al_1)^{-(\nu_1-\al_1)}
    \le (\delta \sqrt n/3)^{-(\nu_1-\al_1)}
    % \le d_4^{\sqrt n} n^{- \delta \sqrt n / 6}
\]
and
\[
    \frac{f^{\nu/\al} f^\al}{f^\nu}
    \le d_3^{\sqrt n}n^{(\nu_1/2-\al_1)} (\delta \sqrt n/3)^{-(\nu_1-\al_1)}
    \le d_4^{\sqrt n} n^{-\al_1/2}
    \le d_4^{\sqrt n} n^{-\delta \sqrt n/2} .
\]
Thus, in both cases, 
\[
    \frac{f^{\nu/\al} f^\al}{f^\nu}
    \le d_5^{\sqrt n}n^{-d_6\sqrt n}
\]
for positive constants $d_5,d_6$.
Using the same reasoning as in~\eqref{eq:fromR2} and the assumption $t \le g(n) = o(\sqrt n)$, formula~\eqref{eq:expansionR3} yields
\[
    |R_3|
    \le n^t\cdot t\cdot C\cdot d_5^{\sqrt n}n^{-d_6\sqrt n}
    \le C\cdot e^{g(n)\ln n + \ln g(n) + d_7\sqrt n-d_6\sqrt n\ln n}
\]
so that, as $n\to\infty$,
\begin{equation}\label{eq:R3}
   |R_3| 
   \le e^{-d_6\sqrt n\ln n(1+o(1))}
   \to 0 \,.
\end{equation}

It can be easily seen from the bounds \eqref{eq:R1},  \eqref{eq:R2}, and~\eqref{eq:R3}
%on $R_1,R_2,R_3$ 
that $R_2,R_3=o(R_1)$. Thus, the total error term
$R=R_1+R_2+R_3=O(R_1)=O(g(n)f^{-s}(n))=o(f^{-(s-1)}(n))$.
%It remains to observe that the obtained bounds \eqref{eq:R1},  \eqref{eq:R2}, \eqref{eq:R3}
%on $R_1,R_2,R_3$ are uniform over $\nu$. This completes the proof of the theorem.
\end{proof}

%\begin{remark}
%    The claim of this theorem holds under weaker assumptions on $\nu$; namely, instead of taking a.e.\ $\nu$ with respect to Plancherel measure, it suffices to assume that  $\nu_1, \nu'_1 \sim c \sqrt n$ for an arbitrary constant~$c$.
%\end{remark}



\section{A gluing lemma}\label{sec:gluing}

The following ``Gluing Lemma'' is essential to the proof of Theorem~\ref{thm:main1}.

\begin{lemma}[Gluing Lemma]\label{l:gluing}
    Let $\la\vdash n$ be the disjoint union $\la=\mu\cup\tau$ of diagrams $\mu\vdash n-k$ and $\tau\vdash k$. Assume that 
    $k=o(n^{1/4})$,
    the smallest part of $\mu$ is larger than the largest part of $\tau$, 
    and the character $\psi^\mu$ tends to be regular strongly and uniformly,  
    in the sense that for any balanced $\al\vdash n-k$
\[
    \langle L_\mu,s_\al\rangle=\frac{f^\al}{z_\mu}(1+R_\al) 
    \quad \text{with} \quad |R_\al| < \omega(n-k) 
    \quad \text{as} \quad n \to \infty,
\]
    for some function $\omega(n)=o(1)$.
    %as $n\to\infty$. 
    Then $\psi^\la$ tends to be regular. Furthermore,
    \[
        \langle L_\la,s_\nu \rangle
        \sim \frac{f^\nu}{z_\mu z_\tau}
        \quad \text{as} \quad n \to \infty,
    \] 
    for balanced (and hence for Plancherel-typical) $\nu\vdash n$. More precisely, 
    \[
        \langle L_\la,s_\nu\rangle=\frac{f^\nu}{z_\mu z_\tau} (1+R_{\la,\nu})
        \quad \text{with} \quad
        |R_{\la,\nu}| \le O(\om(n-k) + k^2n^{-1/2}).
    \]
\end{lemma}

\begin{proof}
Since $\mu$ and $\tau$ have no common part, 
by~\eqref{eq:Hopf} we have
\[
    \langle L_{\mu\cup\tau},s_\nu \rangle
    = \langle L_\mu L_\tau,s_\nu \rangle
    = \sum_{\al \,\vdash n-k,\,\beta \,\vdash k} c^\nu_{\al\beta} \langle L_\mu,s_\al \rangle \langle L_\tau,s_{\beta} \rangle.
\]
The Littlewood--Richardson coefficient $c^\nu_{\al\beta}$ is nonzero only for $\al \subseteq \nu$. Since $\nu$ is balanced, the rows and columns of $\nu$ are $O(\sqrt{n})=O(\sqrt{n-k})$, hence the same is true for $\al \subseteq \nu$, i.e., $\al\vdash n-k$ is balanced.  
By the assumptions of the theorem,
\[
    \langle L_\mu,s_\al \rangle
    = \frac{f^\al}{z_\mu} (1+R_\al)
\]
with $|R_\al| \le \om(n-k) = o(1)$. 
Therefore
\begin{eqnarray*}%\label{cup}
    \langle L_{\mu\cup\tau},s_\nu \rangle
    &=& z_\mu^{-1} \sum_{\al \,\vdash n-k,\,\beta \,\vdash k} c^\nu_{\al\beta} f^\al (1+R_\al) \langle L_\tau,s_{\beta} \rangle \nonumber \\
    &=& (1 + R_1) \cdot z_\mu^{-1} \sum_{\al \,\vdash n-k,\,\beta \,\vdash k} c^\nu_{\al\beta} f^\al \langle L_\tau,s_{\beta} \rangle,
\end{eqnarray*} 
where $|R_1| \le \om(n-k)$. 
% Denote $\om_1(n) := \om(n-k)$.
Recalling that $f^\al=\langle h_1^{n-k},s_\al\rangle$, it follows from 
% \eqref{cup} and~
\eqref{eq:Hopf} that
% \begin{eqnarray}
%     \langle L_{\mu\cup\tau},s_\nu \rangle
%     &=& z_\mu^{-1} \sum_{\al \,\vdash n-k,\,\beta \,\vdash k} c^\nu_{\al\beta}f^\al (1 + R_\al) \langle L_\tau,s_{\beta} \rangle 
%     = z_\mu^{-1} \sum_{\al \,\vdash n-k,\,\beta \,\vdash k} c^\nu_{\al\beta} (1 + R_\al) \langle h_1^{n-k},s_\al \rangle \langle L_\tau,s_{\beta} \rangle \nonumber \\
%     &\le& (1+\om_1(n)) z_\mu^{-1}
%     \sum_{\al \,\vdash n-k,\,\beta \,\vdash k}
%     c^\nu_{\al\beta} \langle h_1^{n-k},s_\al \rangle \langle L_\tau,s_{\beta} \rangle
%     %= (1+\om_1(n)) z_\mu^{-1} \cdot \langle h_1^{n-k}L_\tau,s_\nu \rangle 
%     \nonumber \\
%     &=& (1+\om_1(n)) z_\mu^{-1} \sum_{\beta\,\vdash k} \langle L_\tau,s_\beta \rangle \langle h_1^{n-k},s_{\nu/\beta} \rangle
%     = (1+\om_1(n)) z_\mu^{-1} \sum_{\beta\,\vdash k} \langle L_\tau,s_\beta \rangle f^{\nu/\beta}.
%     \label{eq:beforetau}
% \end{eqnarray}
\begin{eqnarray*}
    % \langle L_{\mu\cup\tau},s_\nu \rangle
    % &=& 
    \langle L_{\mu\cup\tau},s_\nu \rangle
    &=& (1 + R_1) \cdot z_\mu^{-1} \sum_{\al \,\vdash n-k,\,\beta \,\vdash k} c^\nu_{\al\beta}f^\al\langle L_\tau,s_{\beta} \rangle \nonumber \\
    &=& (1 + R_1) \cdot z_\mu^{-1} \sum_{\al \,\vdash n-k,\,\beta \,\vdash k} c^\nu_{\al\beta}\langle h_1^{n-k},s_\al \rangle \langle L_\tau,s_{\beta} \rangle \nonumber \\
    % &\le& (1+\om_1(n)) z_\mu^{-1}
    % \sum_{\al \,\vdash n-k,\,\beta \,\vdash k}
    % c^\nu_{\al\beta} \langle h_1^{n-k},s_\al \rangle \langle L_\tau,s_{\beta} \rangle
    %= (1+\om_1(n)) z_\mu^{-1} \cdot \langle h_1^{n-k}L_\tau,s_\nu \rangle 
    % \nonumber \\
    &=& (1 + R_1) \cdot z_\mu^{-1} \sum_{\beta\,\vdash k} \langle L_\tau,s_\beta \rangle \langle h_1^{n-k},s_{\nu/\beta} \rangle \\
    &=& (1 + R_1) \cdot z_\mu^{-1} \sum_{\beta\,\vdash k} \langle L_\tau,s_\beta \rangle f^{\nu/\beta}.
\end{eqnarray*}
Since $k = o(n^{1/4})$, Lemma~\ref{lem:strong_Vershik_Kerov} yields, for balanced $\nu \vdash n$,
\begin{equation*}%\label{eq:R_nu_beta}
    f^{\nu/\beta} 
    = \frac{f^\nu f^\beta}{k!} (1 + R_{\nu,\beta})    
\end{equation*}
where $|R_{\nu,\beta}| \le \om_2(n,k) = O(k^2n^{-1/2})$.
Therefore
\begin{eqnarray*}
    \langle L_{\mu\cup\tau},s_\nu \rangle
    &=& (1 + R_1) \cdot z_\mu^{-1} \sum_{\beta\,\vdash k} \langle L_\tau,s_\beta \rangle \frac{f^\nu f^\beta}{k!} (1 + R_{\nu,\beta}) \\
    &=& (1 + R_1) (1 + R_2) \cdot z_\mu^{-1} \sum_{\beta\,\vdash k} \langle L_\tau,s_\beta \rangle \frac{f^\nu f^\beta}{k!} ,
\end{eqnarray*}
where $|R_2| \le \om_2(n,k) = O(k^2n^{-1/2})$.
Now 
\begin{equation*}%\label{eq:tau}
    \frac{f^\nu}{k!} \sum_{\beta\,\vdash k} \langle L_\tau,s_\beta \rangle f^\beta
    = \frac{f^\nu}{k!} \langle L_{\tau},h_1^k \rangle
    = \frac{f^\nu}{z_\tau},
\end{equation*}
since 
\[
    \langle L_{\tau},h_1^k\rangle=\dim\psi^\tau=|C_\tau|=\frac{k!}{z_\tau}.
\]
% Therefore
% \begin{equation}\label{eq:tau}
%     \sum_{\beta\vdash k}\langle L_\tau,s_\beta\rangle f^{\nu/\beta}
%     = (1+\om_2(n))\frac{f^\nu}{k!}\sum_{\beta\vdash k}\langle L_\tau,s_\beta\rangle f^\beta
%     = (1+\om_2(n)) \frac{f^\nu}{k!}\langle L_{\tau},h_1^k\rangle
%     = \frac{f^\nu}{z_\tau}(1+\om_2(n)),
% \end{equation}
% since 
% \[
%     \langle L_{\tau},h_1^k\rangle=\dim\psi^\tau=|C_\tau|=\frac{k!}{z_\tau}.
% \]
% Substituting~\eqref{eq:R_nu_beta} into~\eqref{eq:beforetau} and
% Combining~\eqref{cup}, \eqref{eq:R_nu_beta}, 
% \eqref{eq:beforetau} and~\eqref{eq:tau}, we have
We conclude that
%\[
%\langle L_{\mu\cup\tau},s_\nu \rangle = (1+\max_\alpha\{ R_\alpha)(1+ R_{\nu,\tau})\frac{f^\nu}{z_\mu z_\tau}
%\]
%Thus
\[
    %\langle L_{\la},s_\nu \rangle = 
    \langle L_{\mu\cup\tau},s_\nu \rangle
    = (1 + R_1) (1+ R_2) \frac{f^\nu}{z_\mu z_\tau} 
    = (1 + R_3) \frac{f^\nu}{z_\mu z_\tau} ,
\]
where 
\[
    |R_3| \le O(\om(n-k) + k^2 n^{-1/2}) = o(1),
\]
as required.
\end{proof}

\begin{corollary}\label{cor:hooks}
Any sequence of higher Lie characters indexed by 
hook shapes $(n-k,1^k)$ with $k=o(n^{1/4})$ tends to be regular.    
\end{corollary}
\begin{proof}
    Immediately follows from Theorem~\ref{prop:Lie} and Lemma~\ref{l:gluing}.
\end{proof}



\section{The asymptotics of a random higher Lie character}\label{sec:asymp}


In this section we consider the asymptotic behavior of the higher Lie character $\psi^\la$ for a {\em random} partition $\la \vdash n$. 
As explained in Subsection~\ref{sec:main_results},
it is natural to use the probability distribution
\[
    \Prob(\la) 
    = \frac{\dim \psi^\la}{\dim \Reg_n}
    = \frac{|C_\la|}{n!}
    \qquad (\forall\,\la \vdash n),
\]
where $C_\la$ is the conjugacy class consisting of all permutations in $\S_n$ with cycle type $\la$.
% for $\la\vdash n$ to be distributed according to the ``conjugacy class'' measure 
% $$
%     \conj(\la)=\frac{|C_\la|}{n!}.
% $$
Thus, $\la=\la(\si)$ can be described as the cycle structure of a permutation $\si$ chosen uniformly at random from $\S_n$. 
This equivalent description will be used in the statements and proofs in this section. 
On the other hand, the irreducible characters $\chi^\nu$ (or, equivalently, the Schur function $s_\nu$) will be taken, as usual, according to the Plancherel distribution.

%Recall Definition~\ref{def:tends_in_probability} of a sequence of random characters tending in probability to be regular. 


\begin{definition}\label{def:tends_in_probability}
    A random sequence of characters $\rho = (\rho_n)$, with $\rho_n$ a character of $\S_n$, {\em tends in probability to be regular} if the following holds: 
    % defining the random variable $R_n$ by 
    if
    \[
        \langle \rho_n,\chi^{\nu_n} \rangle 
        = \frac{\dim \rho_n}{n!} \cdot f^{\nu_n} \cdot (1 + R_n),
    \]
    then, for every $\ve > 0$,
    \[
        \lim_{n \to \infty} \Prob\,(|R_n| > \ve) = 0
    \]
    for Plancherel-typical $\nu=(\nu_1,\nu_2,\ldots)$. 
\end{definition}


The main result of this section is the following.



%Recall Definition~\ref{def:tends_in_probability}. 

%\todo{An alternative theorem.

\begin{theorem}\label{thm:main1}
    If $\si_n$ is chosen uniformly at random in~$\S_n$ then $\psi^{\la(\si_n)}$ tends in probability to be regular. 
    Furthermore, 
    % defining the random variable $R_n$ by  
    if
    %then with probability at least $1-o(n)$
    \[
        \langle L_{\la(\si_n)},s_{\nu_n} \rangle 
        = \frac{f^{\nu_n}}{z_{\la(\si_n)}} (1+R_n) 
        % \qquad \text{with} \qquad
        % % |R_n| < \ve
        % \Prob\,(|R_n| > \ve) = o(n^{-b}) %\rightarrow 0,
    \]
    where $\si_n$ is chosen uniformly at random in~$\S_n$, then, for every $\ve >0$,
    \[
        \lim_{n \to \infty} \Prob\,(|R_n| > \ve) = 0
    \]
    for balanced (and hence for Plancherel-typical) $\nu=(\nu_1,\nu_2,\ldots)$. 
\end{theorem}

% \todo{
%     NT and RA: Maybe change the wording of this theorem to fit the statement of Theorem~\ref{thm:main}.
% }
The first part of Theorem~\ref{thm:main1} is Theorem~\ref{thm:main} from the Introduction.
In fact, we prove the following more general version.

\begin{theorem}\label{thm:main_general}
    Let $f(n)$ and $g(n)$ be two functions satisfying the following conditions as $n\to\infty$:
    $f$~is (eventually) monotone increasing, %$f(n)\to\infty$, 
    $g(n) / \ln n\to\infty$, $g(n) = O(f(n) / \ln n)$, %$g(n) = o(\sqrt n)$, 
    and $f(n)g(n)=o(n^{1/4})$. Denote, for any $s>0$,
    \[
        \ve_n := \frac1{f^s(n/2)} + \frac{f^2(n)g^2(n)}{\sqrt n},
        \quad 
        \theta_n
        %= \frac{1}{f(n)}+\frac{\ln n}{2g(n)}.
        := \frac{3 \ln n}{g(n)}.
    \]
    If $\si_n$ is chosen uniformly at random in~$\S_n$, and
    \[
        \langle L_{\la(\si_n)},s_{\nu_n} \rangle 
        = \frac{f^{\nu_n}}{z_{\la(\si_n)}} (1+R_n), 
    \]
    then, for some positive constant $C$, 
    \[
        \Prob\,(|R_n| > C \ve_n) <\theta_n
    \]
    for sufficiently large $n$ and for balanced (and hence for Plancherel-typical) $\nu=(\nu_1,\nu_2,\ldots)$. 
\end{theorem}


In order to prove Theorem~\ref{thm:main_general}, 
%{thm:main1} 
we need the following lemma.

\begin{lemma}\label{lem:equalcycles}
    The probability for a random permutation, uniformly distributed on~$\S_n$, to have two cycles of the same length $m$ for some $m > M$  is at most $1/M$.
\end{lemma}
\begin{proof}
    Consider the number of permutations $\si\in\S_n$ that have two cycles of a fixed length $m$. To determine such a permutation, we can choose the $2m$ elements that belong to these two cycles, choose how they are distributed between the cycles and how they are organized inside the cycles, and then choose a~permutation of the remaining $n-2m$ elements. Thus, the number in question is at most
    \[
        \binom{n}{2m} \cdot \frac{1}{2} \binom{2m}{m} ((m-1)!)^2\cdot (n-2m)!
        = \frac{n!}{2m^2}.
    \]
    Therefore, the probability for a random permutation to have two cycles of the same length $m$ is at most~$m^{-2}$. It follows that the probability to have two cycles of arbitrary length $m>M$ is at most
    \[
        \sum_{m=M+1}^{\lfloor n/2\rfloor}\frac 1{m^2}
        \le \int_{M}^{\infty}\frac1{x^2}dx 
        = \frac1{M}.
        \qedhere
    \]
\end{proof}


% \begin{proof}[Proof of Theorem~\ref{thm:main1}.] 
% %\todo{An alternative proof.
% %\begin{proof}[Proof of Theorem~\ref{thm:main1}.] 
% Let $0<\delta'<\frac14$ to be chosen, and let \magenta{$M_n:=n^{\delta'} / (\ln n)^2$}.
% Choose $\si_n$ uniformly at random in $\S_n$, and write its cycle type $\la=\la(\si_n)$ as a disjoint union $\la=\mu\cup\tau$, where $\mu$ is the part of $\la$ consisting of the rows of length $> M_n$ and $\tau$ is the part consisting of the rows of length $\le M_n$.

% By Lemma~\ref{lem:equalcycles}, for sufficiently large $n$ all rows of $\mu$ are distinct with probability at least $1-\frac{1}{M_n}$. Denote by $c(\si)$ the number of cycles of a permutation $\si$.
% It is well known (see, e.g., \cite[(6.5)]{Feller}) that for a random permutation $\si_n\in\S_n$ we have 
% $Ec(\si_n)\sim\ln n$ as $n\to\infty$, where by $E$ we denote the expectation of a random variable. 
% % Take a number $\kappa >3/\ve$. 
% Then, by Markov's inequality, for sufficiently large $n$,
% \begin{equation}\label{eq:Markov2}
%     \operatorname{Prob}(c(\si_n) > (\ln n)^2)
%     \le \frac{Ec(\si_n)}{(\ln n)^2}
%     < \frac{1}{2 \ln n}.
% \end{equation}
% %Finally, it follows from the Vershik--Kerov limit shape theorem~\cite{VershikKerovPlanch} that for Plancherel-typical $\nu$ we have $c_1\sqrt n<\nu_1,\nu_1'<c_2\sqrt n$ where constants $0<c_1<2<c_2$ can be chosen arbitrarily close to $2$.
% Thus, for sufficiently large $n$, with probability at least $1-\frac{1}{M_n}-\frac{1}{2\ln n}$, all the assumptions of Theorem~\ref{thm:distinct} are satisfied
% with $f(n) = n^{\delta'} / (\ln n)^2$ and $g(n) = (\ln n)^2$.
% %Comparing the bounds \eqref{eq:R1},  \eqref{eq:R2}, \eqref{eq:R3} on the error terms $R_1,R_2,R_3$ in the proof of Theorem~\ref{thm:distinct}, it is not difficult to see that for such $f,g$ the leading error term is $R_1=O(g(n)f^{-s}(n))=o(n^{-1/2})$ for $s$ large enough. 
% Then
%    \[
%         \langle L_{\mu},s_\nu \rangle = \frac{f^\nu}{z_{\mu}}(1+R_n)
%         \quad \text{where} \quad
%         |R_n| \le \om(n) = o(n^{-1/2}).
%     \] 
    
% As for $\tau$, all its rows are of length at most $n^{\delta'} / (\ln n)^2$ and, by inequality~\eqref{eq:Markov2}, 
% % with probability $1-\ve/2$  
% there are at most $(\ln n)^2$ of them. So, the total size of $\tau$ is $k = O(n^{\delta'})$. Since $\delta'<\frac14$, the assumptions of the Gluing Lemma (Lemma~\ref{l:gluing}) are satisfied with $\om(n)=o\!\left( n^{-1/2} \right)$, so
% \[
%     \langle L_{\la},s_\nu \rangle 
%     = \frac{f^\nu}{z_\la}(1+R'_n)
% \]
% where, for any $\delta > \delta'$, 
% $|R'_n| = O((n-k)^{-1/2} + n^{-1/2+\delta'}) = o(n^{-1/2+\delta})$.
% Thus, for any fixed $\ve > 0$, we have $|R'_n|<\ve$ for sufficiently large $n$ with probability at least $1-\frac{1}{M_n}-\frac{1}{\ln n} \to 1$, so
% \[
%     \lim_{n \to \infty} \Prob\,(|R'_n| > \ve) = 0 .
%     \qedhere
% \]
% \end{proof}
% }

% \begin{definition}
%         A sequence of characters $\rho = (\rho_n)$, with $\rho_n$ a character of $\S_n$, is {\em essentially regular} if for every $\ve>0$ there exists $N>0$ such that for every $n\ge N$, 
%     \[
%         % |\langle \rho_n,\chi^{\nu_n} \rangle - \frac{\dim \rho_n}{n!} \cdot f^{\nu_n}| < \ve 
%         \langle \rho_n,\chi^{\nu_n} \rangle 
%         = \frac{\dim \rho_n}{n!} \cdot f^{\nu_n} \cdot (1 + R_n)
%         \qquad \text{with} \qquad
%         % |R_n| < \ve
%         \Prob\,(|R_n| < \ve) \ge 1 - \ve,
%     \]
%         % with probability at least $1-\ve$, 
%         for almost every $\nu=(\nu_1,\nu_2,\ldots)$ with respect to the Plancherel measure. 
% \end{definition}

% \begin{proof}
% It suffices, in the proof of Theorem~\ref{thm:main1}, to take $M_n=f(n)$ and use Markov's inequality~\eqref{eq:Markov2} in the form
% \[
%     \operatorname{Prob}(c(\si_n) > g(n))
%     \le \frac{Ec(\si_n)}{g(n)}
%     \sim \frac{\ln n}{g(n)}
%     < \frac{\ln n}{2g(n)},
% \]
% for sufficiently large $n$.
% Then it is easy to see that, under our assumptions on $f(n)$ and $g(n)$, all the assumptions of Theorem~\ref{thm:distinct} and Lemma~\ref{l:gluing} are satisfied and give the estimates we want.
% \end{proof}

\begin{proof}[Proof of Theorem~\ref{thm:main_general}.] 
    % Let $0<\delta'<\frac14$ to be chosen, and let \magenta{$M_n:=n^{\delta'} / (\ln n)^2$}.
    Choose $\si_n$ uniformly at random in $\S_n$, and write its cycle type $\la=\la(\si_n)$ as a disjoint union $\la=\mu\cup\tau$, where $\mu$ is the part of $\la$ consisting of the rows of length $> f(n)$ and $\tau$ is the part consisting of the rows of length $\le f(n)$.
    
    By Lemma~\ref{lem:equalcycles}, for sufficiently large $n$ all rows of $\mu$ are distinct with probability at least $1-\frac{1}{f(n)}$. 
    
    Denote by $c(\si)$ the number of cycles of a permutation $\si$. It is well known (see, e.g., \cite[(6.5)]{Feller}) that for a random permutation $\si_n\in\S_n$ we have 
    $Ec(\si_n)\sim\ln n$ as $n\to\infty$, where by $E$ we denote the expectation of a random variable. 
    Then, by Markov's inequality, for sufficiently large $n$,
    \[
        \operatorname{Prob}(c(\si_n) > g(n))
        \le \frac{Ec(\si_n)}{g(n)}
        < \frac{2 \ln n}{g(n)}.
    \]
    Assume, from now on, that all the rows of $\mu$ are distinct and that $c(\si_n) \le g(n)$.
    For sufficiently large $n$, this happens with probability at least $1 - \frac{1}{f(n)} - \frac{2 \ln n}{g(n)} > 1 - \frac{3 \ln n}{g(n)} = 1-\theta_n$.
    Note that, by assumption, $g(n) \to \infty$ and $f(n) g(n) = o(n^{1/4})$, thus $f(n) = o(n^{1/4})$, hence $g(n) = O(f(n) / \ln n) = o(n^{1/4}) = o(n^{1/2})$.
    Thus all the assumptions of Theorem~\ref{thm:distinct} are satisfied, so that
    \[
        \langle L_{\mu},s_\nu \rangle = \frac{f^\nu}{z_{\mu}}(1+R_n)
        \quad \text{where} \quad
        |R_n| \le \om(n) = o(f(n)^{-q}),
    \] 
    for any fixed (arbitrarily large) $q$.
    
    As for $\tau$, all its rows are of length at most $f(n)$ and
    %, by inequality~\eqref{eq:Markov2}, for sufficiently large $n$, with probability at least $1-\frac{1}{f(n)}-\frac{\ln n}{2 g(n)}$,
    there are at most $g(n)$ of them. 
    The total size of $\tau$ is $k \le f(n) g(n)$, so by assumption $k = o(n^{1/4})$. 
    The assumptions of the Gluing Lemma (Lemma~\ref{l:gluing}) are satisfied with $\om(n) = o\!\left(f(n)^{-q} \right)$, so
    \[
        \langle L_{\la},s_\nu \rangle 
    = \frac{f^\nu}{z_\la}(1+R'_n)
    \]
    where, by monotonicity of $f$,
    \[
        |R'_n| 
        = O(f(n-k)^{-q} + k^2 n^{-1/2})
        = O(f(n/2)^{-q} + f(n)^2 g(n)^2 n^{-1/2}) \le C \ve_n.
    \]
    This completes the proof of the theorem.
%    Thus, for any fixed $\ve > 0$, we have $|R'_n|<\ve$ for sufficiently large $n$ with probability at least $1 - \frac{\ln n}{g(n)}$, so
%    \[
%        \lim_{n \to \infty} \Prob\,(|R'_n| > \ve) = 0 .
%        \qedhere
%    \]
\end{proof}

\begin{proof}[Proof of Theorem~\ref{thm:main1}.] 
It suffices to take any functions $f(n)$ and $g(n)$ satisfying the assumptions of  Theorem~\ref{thm:main1} --- for instance,
$f(n)=\ln^3n$ and $g(n)=\ln^2n$. 
Since
$f(n)\to \infty$ and $f^2(n)g^2(n))=o(\sqrt n)$, we have 
$\ve_n\to 0$. Also, $g(n) / \ln n\to\infty$, thus $\theta_n\to 0$. 
%    Take, in Theorem~\ref{thm:main_general}, $f(n):=n^{\delta'} / (\ln n)^2$ for any $0<\delta'<1/4$ and $g(n)=(\ln n)^2$.
 %   Then $\ve_n=o(n^{-1/2+\delta})$ for any $\delta>0$ and $\theta_n=O(1 / \ln n)$, optimizing the size of the error term. 
\end{proof}

\begin{remark}
    There is a trade-off between the sizes of the error term $\ve_n$ and the error probability~$\theta_n$. Choosing the functions $f(n)$ and $g(n)$ appropriately, one can optimize whatever one wants. 
\end{remark}

For example, to achieve a good error term $\ve_n$, take $f(n) = n^{\delta'}$ for an arbitrarily small $0<\delta'<1/4$ and $g(n) = 3 (\ln n)^2$. 
On the other hand, to get a good error probability $\theta_n$, take  $f(n)=n^{1/8}$ and $g(n)=n^{1/8}/\ln n$. 
Thus Theorem~\ref{thm:main_general} yields the following.

% For instance, taking $f(n):=n^{\delta'} / (\ln n)^2$ for any 
% $0<\delta'<1/4$ and $g(n)=(\ln n)^2$, as in the proof of Theorem~\ref{thm:main1}, we obtain $\ve_n=o(n^{-1/2+\delta})$ for any $\delta>0$ and 
% $\theta_n=O(1 / \ln n)$ (optimizing the size of the error term). 

\begin{corollary}\label{thm:main_cor1}
    Assume that $\si_n$ is chosen uniformly at random in~$\S_n$ and denote
    \[
        \langle L_{\la(\si_n)},s_{\nu_n} \rangle 
        = \frac{f^{\nu_n}}{z_{\la(\si_n)}} (1+R_n). 
    \]
    Then, for arbitrarily small $\delta>0$, 
    \begin{itemize}
        \item[(a)] 
        \[
            \Prob\,(|R_n| > n^{-1/2+\delta}) 
            < 1 / \ln n 
        \]
        %and
        \item[(b)]
        and
        %, \magenta{for some positive constant $C$}, 
        \[
            \Prob\,(|R_n| > 1/\ln n) 
            % < (\ln n)^2n^{-1/8} 
            < n^{-1/8 + \delta}
        \]

    \end{itemize}
    for sufficiently large $n$ and for balanced (and hence for Plancherel-typical) $\nu=(\nu_1,\nu_2,\ldots)$. 
\end{corollary}

% \todo{
% \begin{theorem}\label{thm:main}
% % If $\si_n$ is chosen uniformly at random in~$\S_n$ then
% % $L_{\la(\si_n)}$ is essentially regular. 
% %More precisely, 
%     for every $\ve>0$ there exists $N>0$ such that for every $n\ge N$, if $\si_n$ is chosen uniformly at random in~$\S_n$, then with probability at least $1-o(n^{-b})$, where $b=\frac{1}{2}+\epsilon ??$
%     \[
%         \langle L_{\la(\si_n)},s_{\nu_n} \rangle = \frac{f^{\nu_n}}{z_{\la(\si_n)}}\left(1+o(n^{-1/2+\delta})\right)
%     \]
% %    for sufficiently large $n$,
%     for arbitrary small $\delta>0$,
%     for almost every $\nu=(\nu_1,\nu_2,\ldots)$ with respect to the Plancherel measure. 

 
% %    For every $\ve>0$, with probability $1-\ve$
% %    \[
% %        \langle L_{\la(\si_n)},s_{\nu_n} \rangle = \frac{f^{\nu_n}}{z_{\la(\si_n)}}(1+R)
% %        \quad\text{ where }\quad |R|<\ve
% %    \]
% %    for sufficiently large $n$,
% %    for almost every $\nu=(\nu_1,\nu_2,\ldots)$ with respect to the Plancherel measure.  
% \end{theorem}

% \begin{proof} 
% In the proof of Theorem~\ref{thm:main1} 
% choose $\kappa:=n^{-\epsilon} ??$ and $M_n:=n^{\delta'} / \ln n$ ??.
% \end{proof}
% }

%\todo{YR: Is the following formulation correct?
%\begin{theorem}
%For a random permutation $\sigma_n\in \S_n$, asymptotically almost surely  
%%and a random partition $\nu^{(n)}\vdash n$
%$$
%\psi^{\la(\sigma_n)}\sim_{prop} \Reg_n.
%$$ 
%More precisely, 
%    for every $\ve>0$  %with probability $1-\ve$
%    \[
%    \langle \psi^{\la(\sigma_n)}, \chi^{\nu_n} \rangle = \frac{|C_{\lambda(\sigma_n)}|}{n!}\cdot \dim \chi^{\nu(n)} \cdot (1+R)
%   \quad\text{ where }\quad |R|<\ve
%    \]
%    for sufficiently large $n$, 
%    for $\sigma_n\in \S_n$ with probability $1-\frac{1}{n^{1/4+\ve}}$, 
%    and 
%    almost every $\nu_n\vdash n$ with respect to the Plancherel measure.  	
%\end{theorem}
%}



% \todo{

% \begin{proof}[Old proof.]
% Let $0<\delta'<\frac14$ to be chosen, and let $M_n:=n^{\delta'} / \ln n$.
% Choose $\si_n$ uniformly at random in $\S_n$, and write its cycle type $\la=\la(\si_n)$ as a disjoint union $\la=\mu\cup\tau$, where $\mu$ is the part of $\la$ consisting of the rows of length $> M_n$ and $\tau$ is the part consisting of the rows of length $\le M_n$.

% Fix $\ve>0$. 
% By Lemma~\ref{lem:equalcycles}, for sufficiently large $n$ all rows of $\mu$ are distinct with probability at least $1-\frac{1}{M_n}\ge 1-\ve/2$. Denote by $c(\si)$ the number of cycles of a permutation $\si$.
% It is well known \red{(see, e.g., \cite[(6.5)]{Feller})} that for a random permutation $\si_n\in\S_n$ we have 
% $Ec(\si_n)\sim\ln n$ as $n\to\infty$, where by $E$ we denote the expectation of a random variable. Take a number $\kappa >3/\ve$. Then, by Markov's inequality, for sufficiently large $n$,
% \begin{equation}\label{eq:Markov}
%     \operatorname{Prob}(c(\si_n) > \kappa \ln n)
%     \le \frac{Ec(\si_n)}{\kappa \ln n}
%     <\frac\ve2.    
% \end{equation}
% %Finally, it follows from the Vershik--Kerov limit shape theorem~\cite{VershikKerovPlanch} that for Plancherel-typical $\nu$ we have $c_1\sqrt n<\nu_1,\nu_1'<c_2\sqrt n$ where constants $0<c_1<2<c_2$ can be chosen arbitrarily close to $2$.
% Thus, with probability $1-\ve$, for sufficiently large $n$, all the assumptions of Theorem~\ref{thm:distinct} are satisfied
% with $f(n) = n^{\delta'} / \ln n$ and $g(n) = \kappa \ln n$.
% %Comparing the bounds \eqref{eq:R1},  \eqref{eq:R2}, \eqref{eq:R3} on the error terms $R_1,R_2,R_3$ in the proof of Theorem~\ref{thm:distinct}, it is not difficult to see that for such $f,g$ the leading error term is $R_1=O(g(n)f^{-s}(n))=o(n^{-1/2})$ for $s$ large enough. 
% Thus, as $n\to\infty$, 
%    \[
%         \langle L_{\mu},s_\nu \rangle = \frac{f^\nu}{z_{\mu}}(1+\om(n))
%         \quad \text{where} \quad
%         \om(n) =o (n^{-1/2}).
%     \] 

% As for $\tau$, all its rows are of length at most $n^{\delta'} / \ln n$ and, by inequality~\eqref{eq:Markov}, with probability $1-\ve/2$  there are at most $\kappa\ln n$ of them. So, the total size of $\tau$ is $k=O(n^{\delta'})$. Since $\delta'<\frac14$, the assumptions of the Gluing Lemma (Lemma~\ref{l:gluing}) are satisfied with $\om(n)=o\left( n^{-1/2} \right)$, so
% \[
% \langle L_{\la},s_\nu \rangle = \frac{f^\nu}{z_\la}(1+R)
% \]
% where $R = O((n-k)^{-1/2} + n^{-1/2+\delta'}) = o(n^{-1/2+\delta})$.
% \end{proof}
% }


\section{Variations}\label{sec:variations}

\subsection{Virtual permutations} 

If, considering the behavior of random higher Lie characters, we want to achieve convergence almost surely instead of convergence in probability, we should introduce a probability measure on {\em infinite sequences} of higher Lie characters. A natural way to achieve this is to consider the space of virtual permutations, introduced in~\cite{KOV} in the context of harmonic analysis on the infinite symmetric group.

Let $\pi_{n,n+1}:\S_{n+1} \to \S_{n}$ be the map that, given a permutation $\si\in\S_{n+1}$, deletes the element $n+1$ from its cycle. The  {\it space of virtual permutations} $\X:=\varprojlim S_n$ is the inverse limit of finite symmetric groups $\S_n$ with respect to these maps. In other words, a virtual permutation, i.e., an element of $\X$, is a sequence $\omega =(\si_1,\si_2,\ldots) \in \S_1\times\S_2\times\ldots$ such that $\si_{n+1}$ is obtained from $\si_n$ by either inserting $n+1$ into one of the existing cycles, or adding $n+1$ as a fixed point; we say that such a sequence of permutations is {\it coherent}.
It is easy to see that the sequence of uniform measures on $\S_n$ agrees with the projections $\pi_{n,n+1}$, so the inverse limit of these measures is a probability measure $M$ on $\X$, which is called the {\it Haar measure} on $\X$. By definition, the projection of $M$ to $\S_n$ coincides with the uniform measure.

\begin{proposition}\label{lem:virtual}
    Let $m_n$ be an increasing sequence of positive integers such that $m_n = \Omega(n^{9})$. 
    Then, for almost every (with respect to the Haar measure) virtual permutation $\sigma =(\si_1,\si_2,\ldots) \in \S_1\times\S_2\times\ldots$,  
    \[
        \langle \psi^{\la(\si_{m_n})},\chi^{\nu_{m_n}} \rangle
        \sim \frac{f^{\nu_{m_n}}}{z_{\la(\si_{m_n})}} 
        \quad \text{as} \quad n \to \infty
   \]
   for balanced (and hence for Plancherel-typical) $\nu=(\nu_1,\nu_2,\ldots)$.
\end{proposition}

\begin{proof}
By the definition of the Haar measure on $\X$, for each $n$ the permutation $\si_n$ is uniformly distributed on~$\S_n$. Let
    \[
        \langle \psi^{\la(\si_{m_n})},\chi^{\nu_{m_n}} \rangle 
        = \frac{f^{\nu_{m_n}}}{z_{\la(\si_{m_n})}} (1+R_{m_n}(\sigma)). 
    \]
Fix an arbitrary $\ve>0$ and denote by $A_{m_n}$ the set of all virtual permutations $\si$
for which $|R_{m_n}(\si)|\ge\ve$. 
By Corollary~\ref{thm:main_cor1}(b),  
%Theorem~\ref{thm:main2}, 
for sufficienlty large $n$
we have (for a suitable $\delta > 0$)
    \[
        \Prob\,(A_{m_n}) 
        < {m_n}^{-1/8 + \delta}
        < n^{-10/9}.
    \]
It follows that the series 
$\sum_{n=1}^\infty \Prob\,(A_{m_n})$ converges. Then, by the Borel--Cantelli lemma,
with probability~$1$ only finitely many of the events $A_{m_n}$ occur, which means that, with probability $1$, for all sufficiently large $n$ we have $|R_{m_n}(\si)|<\ve$. Thus, 
$R_{m_n}(\si)\to0$ as $n\to\infty$
for a.e.\ $\si$.
\end{proof}

\begin{conjecture}\label{conj:virtual}
    For almost every virtual permutation $\si =(\si_1,\si_2,\ldots) \in \S_1\times\S_2\times\ldots$ (with respect to the Haar measure) and for almost every $\nu$ (with respect to the Plancherel measure), 
    \[
        \langle \psi^{\la(\si_n)},\chi^{\nu_n} \rangle
        \sim \frac{f^{\nu_n}}{z_{\la(\si_n)}} 
        \quad \text{as} \quad n \to \infty .
    \]
\end{conjecture}

% \iffalse
% \todo{YR: Three possible applications:
% \begin{itemize}
% \item It seems that all results above hold for conjugacy character on conjugacy classes; this is the character induced from trivial character of the centralizers (and to other coherent families of characters induced from linear characters of centralizers).
% \magenta{NT: yes, see Section~\ref{sec:conjugacy}.}

% \item {\bf Q:} If so, does an analogue of Theorem~\ref{thm:main} imply the main results in~\cite{AF} and ~\cite{R97}? \magenta{NT: no.}

% \item {\bf Q:} Does Theorem~\ref{thm:main} imply Theorem~\ref{thm:Regev}, namely main result in \cite{Regev99}? \magenta{NT: no.}
    
% \end{itemize}
% }
% \fi 


\subsection{Conjugacy characters}\label{sec:conjugacy}

Replacing, in the construction of higher Lie characters described in~\eqref{eq:wreathproduct}, the primitive irreducible character 
$\zeta_k$ of $\bbz_k$ by any character of the form $\zeta_k^r$ for $r \in \bbz$,
we obtain the more general character 
\[
    \psi^\la_r:=\Ind_{Z_\la}^{\S_n}\left(\id_{m_1}[\zeta_1^r]\otimes\id_{m_2}[\zeta_2^r]\otimes\ldots\otimes\id_{m_n}[\zeta_n^r]\right),
\]
where $\id_{m_i}[\zeta_i^r]$ is the wreath product of $\zeta_i^r$ by the trivial character $\id_{m_i}$ of $\S_{m_i}$. Clearly, an analog of~\eqref{eq:hlcformula} also holds:
\[
   \psi^\la_r =\Ind_{\S_{m_1}[\S_1]\times\S_{m_2}[\S_2]\times\ldots}^{\S_n}
\left(\id_{m_1}[\psi_r^{(1)}]\otimes\id_{m_2}[\psi_r^{(2)}]\otimes\ldots\right),
\]
where $\S_{m_i}[\S_i]$ is the wreath product of $\S_i$ by $\S_{m_i}$ and $\id_{m_i}[\psi_r^{(i)}]$ is the wreath product of the character $\psi_r^{(i)}$ of $\S_i$ 
(corresponding to the one-row diagram $(i)$)
by the trivial character $\id_{m_i}$ of $\S_{m_i}$.


\begin{theorem}\label{thm:conjugacy}
All 
%results in Sections~\ref{sec:1row}--\ref{sec:asymp}
the previous results in this paper, in particular Theorems~\ref{main:rectangular} and~\ref{thm:main},
are valid for $\psi_r^\la$, for arbitrary $r$, instead of the higher Lie character~$\psi^\la = \psi_1^\la$.
\end{theorem}
\begin{proof}
It is easy to observe that all the proofs in the previous sections use only the following two facts: 
(a) Theorem~\ref{prop:Lie} about the asymptotics of the ordinary Lie character $\psi^{(n)}$ corresponding to the one-row diagram $\la=(n)$; 
(b) the wreath product construction~\eqref{eq:wreathproduct} of a higher Lie character $\psi^\la$, for an arbitrary $\la$, from ordinary Lie characters. By construction, the structure~(b) is the same for $\psi_r^\la$. 
The proof of Theorem~\ref{prop:Lie}, in turn, uses only the Kra\'skiewicz--Weyman formula (Theorem~\ref{thm:KW})
\[
    \langle\psi^{(n)},\chi^\la\rangle
    =\#\{T\in\operatorname{SYT}(\la)\,\colon \operatorname{maj}T\equiv 1 \!\!\!\pmod n\} =: a_{n,1}
\]
    and Swanson's result (Theorem~\ref{thm:Swanson}) about the asymptotics of $a_{n,1}$. However, both results hold in a more general situation. Namely,  according to Kra\'skiewicz--Weyman~\cite[second corollary to Theorem 1]{KraskiewiczWeyman},
\[
    \langle\psi^{(n)}_r,\chi^\la\rangle=
    \#\{T\in\operatorname{SYT}(\la)\,\colon \operatorname{maj}T\equiv r \!\!\!\pmod n\} =: a_{n,r}
\]
holds for any $r$, 
and Swanson's result is already stated for any $r$. Thus, Theorem~\ref{prop:Lie} holds for $\psi_r^\la$ as well, which implies, according to the observation above, that all the previous results in this paper also hold for the character $\psi_r^\la$.
\end{proof}

In particular, the case $r=0$ corresponds to the conjugacy characters of $\S_n$. Namely, consider the linear
% conjugacy
representation $\pi_C$ of $\S_n$ on the group algebra $\CC[\S_n]$ defined by the action of $\S_n$ on itself by conjugation: $\pi_C(g)(h)= ghg^{-1}$.
Clearly, for each $\la\vdash n$, the subspace $H_\la\subset\mathbb C[\S_n]$ spanned by the
permutations with cycle structure $\la$ is invariant under $\pi_C$. The {\em conjugacy character} $\phi^\la$ of $\S_n$ is the character of the restriction
of $\pi_C$ to $H_\la$. It is not difficult to see that
$\phi^\la=\psi_0^\la$, i.e.,
\[
    \phi^\la
    = \Ind_{Z_\la}^{\S_n}\id_{Z_\la}
    = \Ind_{Z_\la}^{\S_n}\left( \id_{m_1}[\iota_1] \otimes \id_{m_2}[\iota_2] \otimes \ldots \otimes \id_{m_n}[\iota_n] \right)
\]
where $\iota_k=\zeta_k^0$ is the trivial character of~$\bbz_k$.

\begin{corollary}\label{cor:conjugacy}
   All the previous results in this paper are valid for the conjugacy character $\phi^\la$ instead of the higher Lie character~$\psi^\la$. 
\end{corollary}



\subsection{Characters arising from symmetric group actions on (co)homology}

Characters close to higher Lie characters appear in natural symmetric group actions on certain homology and cohomology groups.

For instance, Stanley~\cite{Stanley1982} proved that the character $\pi_n$ of the action of $\S_n$ on the top component of the cohomology ring of the partition lattice $\Pi_n$,  
%\magenta{(
which is isomorphic to the cohomology ring of the pure braid group and to that of the Orlik--Solomon algebra of the braid hyperplane arrangement, 
%)%} 
is equal to $\ve_n\psi^{(n)}$, where $\ve_n$ is the sign character of $\S_n$. 


\begin{corollary}
    The sequence of characters $\pi_n$ tends to be regular.    
\end{corollary}
\begin{proof}
    Consider the classical involution $\omega$ on symmetric functions.
    By Theorem~\ref{prop:Lie}, for Plancherel-typical $\nu\vdash n$,
    \[
        \langle \ch\,\pi_n,s_\nu \rangle
        =\langle \omega\Lie_n,s_\nu \rangle
        =\langle \Lie_n,\omega s_\nu \rangle
        =\langle \Lie_n,s_{\nu'} \rangle
        \sim \frac{f^{\nu'}}{n}
        =\frac{f^{\nu}}{n}
        \quad \text{as} \quad n \to \infty. \qedhere
    \]
\end{proof}

%Note that the character $\pi_n$ is  is isomorphic to the cohomology ring of the pure braid group and that of the Orlik--Solomon algebra of the braid hyperplane arrangement. 

More generally, one can show that all the previous proofs in this paper remain valid if in~\eqref{eq:hlcformula} we replace the higher Lie character $\psi^{(i)}$ by $\pi^{(i)}:=\varepsilon_i\psi^{(i)}$, and/or replace some of the trivial characters~$\id_{m_i}$ by~$\varepsilon_{m_i}$. In particular, all our results are valid for the characters
\[
    \Ind_{\S_{m_1}[\S_1]\times\S_{m_2}[\S_2]\times\ldots}^{\S_n}
    \left(\id_{m_1}[\pi^{(1)}]\otimes\varepsilon_{m_2}[\pi^{(2)}]\otimes\ldots
    \otimes\id_{m_{2i-1}}[\pi^{(2i-1)}]\otimes\varepsilon_{m_{2i}}[\pi^{(2i)}]\otimes\ldots\right),  
\]
which appear in the description of the actions of $\S_n$ on the Whitney homology of the partition lattice (see~\cite[Theorem~2.14]{Hersh}) and on the cohomology of the configuration space of points in $\bbr^d$ (see~\cite[Theorem~2.7]{Hersh}).


\section{Further questions}\label{sec:questions}

%Similar characters}
\subsection{Thresholds}

In Theorem~\ref{thm:rectangular}, regarding rectangular diagrams, it is assumed that the height of the rectangle is asymptotically smaller than its width. 
%Otherwise the sequence of rectangles may not tend to be regular, as the following result shows. 
Indeed, there are families of rectangles which do not tend to be regular. Here is an important example. 

\begin{proposition}\label{prop:2xn}
    The sequence of higher Lie characters indexed by rectangular diagrams $\la=(2^k)$ with constant width $m=2$ and height $k\to\infty$ does not tend to be regular.
\end{proposition}
\begin{proof}
    It is easy to see that $\Lie_2=e_2$. By~\eqref{eq:hlcformula}, we thus have 
    \[
        L_{(2^k)} = h_k[\Lie_2] = h_k[e_2].
    \]
    By \cite[Ex.~I.8.6]{Macdonald},
    \[
        h_k[e_2]
        = \sum_{\nu \vdash 2k,\, \text{all columns even}} s_\nu,
    \]
    so $\langle L_{(2^k)},s_\nu\rangle$ is either $0$ or $1$ for every $\nu$, hence \eqref{claim} obviously fails. 
\end{proof}

\begin{problem}
    What is the threshold for the height $k$, as a function of the width $m$, for a sequence of higher Lie characters indexed by rectangular diagrams $(m^k)$ to tend to be regular?
\end{problem}

% \todo{NT:
% Recall Theorem~\ref{thm:rectangular} and Corollary~\ref{cor:hooks}. 

% {\bf Q:} Regarding hooks and rectangles, what is the threshold for $k$ (as a function of $n$) for good behavior? 

% %\begin{conjecture}\label{conj1}
% %The assumption $k=o(m)$ 
% %%(or, which is the same, $k=o(\sqrt n)$) 
% %in Theorem~\ref{thm:rectangular} 
% %is necessary and sufficient.   
% %\end{conjecture}

% }

%Following Corollary~\ref{cor:hooks}, 
An analogous problem can be posed for hook shapes.

\begin{problem}
    What is the threshold for the leg length $k$, as a function of the total size $n$, for a sequence of higher Lie characters indexed by hook shapes $(n-k,1^k)$ to tend to be regular?
\end{problem}

%\todo{YR: 
Corollary~\ref{cor:hooks} provides a lower bound for the threshold. Here is an upper bound.

\begin{proposition}\label{prop:hooks_t}
    For any $\ve > 0$, any sequence of higher Lie characters indexed by hook shapes $(n-k,1^k)$ with $k> (2 + \ve) \sqrt n$ for infinitely many values of $n$ does not tend to be regular.    
\end{proposition}

\begin{proof}
    Noting that $L_{(n-k,1^k)} = L_{(n-k)} L_{(1^k)} = L_{n-k} h_k$, the Pieri  
    %Littlewood -Richardson 
    rule implies that for any $\nu^{(n)} \vdash n$, $\langle L_{(n-k,1^k)}, s_{\nu^{(n)}} \rangle \ne 0 \Longrightarrow \nu^{(n)}_1 \ge k$. 
    %By assumption, $k/2>2 \sqrt n$. 
    Recalling that $\nu^{(n)}_1 \sim 2\sqrt n$ for a Plancherel-typical sequence $\nu$, one deduces that 
    if $k> (2 + \ve) \sqrt n$ for infinitely many values of $n$, then
    %, for any Plancherel-typical sequence $\nu$, 
    $\langle L_{(n-k,1^k)}, s_{\nu^{(n)}} \rangle = 0$ for infinitely many values of $n$.
    Thus $\psi^{(n-k,1^k)}$ does not tend to be regular.
\end{proof}


% It should be noted that any sequence of higher Lie characters indexed by  hook shapes $(n-k,1^k)$ with $k=\Omega(n^{1/2})$ 
% does not tend to be regular. {\bf to be explained??}. 

% Combining Corollary~\ref{cor:hooks} 
% with  Proposition~\ref{prop:hooks_t} 
% implies that the threshold function for hook shapes satisfies $ o(n^{1/4}) < f(n)< 4 \sqrt n$.
%}



\subsection{Conjugacy characters %on conjugacy classes
in other groups}

An old problem in group representation theory is to determine the irreducible modules and their multiplicities in a conjugacy action of a finite group; see, e.g., \cite{Heide} and references therein.

It is well known that the conjugacy action of the symmetric group on itself tends to be regular; see, e.g., \cite{AF, R96}. 
This result is refined in Section~\ref{sec:conjugacy} above: 
by Theorem~\ref{thm:conjugacy}, the conjugacy action of $S_n$ on a random conjugacy class tends to be regular. 
It is natural to ask whether a similar phenomenon holds for finite groups of Lie type.  

Of special interest are the special linear groups.   
%\todo{
For an element $x\in G=SL_n(\FF_q)$, consider the action of $G$ by conjugation on the conjugacy class  of $x$. 
      
\begin{problem} 
    For which sequences of elements does this character tend 
    %asymptotically equal 
    to be regular, 
    %character of $SL_n(\FF_q)$,  
    as $n\rightarrow \infty$ or $q\rightarrow \infty$ ?
\end{problem}

In particular, for a fixed finite field $\FF_q$, does an analogue of Theorem~\ref{thm:main1} hold for a sequence of uniformly random elements $x_n \in SL_n(\FF_q)$ ?


\subsection{Descent set distribution}

Definition~\ref{def:regularity} uses inner products with irreducible characters (equivalently, Schur functions). 
One can ask what happens if these are replaced, for example, by {\em ribbon Schur functions}, indexed by compositions.
%(or, equivalently, subsets of $[n-1]$) and 
Their inner products with higher Lie characters enumerate, by the Gessel--Reutenauer theorem~\cite[Theorem 2.1]{GR93}, permutations with a given descent set and a given cycle structure. 

\begin{problem}
    Do results analogous to the ones presented in this paper hold for ribbon Schur functions?
\end{problem}

%\todo{
%A closely related problem is to obtain  
    %An affirmative solution may apply 
    Asymptotic estimates on the descent set distribution of permutations in random conjugacy classes may follow. 
    %certain/typical/random  conjugacy classes 
    %with a prescribed (typical?) descent set?
%}
% \todo{Definition: descent and cyclic descent set  of a permutation.

%  Gessel-Reutenauer theorem for fiber sizes of Des (cDes) on Schur-positive sets. 

% Special cases: $S_n$ and its conjugacy classes. 
% }

% % \todo{
% % \[
% % \lim_{n\to \infty}\frac{\langle L_{\la^{(n)}}, s_{J_n}\rangle}{\langle \Reg_n, s_{J_n}\rangle} = ??
% % \]

% % }

% \todo{\begin{question}
%     Do analogous results to main results in the paper hold in this context?\\   
%     Can one obtain asymptotic estimates on the distribution of permutations in certain/typical/random  conjugacy classes with a prescribed (typical?) descent set?
% \end{question}

% NT: At the moment, I do not see how to approach this.}


\begin{thebibliography}{}
    
\bibitem{AAER}
    R.\ M.\ Adin, C.\ A.\ Athanasiadis, S.\ Elizalde, and Y.\ Roichman,
    {\em Character formulas and descents for the hyperoctahedral group},
    Adv.\ Appl.\ Math.~{\bf 87} (2017), 128--169.


\bibitem{AF} 
    R.\ M.\ Adin and A.\ Frumkin,
    {\em The conjugacy character of  $S_n$  tends to be regular}, 
    Israel J.\ Math.~{\bf 59} (1987), 234--240.

%\red{
\bibitem{AHR}
    R.\ M.\ Adin, P.\ Heged\"{u}s and Y.\ Roichman, 
    {\em Higher Lie characters and cyclic descent extension on conjugacy classes}, 
    Algebr.\ Comb.~{\bf 6} (2023), 1557--1591.
%}

%\red{\bibitem{AR15} 
%    R.\ M.\ Adin and Y.\ Roichman, 
%    {\em Matrices, characters and descents}, 
%    Linear Algebra Appl.~{\bf 469} (2015), 381--418.
%}  

%\red{
\bibitem{AS_Thrall} 
    C.\ Ahlbach and J.\ P.\ Swanson,
    {\em Cyclic sieving, necklaces, and branching rules related to Thrall's problem}, 
    Electron.\ J.\ Combin.~{\bf 25} (2018), no.\ 4, Paper No. 4.42, 38 pp.  
%}

\bibitem{Almousa} 
    A.\ Almousa, V.\ Reiner and S.\ Sundaram, 
    {\em Koszulity, supersolvability, and Stirling representations}, 
    Ann.\ Represent.\ Theory~{\bf 2} (2025), 173--247.
    %preprint 2024, {\tt arXiv:2404.10858}. 

%\red{
%\bibitem{GilAlon} 
%    G.\ Alon, 
%    {\em Eigenvalues of the Adin--Roichman matrices}, 
%    Linear Algebra Appl.~{\bf 450} (2014), 280--292.
%}

%\red{
\bibitem{Amendola}
    C.\ Am\'endola, F.\ Galuppi, \'A.\ D.\ R.\ Ortiz, P.\  Santarsiero and T.\ Seynnaeve, 
    {\em Decomposing tensor spaces via path signatures}, 
    J.\ Pure Appl.\ Algebra~{\bf 229} (2025), Paper No.~107807, 19 pp.
%}

%\red{
    \bibitem{Armon}
    S.\ Armon and J.\ P.\ Swanson, 
    {\em Super major index and Thrall's problem}, 
    Algebr.\ Comb.~{\bf 8} (2025), 795--815.
%}

% \magenta{
% \bibitem{BKS_AdvApplMath} 
%     S.\ C.\ Billey, M.\ Konvalinka and J.\ P.\ Swanson, 
%     {\em Asymptotic normality of the major index on standard tableaux}, 
%     Adv.\ Appl.\ Math.~{\bf 113} (2020), 101972, 36 pp.
% }

%\magenta{
\bibitem{Billey} 
    S.\ C.\ Billey, M.\ Konvalinka and J.\ P.\ Swanson, 
    {\em Tableau posets and the fake degrees of coinvariant algebras}, 
    Adv.\ Math.~{\bf 371} (2020), 107252, 46 pp.
%}

%\magenta{
\bibitem{Brandt} 
    A.\ Brandt, 
    {\em The free Lie ring and Lie representations of the full linear group}, 
    Trans.\ Amer.\ Math. Soc.~{\bf 56} (1944), 528--536.
%}

% \magenta{
% \bibitem{BKS_SLC}
%     S.\ C.\ Billey, M.\ Konvalinka and J.\ P.\ Swanson, 
%     {\em On the distribution of the major index on standard Young tableaux}, 
%     S\'em.\ Lothar.\ Combin.~{\bf 84B} (2020), Art. 44, 12 pp.
% }

%\red{
\bibitem{Commins}
    P.\ Commins, 
    {\em Invariant theory for the face algebra of the braid arrangement}, 
    Algebr.\ Comb.~{\bf 8} (2025), 421--468.
%}

\bibitem{DesarmenienWachs}
    J.\ D\'esarm\'enien and M.\ Wachs, 
    {\em Descent classes of permutations with a given number of fixed points}, 
    J.\ Combin.\ Theory Ser.\ A~{\bf 64} (1993), 311--328.

\bibitem{DousseFeray}
    J.\ Dousse and V.\ F\'eray,
    {\em Asymptotics for skew standard Young tableaux via bounds for characters},
    Proc.\ Amer.\ Math.\ Soc.~{\bf 147} (2019), 4189--4203.

\bibitem{Feller}
    W.\ Feller,
    An Introduction to Probability Theory and Its Applications, Vol.\ I,
    Wiley, 1950.

%\bibitem{FeraySniady}
%    V.\  F\'eray and P. \ \'Sniady, 
%    {\em Asymptotics of characters of symmetric groups related to Stanley
%    character formula}, 
%    Ann.\ Math.~{\bf173} (2011), 887--906.  

%\bibitem{FominLulov}
%    S.\ V.\ Fomin and N.\ Lulov,
%    {\em On the number of rim hook tableaux},
%    J.\ Math.\ Sci.\ (New York)~{\bf 87}, No.~6 (1997), 4118--4123.

\bibitem{Garsia} 
    A.\ M.\	Garsia, 
    {\em Combinatorics of the free Lie algebra and the symmetric group}, 
    in: Analysis, et cetera, 309--382, Academic Press, Boston, MA, 1990.

\bibitem{GR93}
    I.\ M.\ Gessel and C.\ Reutenauer, 
    {\em Counting permutations with	given cycle structure and descent set}, 
    J.\ Combin.\ Theory Series A~{\bf~64} (1993), 189--215.

\bibitem{GutierrezRosas}
    A.\ Guti\'errez and M.\ Rosas,
    {\em Necessary conditions for the positivity of Littlewood–Richardson and plethystic coefficients},
    Comptes Rendus Math\'ematique~{\bf 361} (2023), 1163--1173.

\bibitem{Heide} 
    G.\ Heide, J.\ Saxl, P.\ H.\ Tiep, and A.\ Zalesski, 
    {\em Conjugacy action, induced representations and the Steinberg square for simple groups of Lie type}, 
    Proc.\ Lond.\ Math.\ Soc.~{\bf  106} (2013), 908--930.
        
\bibitem{Hersh}
    P.\ Hersh and V.\ Reiner, 
    {\em Representation stability for cohomology of configuration spaces in $\bbr^d$}, 
    Int.\ Math.\ Res.\ Notices IMRN, 2016, pp.\ 1--54.

\bibitem{Inglis} 
    N.\ F.\ J.\ Inglis, R.\ W.\ Richardson, J.\ Saxl, 
    {\em An explicit model for the complex representations of $S_n$} 
    Arch.\ Math.\ (Basel)~{\bf 54} (1990), no.~3, 258--259.

\bibitem{KOV}    
   S.\ V.\ Kerov, G.\ I.\ Olshanski, and A.\ M.\ Vershik,
   {\em Harmonic analysis on the infinite symmetric group},
   Comptes  Rend.\ Acad.\ Sci.\ Paris.~{\bf 316} (1993), 773--778.
   
%\bibitem{Kingman}
%   J.\ F.\ C.\ Kingman, 
%   {\em The coalescent}, 
%   Stochastic Process.\ Appl.~{\bf 13} (1982), 235--248.

\bibitem{Klyachko} 
    A.\ A.\ Klyachko, 
    {\em Lie elements in a tensor algebra}, 
    Siberian Math.\ J.~{\bf 15} (1974), 914--921.

\bibitem{KraskiewiczWeyman}
    W.\ Kra\'skiewicz and J.\ Weyman, 
    {\em Algebra of coinvariants and the action of a Coxeter element}, 
    Bayreuth.\ Math.\ Schr.~{\bf 63} (2001), 265--284 (preprint, 1987).

%\red{
\bibitem{Lim}
    K.\ J.\ Lim, 
    {\em Modular idempotents for the descent algebras of type A and higher Lie powers and modules}, 
    %Lim, Kay Jin
    J.\ Algebra~{\bf  628} (2023), 98--162.
%}


\bibitem{Macdonald}
    I.\ G.\ Macdonald,
    Symmetric Functions and Hall Polynomials, 
    Oxford University Press, 1995.

%\red{
\bibitem{Pak}
    I. Pak, 
    {\em What is a combinatorial interpretation?}
    Proc.\ Sympos.\ Pure Math.~{\bf 110}, 
    American Mathematical Society, Providence, RI, 2024, 191--260. 
%}

\bibitem{Regev99}
    A.\ Regev, 
    {\em Asymptotics of degrees of some $\S_n$-sub regular representations}, 
    Isr.\ J.\ Math.~{\bf 113}, 15--28 (1999). 


\bibitem{Reiner-Banff} 
    V.\ Reiner, 
    {\em Thrall's problem and coarsenings}, 
    Banff workshop on positivity in algebraic combinatorics,  
    lecture slides, 2015, \url{https://www-users.cse.umn.edu/~reiner/Talks/ThrallsProblem.pdf}. 

\bibitem{Saliola} 
    V.\ Reiner, F.\ Saliola and V.\ Welker, 
    {\em Spectra of symmetrized shuffling operators}, 
    Mem.\ Amer.\ Math.\ Soc.~{\bf  228} (2014), no. 1072, vi+109 pp. 

\bibitem{ReinerWebb}
    V.\ Reiner and P.\ Webb, 
    {\em The combinatorics of the bar resolution in group cohomology}, 
    J.\ Pure Appl.\ Algebra~{\bf 190} (2004), 291--327.

\bibitem{Reutenauer} 
    C.\ Reutenauer, 
    Free Lie Algebras, 
    London Mathematical Society Monographs Volume 7, New Series, 
    The Clarendon Press, Oxford University Press, New York, 1993.

\bibitem{R96} 
    Y.\ Roichman, 
    {\em Upper bound on the characters of the symmetric groups}, 
    Invent.\  Math.~{\bf 125} (1996), 451--485.
    
\bibitem{R97} 
    Y.\ Roichman, 
    {\em Decomposition of the conjugacy representation of the symmetric groups}, 
    Israel J.\ Math.~{\bf 97} (1997), 305--316.

\bibitem{Schocker}
    M.\ Schocker, 
    {\em Multiplicities of higher Lie characters}, 
    J.\ Aust.\ Math.\ Soc.~{\bf 75} (2003), 9--21.

\bibitem{Schur} 
    I.\ Schur, 
    {\em \"Uber die rationalen Darstellungen der allgemeinen linearen Gruppe}, 
    Sitzungsber.\ Pr.\ Akad.\ Wiss.\ (1927), 58--75.

\bibitem{Spencer}
    J.\ Spencer and L.\ Florescu,
    Asymptopia, 
    Student Mathematical Library, Vol.\ 71, Amer.\ Math.\ Society, 2014.

\bibitem{Stanley1982}
    R.\ P.\ Stanley, 
    {\em Some aspects of groups acting on finite posets}, 
    J.\ Combin.\ Theory Ser.~A~{\bf} 32 (1982), 132--161.


%\bibitem{Stanley_ECI}
%    R.\ P.\ Stanley, 
%    Enumerative Combinatorics, Vol.\ 1, 
%    Cambridge Studies in Adv.\ Math.\ 62, Cambridge Univ.\ Press, Cambridge, 1997.


\bibitem{Stanley_ECII}%{ECII} 
    R.\ P.\ Stanley, 
    Enumerative Combinatorics, Vol.\ 2, 
    Cambridge Studies in Adv.\ Math.\ 62, Cambridge Univ.\ Press, Cambridge, 1999.  

% \red{
% \bibitem{Stanley}
%     R.\ P.\  Stanley, 
%     {\em On the enumeration of skew Young tableaux}, 
%     Adv.\ Appl.\ Math.~{\bf 30} (2003), 283--294.
% }

%\bibitem{Sundaram94} 
%    S.\ Sundaram, 
%    {\em The homology representations of the symmetric group on Cohen--Macaulay subposets of the partition lattice}, 
%    Adv.\ Math.~{\bf 104} (1994), 225--296.    

%\red{
\bibitem{Sundaram94}
    S.\ Sundaram, 
    {\em The homology representations of the symmetric group on Cohen-Macaulay subposets of the partition lattice}, 
    Adv.\ Math.~{bf 104} (1994), 225--296. 
%}

%\red{
\bibitem{Sundaram18}
    S.\ Sundaram, 
    {\em On conjugacy classes of  $S_n$  containing all irreducibles}, 
    Israel J.\ Math.~{\bf 225} (2018), 321--342.
%}

%\red{
%\bibitem{Sundaram21} 
%    S.\ Sundaram, 
%    {\em Prime power variations of higher Lie modules}, 
%    J.\ Combin.\ Theory Ser.\ A~{\bf 184} (2021), Paper No.\ 105512, 25 pp.
%}

\bibitem{Swanson} 
    J.\ P.\ Swanson, 
    {\em On the existence of tableaux with given modular major index}, 
    Algebr.\ Comb.~{\bf 1} (2018), no. 1, 3--21.

\bibitem{Thrall} 
    R.\ Thrall, 
    {\em On symmetrized Kronecker powers and the structure of the free Lie ring}, 
    Amer.\ J.\ Math.\ {\bf 64} (1942), 371--388. 

%\bibitem{Tsilevich}
%   N.\ V.\ Tsilevich,
%   {\em Distribution of cycle lengths of infinite permutations},
%   J.\ Math.\ Sci.\ (N.Y.)~{\bf 87}, No.~6 (1997), 4072--4081.

\bibitem{VershikKerovPlanch}
    A.\ M.\ Vershik and S.\ V.\ Kerov,
    {\em Asymptotics of the Plancherel measure of the symmetric group and the limiting form of Young tableaux},
    Sov.\ Math.\ Dokl.\ {\bf 18} (1977), 527--531.

\bibitem{VershikKerovChar}
    A.\ M.\ Vershik and S.\ V.\ Kerov,
    {\em Asymptotic theory of characters of the symmetric group},
    Funct.\ Anal.\ Appl.\ {\bf 15}, No.~4 (1981), 246--255.

\end{thebibliography}
    

\end{document}


\section{Old Section: $n$-cycles}

\subsection{Preliminaries}

\begin{theorem}[Kra\'skiewicz-Weyman Theorem]\cite[Theorem 8.4]{Garsia}\label{thm:KW}
    For every partition $\nu\vdash n$, the multiplicity $m_{\nu,(n)} := \langle \psi^{(n)},\chi^{\nu} \rangle$ is equal to the cardinality of the set
    \[
        %A_{k,n,1}=
        \{T \in \SYT(\nu):\ \maj(T) \equiv 1 \pmod n\}.
    \]	
\end{theorem}



\begin{theorem}\cite[Theorem 1.8]{Swanson}\label{thm:Swanson}
    For every partition $\nu\vdash n\ge 1$ and all $r$,
    \[
\frac{|\{T\in \SYT(\lambda):\ \maj(T) \equiv r \pmod n\} |}{f^\lambda}\le \frac{2n^{3/2}}{\sqrt {f^\lambda}}.
    \]	
\end{theorem}

For stronger versions see~\cite{BKS_SLC, BKS_AdvApplMath, Billey}. 

\subsection{The result}

Denote the regular character 
of the symmetric group $\S_n$ by $\rho_n$. 
We claim that the higher Lie character $\psi^{(n)}$ is asymptotically equal to 
$\frac{1}{n} \rho_n$.

\begin{proposition}
For a generic partition $\lambda$, the multiplicity of  of the irreducible character  $\chi^\lambda$ 
in $\psi^{(n)}$ is 
$(1+o(1))%\frac{|C_{(n)}}{|\S_n|}
\frac{1}{n}f^\lambda$. 
\end{proposition}

\begin{proof}
    By the Kra\'skiewicz-Weyman Theorem, the multiplicity of $\chi^\lambda$ in $\psi^{(n)}$ is equal to the number of SYT of shape $\lambda$ and 
    $\maj (T) \mod n\equiv 1$.
By~\cite[Theorem 1.8]{Swanson}, 
this number tends to $1/n f^\lambda$ for all partitions $\lambda$ with $f^\lambda > n^5$.
\end{proof}

\subsection{The cycle type $(n-1,1)$}

\begin{proposition}
The higher Lie character $\psi^{(n)}$ is asymptotically equal to 
$\frac{1}{n} \rho_n$. 
\end{proposition}

\begin{proof}
By Frobenius reciprocity together with the branching rule,
\[
\langle \psi^{(n-1,1)},\chi^\lambda\rangle = 
\langle \psi^{(n-1)},\chi^\lambda\downarrow_{\S_{n-1}}\rangle= \sum\limits_{\mu\in \lambda\downarrow} 
\langle \psi^{(n-1)},\chi^\mu\rangle.
\]
Here  the RHS is a sum over all $\mu\vdash n-1$, which are obtained by deleting a corner box from $\lambda$. If $\lambda$ is generic then the partitions which appear in the sum are also generic. Swanson Theorem then implies that the RHS is asymptotically equal to
\[
\sum\limits_{\mu\in \lambda\downarrow} (1+o(1))\frac{1}{n-1} f^\mu= (1+o(1))\frac{1}{n-1} f^\lambda  = (1+o(1)) \frac{|C_{(n-1,1)}|}{n!}f^\lambda.
\]
\todo{To be cleaned!}
\end{proof}



\subsection{Two distinct large cycles}

\begin{proof}
By Frobenius reciprocity,
\[
\langle \psi^{(n-k,k)},\chi^\lambda\rangle = 
\langle \psi^{(n-k)}\psi^{(k)},\chi^\lambda\downarrow_{\S_{n-k}\times S_k}\rangle.
\]
%Now,  by the Littlewood-Richardson rule, 
Recall that 
for every pair of partitions $\mu\vdash n-k$  and $\nu\vdash k$, 
\[
\langle \chi^\lambda\downarrow_{\S_{n-k}\times S_k}, \chi^\mu\chi^\nu\rangle= c^\lambda_{\mu,\nu} 
\]
the Littlewood-Richardon coefficient. Hence, by~\cite[Theorem 1.8]{Swanson},  
\[
\langle \psi^{(n-k,k)},\chi^\lambda\rangle = 
\sum\limits_{\mu\vdash n-k, \nu\vdash k} c^\lambda_{\mu,\nu}\langle \psi^{(n-k)},\chi^\mu\rangle \langle \psi^{(k)},\chi^\nu\rangle\sim \frac {1}{(n-k)(k)}\sum\limits_{\mu\vdash n-k, \nu\vdash k} c^\lambda_{\mu,\nu}  f^\mu f^\nu=\frac {1}{(n-k)(k)} f^\lambda.  
\]
\todo{Note that the contribution of non-generic $\mu$ and $\nu$ in the RHS sum is tiny}

\todo{To be cleaned!}
\end{proof}

%\subsection{Small number of large cycles of distinct lengths}\ \\

%\todo{The proof should be similar in previous subsection}


%, based on recent results of Billey, Konovolinka and Swanson~\cite{Billey} and others.  

%\todo{Add asymptotic approach a la Regev Swanson}

\section{Old Section: Main Conjecture}
%%%%%%%%%%%%%%%%%%%%%%%%%

Denote the higher Lie character of $\S_n$ indexed by $\mu\vdash n$ by $\psi^\mu$.

\begin{conjecture}
    For every sequence of partitions  $\mu^n \vdash n$  with no bounded parts, the higher Lie characters $\psi^{\mu^n}$ are asymptotically norm equal to the character of the regular representation.
\end{conjecture}


%\section{Old Section: The general case ??}

\todo{Conjecture should hold for small number of large cycles of equal length, as well.}

===
\subsection{Hook shapes}

Now let $\la=(n-k,1^k)$ be a hook diagram.

\begin{proposition}\label{prop:hook}
If $k=o(n^{1/4})$, then the sequence of hook diagrams is regular:
$$
\langle L_{(n-k,1^k)},s_\nu\rangle\sim\frac{f^\nu}{(n-k)k!} \text{ as }n\to \infty
$$ 
for typical Plancherel diagram $\nu\vdash n$.
\end{proposition}

\begin{proof}
We have
$$
\psi^{(n-k,1^k)}=\Ind_{\S_{n-k}\times \S_k}(\psi^{(n-k)}\otimes\id_k), 
$$
so
\begin{eqnarray*}
\langle L_{(n-k,1^k)},s_\nu\rangle &=& \langle h_k\Lie_{n-k},s_\nu\rangle
=\sum_{\nu_1\vdash n-k,\,\nu_2\vdash k}c^\nu_{\nu_1\nu_2}\langle\Lie_{n-k},s_{\nu_1}\rangle\langle h_{k},s_{\nu_2}\rangle\\
&=&\sum_{\nu_1\vdash n-k}c^\nu_{\nu_1,(k)}\langle\Lie_{n-k},s_{\nu_1}\rangle
=\sum_{\nu_1\vdash n-k}c^\nu_{\nu_1,(k)}\frac{f^{\nu_1}}{n-k}(1+o(\frac1{n-k})).
\end{eqnarray*}
Further,
\begin{eqnarray*}
\sum_{\nu_1\vdash n-k}c^\nu_{\nu_1,(k)}f^{\nu_1}&=&\sum_{\nu_1\vdash n-k,\,\nu_2\vdash k}c^\nu_{\nu_1\nu_2}\langle h_1^{n-k},s_{\nu_1}\rangle\langle h_{k},s_{\nu_2}\rangle=\langle h_kh_1^{n-k},s_\nu\rangle\\
&=&\langle h_1^{n-k},s_{\nu/(k)}\rangle=f^{s_{\nu/(k)}}.
\end{eqnarray*}
Thus, it suffices to prove that 
\begin{equation}\label{planch}
\frac{f^{\nu/(k)}}{f^\nu}\sim\frac1{k!}
\end{equation}
for Plancherel-a.e.\ partition $\nu$. Let $\Pl$ be the Plancherel measure on partitions: $\Pl(\la)=\frac{(f^{\la})^2}{k!}$ for $\la\vdash k$.
If $k$ is fixed, then it follows from~\cite{VershikKerovChar} that
$$
\frac{f^{\nu/(k)}}{f^\nu}\sim\frac{\Pl((k))}{f^{(k)}}=\frac1{k!},
$$
so \eqref{planch} holds. 

Now let $k=o(n^{1/4})$. Stanley~\cite{Stanley} proved the following formula for $f^{\nu/\tau}$ where $\tau\vdash k$:
$$
\frac{k!f^{\nu/\tau}}{f^\nu f^\tau}=\sum_{\si\in\S_k}\frac{\chi^\nu(\si)}{f^\nu}\frac{\chi^\tau(\si)}{f^\tau}.
$$
In our case $\tau=(k)$, so $f^\tau=1$ and $\chi^\tau(\si)=1$ for all $\si$, whence
$$
\frac{k!f^{\nu/(k)}}{f^\nu}=\sum_{\si\in\S_k}\frac{\chi^\nu(\si)}{f^\nu}.
$$
To estimate the right-hand side, we use the following estimate due to F\'eray and \'Sniady~\cite{FeraySniady}: there exists a constant $a>1$ such that for every $\nu\vdash n$ and every $\si\in\S_k$
$$
\left|\frac{\chi^\nu(\si)}{f^\nu}\right|\le\left[a\max\left(\frac{\nu_1}{n},\frac{\nu'_1}{n},\frac{n-c(\si)}{n}\right)
\right]^{n-c(\si)},
$$
where $c(\si)$ is number of cycles of $\si$.
Note that $\frac{\nu_1}{n},\frac{\nu'_1}n\sim\frac2{\sqrt n}$ and $\frac{n-c(\si)}{n}\le \frac kn=o(\frac1{\sqrt n})$, so
 the right-hand side above is $\le\left( \frac{b}{n^{1/2}}\right)^{n-c(\si)}$. Therefore, since $n-c(e)=0$,
\begin{eqnarray*}
\left|\sum_{\si\in\S_k}\frac{\chi^\nu(\si)}{f^\nu}-1\right|\le
\sum_{\si\in\S_k}\left( \frac{b}{n^{1/2}}\right)^{n-c(\si)}-1
=\sum_{i=0}^k c(k,i)\left( \frac{b}{n^{1/2}}\right)^{n-i}-1\\
=\prod_{j=1}^{k-1}\bigg(1+\frac{jb}{\sqrt n}\bigg)-1\sim
\sum_{j=1}^{k-1}\ln\bigg(1+\frac{jb}{\sqrt n}\bigg)=
\frac{bk(k-1)}{2\sqrt n}(1+o(1))\to 0,
\end{eqnarray*}
where $c(k,i)$ are unsigned Stirling numbers of the first kind and we have used the well-known formula for the generating function of these numbers (see \cite[Lemma~I.3.3]{Stanley}).
 Finally, we obtain
$$
\frac{k!f^{\nu/(k)}}{f^\nu}=1+o(1),
$$
which completes the proof.
\end{proof}

\begin{remark}
    Observe that if $k$ grows faster than $\sqrt n$, then $(k)\not\subset\nu$ for Plancherel-a.e.~$\nu$, so $f^{\nu/(k)}=0$ and the conjecture fails. The case $k=o(n^\al)$ for $\frac14\le\al<\frac12$ is not clear. 
    See also Remark~\ref{rem:counterexample}.
\end{remark}

%\begin{conjecture}
%The assumption $k=o(n^{1/3})$  in Proposition~\ref{prop:hook} is necessary and sufficient. 
%\end{conjecture}

\subsection{Attaching a bounded diagram to a ``good'' infinite part}

  Note  that if the row lengths of diagrams $\la$ and $\mu$ are disjoint, then $z_{\la\cup\mu}=z_\la z_\mu$ where 
$\la\cup\mu$ is the union of  $\la$ and $\mu$. 

Let us say that a sequence of diagrams $\mu\vdash n$ is \textit{nice} as $n\to\infty$ if 
$\langle L_\mu,s_\nu\rangle\sim\frac{1}{z_\mu}f^\nu$ for all $\nu\vdash n$ with $\nu_1,\nu_1'=O(\sqrt n)$. According to Remark~\ref{rem:nice} and the proof of Theorem~\ref{thm:blocks}, diagrams from Theorem~\ref{thm:main} and Theorem~\ref{thm:blocks} are nice.



\begin{proposition}\label{prop:finitepart}
Let $\la$ and $\mu$ be Young diagrams such that {\rm(a)} $\la\vdash k$ where $k$~is fixed;
{\rm(b)} $\mu\vdash n-k$ is nice; {\rm (c)}  $\mu_i>k$ for all $i$.  Then
$$
\langle L_{\la\cup\mu},s_\nu\rangle\sim\frac{1}{z_{\la\cup\mu}}f^\nu
$$
for Plancherel-a.e. $\nu\vdash n$.
\end{proposition}

\begin{proof}
We have
\begin{equation} \label{cup}
\langle L_{\la\cup\mu},s_\nu\rangle=\langle L_{\la}L_{\mu},s_\nu\rangle=\sum_{\al\vdash k,\,\beta\vdash n-k}c^\nu_{\al\beta}\langle L_\la,s_\al\rangle\langle L_\mu,s_{\beta}\rangle.
\end{equation}
For Plancherel-a.e. $\nu\vdash n$, its rows and columns are $O(\sqrt{n})=O(\sqrt{n-k})$, hence the same is true for $\beta\subset\nu$. Thus, 
$$
\langle L_\mu,s_{\beta}\rangle=z_\mu^{-1}\langle h_1^{n-k},s_\beta\rangle (1+o(1)).
$$
Therefore, the right-hand side of~\eqref{cup} equals
\begin{eqnarray*}
(1+o(1))z_\mu^{-1}\sum_{\al\vdash k,\,\beta\vdash n-k}c^\nu_{\al\beta}\langle L_\la,s_\al\rangle\langle h_1^{n-k},s_{\beta}\rangle\\
=
(1+o(1))z_\mu^{-1}\sum_{\al\vdash k}\langle L_\la,s_\al\rangle\sum_{\beta\vdash n-k}c^\nu_{\al\beta}\langle h_1^{n-k},s_{\beta}\rangle
\end{eqnarray*}
For the last sum, we have
$$
\sum_{\beta\vdash n-k}c^\nu_{\al\beta}\langle h_1^{n-k},s_{\beta}\rangle=
\langle h_1^{n-k},\sum_{\beta\vdash n-k}c^\nu_{\al\beta}s_{\beta}\rangle=\langle h_1^{n-k},s_{\nu/\al}\rangle=f^{\nu/\al}.
$$
But for Plancherel-a.e.\ $\nu$ we have
\begin{equation}\label{skewdim}
\frac{f^{\nu/\al}}{f^\nu} =\frac{f^\al}{k!}(1+o(1)),
\end{equation}
so  we obtain
\begin{eqnarray*}
(f^\nu)^{-1}\sum_{\al\vdash k}\langle L_\la,s_\al\rangle\sum_{\beta\vdash n-k}c^\nu_{\al\beta}\langle h_1^{n-k},s_{\beta}\rangle
=(1+o(1))\sum_{\al\vdash k}\langle L_\la,s_\al\rangle\frac{f^\al}{k!}\\
=\frac1{k!}\langle L_{\la},\sum_{\al\vdash k} f^\al s_\al\rangle(1+o(1))
=\frac1{k!}\langle L_{\la},h_1^k\rangle(1+o(1))=\frac1{z_\la}(1+o(1)),
\end{eqnarray*}
since the dimension of $\psi^\la$ is equal to $|C_\la|=\frac{k!}{z_\la}$.
Therefore,
\begin{eqnarray*}
\frac{\langle L_{\la\cup\mu},s_\nu\rangle}{f^\nu}=\frac1{z_\mu}(1+o(1))\frac1{z_\la}(1+o(1))\sim\frac1{z_\la z_\mu}=\frac1{z_{\la\cup \mu}},
\end{eqnarray*}
as required.
\end{proof}

\begin{remark}\label{rem:counterexample}
     It is clear from the proof that in order for the claim of Proposition~\ref{thm:finitepart} to hold, we need the asymptotics~\eqref{skewdim} for dimensions of skew Young diagrams, where $\nu$ is Plancherel-distributed. For $k=o(n)$ these dimensions were studied in~\cite{DousseFeray}, but only in the case where the rows and columns of $\al$ are $O(\sqrt k)$. In particular, it follows from \cite[Corollary~2]{DousseFeray} that under this assumption on $\la$ the claim is true for 
$k=o(n^{1/3})$.

On the other hand, it follows from Table~1 in the same paper that for $k=n^{1/3}$ there is a counterexample to~\eqref{skewdim}, so the claim of Proposition~\ref{thm:finitepart} fails for $k=n^{1/3}$.
\end{remark}



\begin{conjecture}
Proposition~\ref{thm:finitepart} holds for $k=o(n^{1/3})$.  
\end{conjecture}
===

\subsection{The disjoint union of a linear diagram and a sublinear one}

Note that if the row lengths of diagrams $\la$ and $\mu$ are disjoint, then $z_{\la\cup\mu}=z_\la z_\mu$ where $\la\cup\mu$ is the union of $\la$ and $\mu$. 

%Let us say that a sequence of diagrams $\mu\vdash n$ is \textit{nice} as $n\to\infty$ if 
%$\langle L_\mu,s_\nu\rangle\sim\frac{1}{z_\mu}f^\nu$ for all $\nu\vdash n$ with $\nu_1,\nu_1'=O(\sqrt n)$. According to Remark~\ref{rem:nice} and the proof of Theorem~\ref{thm:blocks}, diagrams from Theorem~\ref{thm:rectangular} and Theorem~\ref{thm:blocks} are nice.



\begin{theorem}\label{thm:finitepart}
Let $\la\vdash k$ and $\mu\vdash n-k$ be Young diagrams such that 
{\rm(a)} $k=o(n^{1/4})$; {\rm(b)} $\mu$ is linear; {\rm (c)}~$\mu_i>k$ for all $i$.  Then the sequence of diagrams $\la\cup\mu$ is regular:
$$
\langle L_{\la\cup\mu},s_\nu\rangle\sim\frac{1}{z_{\la\cup\mu}}f^\nu
$$
for  typical  $\nu\vdash n$.
\end{theorem}

\begin{proof}
It follows from condition~(c) and formula~\eqref{eq:hlcformula} that $L_{\la\cup\mu}=L_\la L_\mu$.
Then by~\eqref{eq:product} we have
\begin{equation} \label{cup}
\langle L_{\la\cup\mu},s_\nu\rangle=\sum_{\al\vdash k,\,\beta\vdash n-k}c^\nu_{\al\beta}\langle L_\la,s_\al\rangle\langle L_\mu,s_{\beta}\rangle.
\end{equation}
The Littlewood--Richardson coefficient $c^\nu_{\al\beta}$ is nonzero only for $\beta\subset\nu$. Since $\nu$ is typical,
the rows and columns of $\nu$ are $O(\sqrt{n})=O(\sqrt{n-k})$, hence the same is true for $\beta\subset\nu$, i.e., $\beta$ is typical. Since $\mu$ is linear, by
Corollary~\ref{cor:blocks} we have
$$
\langle L_\mu,s_{\beta}\rangle=z_\mu^{-1}f^\beta(1+o(1))=z_\mu^{-1}\langle h_1^{n-k},s_\beta\rangle (1+o(1)).
$$
Thus, the right-hand side of~\eqref{cup} equals
\[
(1+o(1))z_\mu^{-1}\sum_{\al\vdash k,\,\beta\vdash n-k}c^\nu_{\al\beta}\langle L_\la,s_\al\rangle\langle h_1^{n-k},s_{\beta}\rangle.
\]
But the last sum is equal to
\[
\langle L_\la h_1^{n-k},s_\nu\rangle
=\sum_{\al\vdash k} \langle L_\la,s_\al\rangle\langle h_1^{n-k},s_{\nu/\al}\rangle=\sum_{\al\vdash k}\langle L_\la,s_\al\rangle f^{\nu/\al}.
\]

For a moment assume that $k$ is fixed and $\nu$ is Plancherel-typical.
Then, by the Vershik--Kerov theorem~\cite{VershikKerovChar}, 
\begin{equation}\label{skewdim}
\frac{f^{\nu/\al}}{f^\nu} =\frac{f^\al}{k!}(1+o(1)),
\end{equation}
so  we obtain
\[
(f^\nu)^{-1}\sum_{\al\vdash k}\langle L_\la,s_\al\rangle f^{\nu/\al}
=(1+o(1))\sum_{\al\vdash k}\langle L_\la,s_\al\rangle\frac{f^\al}{k!}
=\frac1{k!}\langle L_{\la},h_1^k\rangle(1+o(1))=\frac1{z_\la}(1+o(1)),
\]
since $\langle L_{\la},h_1^k\rangle=\dim\psi^\la=|C_\la|=\frac{k!}{z_\la}$.
Therefore, 
\begin{eqnarray*}
\frac{\langle L_{\la\cup\mu},s_\nu\rangle}{f^\nu}=\frac1{z_\mu}(1+o(1))\frac1{z_\la}(1+o(1))\sim\frac1{z_\la z_\mu}=\frac1{z_{\la\cup \mu}},
\end{eqnarray*}
as required.

Now assume that $k$ is not fixed, $k=o(n^{1/4})$, and $\nu$ is typical in the sense of Definition~\ref{def:typical}. It is clear from the above that it suffices to show the asymptotics~\eqref{skewdim} of dimensions of skew Young diagrams. 
Stanley~\cite{Stanley} proved the following formula:
$$
\frac{k!f^{\nu/\al}}{f^\nu f^\al}=\sum_{g\in\S_k}\frac{\chi^\nu(g)}{f^\nu}\frac{\chi^\al(g)}{f^\al}=:A_{\nu,\al},
$$
where $\chi^\nu(g)$ is the value of $\chi^\nu$ at the permutation $g\in\S_k$ regarded as a permutation from $\S_n\supset\S_k$.
We want to show that $A_{\nu,\al}\to1$ as $n\to\infty$ for every $\al\vdash k$. 
Note that %$\frac{\chi^\nu(e)}{f^\nu}\frac{\chi^\al(e)}{f^\al}=1$ and
$|\frac{\chi^\al(g)}{f^\al}|\le1$ for all $g\in\S_k$. Then
$$
\left|A_{\nu,\al}-1\right|\le\sum_{g\in\S_k\,g\ne e}\left|\frac{\chi^\nu(g)}{f^\nu}\right|.
$$
To estimate this sum, we use the following result due to F\'eray and \'Sniady~\cite{FeraySniady}: there exists a constant $a>1$ such that for every $\nu\vdash n$ and every $g\in\S_n$
$$
\left|\frac{\chi^\nu(g)}{f^\nu}\right|\le\left[a\max\left(\frac{\nu_1}{n},\frac{\nu'_1}{n},\frac{n-c(g)}{n}\right)
\right]^{n-c(g)},
$$
where $c(g)$ is number of cycles of $g\in\S_n$.
Note that in our case $\frac{\nu_1}{n},\frac{\nu'_1}n=O(\frac1{\sqrt n})$. Further, denoting by $c_k(g)$ the number of cycles of $g$ {\em as a permutation from $\S_k$}, we have $c(g)=c_k(g)+n-k$, so $n-c(g)=k-c_k(g)\le k$,
that is, $\frac{n-c(g)}{n}\le \frac kn=o(\frac1{\sqrt n})$, so
 the right-hand side above is $\le\left( \frac{b}{\sqrt n}\right)^{k-c_k(g)}$. 
For each $i=1,\ldots, k$, the number of $g\in\S_k$ such that $c_k(g)=i$ is the unsigned Stirling numbers of the first kind $c(k,i)$. Also, $k-c_k(e)=0$ and $e$ is the only permutation in $\S_k$ with $k$ cycles. Therefore,
$$
\sum_{g\in\S_k\,g\ne e}\left|\frac{\chi^\nu(g)}{f^\nu}\right|\le
\sum_{{g\in\S_k\,g\ne e}}\left( \frac{b}{n^{1/2}}\right)^{k-c_k(g)}
=\sum_{i=0}^k c(k,i)\left( \frac{b}{n^{1/2}}\right)^{k-i}-1
$$
By \cite[Lemma~I.3.4]{Stanley}), the generating function of Stirling numbers is  $\sum\limits_{i=0}^{k}c(k,i)x^i=x(x+1)\ldots(x+k-1)$, therefore $\sum\limits_{i=0}^{k}c(k,i)x^{k-i}=(1+x)\ldots(1+(k-1)x)$, and the right-hand side above does not exceed
$$
\prod_{j=1}^{k-1}\bigg(1+\frac{jb}{\sqrt n}\bigg)-1\sim
\sum_{j=1}^{k-1}\ln\bigg(1+\frac{jb}{\sqrt n}\bigg)=
\frac{bk(k-1)}{2\sqrt n}(1+o(1))\to 0,
$$
since $k=o(n^{1/4})$,
which completes the proof.
\end{proof}

\begin{remark}\label{rem:counterexample}
     %It is clear from the proof that in order for the claim of Theorem~\ref{thm:finitepart} to hold, we need the asymptotics~\eqref{skewdim} for dimensions of skew Young diagrams, where $\nu$ is Plancherel-distributed. For $k=o(n)$ these dimensions were studied in~\cite{DousseFeray}, but only in the case where the rows and columns of $\al$ are $O(\sqrt k)$. In particular, it follows from \cite[Corollary~2]{DousseFeray} that under this assumption on $\la$ the claim is true for $k=o(n^{1/3})$.

%Observe that if $k$ grows faster than $\sqrt n$, then $(k)\not\subset\nu$ for Plancherel-a.e.~$\nu$, so $f^{\nu/(k)}=0$ and the conjecture fails.

It follows from \cite[Table 1]{DousseFeray}  that for $k=n^{1/3}$ there is a counterexample to~\eqref{skewdim}, so the claim of Theorem~\ref{thm:finitepart} fails for $k=n^{1/3}$.
\end{remark}

\begin{conjecture}
Theorem~\ref{thm:finitepart} holds for $k=o(n^{1/3})$.  
\end{conjecture}

%%%%%%%%%%%%%%%%%%%%%%%%%%%%%%%%%%%%%%%%%%%%%%%%%%%%%%%%%%%%%%%



\subsection{Several rectangular blocks of different lengths}

Now we extend Theorem~\ref{thm:rectangular} to the case of several rectangular blocks.

\begin{theorem}\label{thm:blocks}
Let $\la=(m_1^{k_1} \cdots m_\ell^{k_\ell}) \vdash n=\sum k_im_i$, where $m_1 > \dots > m_\ell$ are distinct and
{\rm(a)} $k_i=o(m_i)$; 
{\rm(b)} $k_im_i\asymp n$; 
{\rm(c)} $\ell = o(n)$. 
Then this sequence of diagrams is regular:
$$
    \langle L_{\la},s_\nu \rangle \sim \frac{f^\nu}{z_\la} \qquad \text{as }n\to \infty
$$ 
for  typical  $\nu\vdash n$.
\end{theorem}


\begin{proof}
It follows from \eqref{eq:hlcformula} that
$$
L_\la=\prod_{i=1}^\ell L_{(m_i^{k_i})}.
$$
We will use the following formula: for arbitrary symmetric functions $f,g$,
\begin{equation}\label{eq:product}
    \langle fg,s_\nu \rangle
    = \sum_{\la\subset\nu} \langle s_\la, f \rangle \langle s_{\la/\mu}, g \rangle 
    =\sum_{\la,\mu\subset\nu} c^\nu_{\la\mu} \langle s_\la, f \rangle \langle s_\mu, g \rangle,
\end{equation}
and its obvious generalization to any number of factors.



Then
\[
    \langle L_\la,s_\nu \rangle
    = \langle \prod_{i=1}^\ell L_{(m_i^{k_i})}, s_\nu \rangle
    = \sum_{\nu_i \vdash k_im_i} c^\nu_{\nu_1,\ldots,\nu_\ell} \prod_{i=1}^\ell \langle L_{(m_i^{k_i})},s_{\nu_i} \rangle,
\]
    where $c_{\nu_1,\ldots,\nu_r}^\nu$ is a ``multi-Littlewood--Richardson coefficient'' equal to the coefficient of $s_\nu$ in the product
    $s_{\nu_1} \cdots s_{\nu_r}$. For typical  $\nu\vdash n$, the rows and columns of $\nu$ are $O(\sqrt n)=O(\sqrt{k_im_i})$ in view of assumption~(b). Since $c_{\nu_1,\ldots,\nu_r}^\nu\ne0$ only if $\nu_j\subset \nu$, the rows and columns of $\nu_i\vdash k_im_i$ are also $O(\sqrt{k_im_i})$.

Then, by Theorem~\ref{thm:rectangular}, for any $\al>0$
$$
\langle L_{(m_i^{k_i})},s_{\nu_i}\rangle=\frac{f^{\nu_i}}{m_i^{k_i}k_i!}(1+R_{\nu_i})\qquad\text{where}\qquad
|R_{\nu_i}|\le m_i^{-\al}
$$
for sufficiently large $m_i$. It follows from assumptions~(a) and~(b) that 
$k_i=o(\sqrt {n})$ and that if $n$ is large enough, then {\em all the $m_i$'s are large enough}.

\todo{NT:
The claim that all the $m_i$'s are large enough is false.

It suffices to assume $\ell = O(\ln n)$.
}

\todo{YR: 
I believe that we do not need all $m_i$-s to be proportional to $n$,  
and that the gap could be resolved. 

%If I understand correctly, for all parts but $\log \log n$ parts,  $m_i \ge \log n$. Letting $\alpha>1$ the sum of the errors on these parts tend to zero (since there are $\sim \log n$ parts).
 
%For the smallest parts, I think that their sizes grow like "a geometric sequence", 
%thus the sum of the errors is bounded. Furthermore, summing the errors over all but finite smallest parts, the sum could be close to 1 as we wish. 
%Is this correct?\\

%\smallskip

Explanation: I believe that the following lemma holds.

\begin{lemma}
for a random permutation, 
the expected size of the $i$-th cycle (when the cycles are ordered by length) is $m_i \sim n  e^{-i}$, as long as this size tends to infinity. 
\end{lemma}

Here is a heuristic:  The expected number of $k$-cycles  is $1/k$, 
hence 
%for any $m$, 
the expected number of cycles of length %$m<k\le e^i m$ is $i$. 
$e^{-i}n\le k\le n$ is $\sim \int_{ne^{-i}}^n {1\over x} dx= i$. 

Another heuristic: Construct the random permutation by choosing uniformly a letter $j$, then $\pi(j)$ then $\pi^2(j)$, etc., one can verify that the expected length of the resulting cycle is $\sim n/e$; continue recursively.

If the lemma is correct, then the sum of the errors is 
\[
\sum m_i^{-\alpha} \sim n^{-\alpha} \sum_{i=1}^{\log n} e^{i\alpha}  \sim 
n^{-\alpha} \int_{1}^{\log n} e^{\alpha x} dx \sim 
\alpha^{-1}. 
\]
This argument should be cleaned up by, 
first, taking in account in the lemma concentration, not just expectation (for this [Arratia-Tavare 92] could be helpful; see 
see also John Baez homepage, random permutations - Part 7). %https://math.ucr.edu/home/baez/permutations/permutations_7.html). 
Second,  in the displayed sum of errors,  
one should take in account 
the part sizes which are bounded.  

}

Thus, 
$m_i^{-1}\asymp k_in^{-1}=o(n^{-1/2})$.
Therefore, 
$$
\prod_{i=1}^\ell(1+R_{\nu_i})-1=\prod_{i=1}^\ell(1+o(n^{-\al/2}))-1=O(\ell n^{-\al/2})
=o(n^{-\al/2-1})
$$
and
$$
\prod\langle L_{(m_i^{k_i})},s_{\nu_i}\rangle=\frac{1}{\prod m_i^{k_i}k_i!}\left(1+o(n^{-\al/2-1})\right)\prod f^{\nu_i}.
$$
Then 
$$
\langle L_{\la},s_\nu\rangle=\frac{1}{\prod m_i^{k_i}k_i!}\left(1+o(n^{-\al/2-1})\right)
\sum_{\nu_i\vdash k_im_i}c^\nu_{\nu_1,\ldots\nu_\ell}\prod f^{\nu_i}.
$$
But
$$
\sum_{\nu_i\vdash k_im_i}c^\nu_{\nu_1,\ldots\nu_\ell}\prod f^{\nu_i}=\sum_{\nu_i\vdash k_im_i}c^\nu_{\nu_1,\ldots\nu_\ell}\prod\langle h_1^{k_im_i},s_{\nu_i}\rangle
=\langle h_1^{\sum k_im_i},s_\nu\rangle=f^\nu,
$$
and the theorem follows.

\todo{YR: Another suggestion for a more refined count, which may fix the gap: 

First, fix $0<\epsilon << 1$. A partition $\nu\vdash n$ is called {\em 
$\epsilon$-good} if 
 with $\nu_1+\nu'_1\le \epsilon n$. Otherwise it is {\em $\epsilon$-bad}.

Notice that for an $\epsilon$-good partition $\nu\vdash n$,  $f^\nu \ge c^n$ for sufficiently large $n$ and constant $c$ which depends on $\epsilon$. 

Hence, by Theorem 3.6 and proof of Proposition 4.1, the error term for such partitions is exponentially decreasing with $n$. 

Next (assume at the moment that all in $\lambda$ parts are distinct)
\[
\langle L_{\la},s_\nu\rangle=\frac{1}{\prod m_i}
\sum_{\nu_i\vdash m_i}c^\nu_{\nu_1,\ldots\nu_\ell}\prod \left(1+Error(\nu_i)\right)f^{\nu_i}.
\]
Thus, for each series of partitions in the sum $\nu_1\vdash m_1,\dots,\nu_\ell\vdash m_\ell$
\[
\prod \left(1+Error(\nu_i)\right)f^{\nu_i}=\prod_i \left(1+c^{-m_i}\right)f^{\nu_i} \prod_{\nu_i\text{ is $\epsilon$-bad}} \left(1+m_i^{-1}\right)f^{\nu_i}.  
\]
It follows that 
\[
\sum_{\nu_i\vdash m_i}c^\nu_{\nu_1,\ldots\nu_\ell}\prod \left(1+Error(\nu_i)\right)f^{\nu_i}
\]
\begin{equation}\label{eq:1111}
\le\sum_{\forall i\, \nu_i\vdash m_i\atop \nu_i \text{ is $\epsilon$-good}}c^\nu_{\nu_1,\ldots\nu_\ell}\prod (1+c^{-m_i}) f^{\nu_i} + \sum_{\exists i\, \nu_i\vdash m_i\atop \nu_i\text{ is $\epsilon$-bad} 
}c^\nu_{\nu_1,\ldots\nu_\ell}\prod (1+m_i^{-1}) f^{\nu_i}. 
\end{equation}
Recall that
\[
\sum_{\nu_i\vdash m_i}c^\nu_{\nu_1,\ldots\nu_\ell}\prod f^{\nu_i}=f^\nu,
\]
Hence, the first sum in \eqref{eq:1111} is  $\le (1+\sum_i c^{- m_i})f^\nu$. 

For the second sum, recall that the summand $c^\nu_{\nu_1,\nu_2} f^{\nu_1}f^{\nu_2}$ is equal to the number of SYT in $\SYT(\nu)$, for which the letters $\le |\nu_1|$ form a sub-tableau of shape $\nu_1$, and the rest of the letters, after operating jdt, form a tableaux of shape $\nu_2$.
Similar phenomenon holds for all series $\nu_1\vdash m_1,\dots,\nu_\ell\vdash m_\ell$. 




The following Observation completes the proof of the upper bound. 

{\bf Observation.} The probability  that the first $\nu_1$ letters in a  SYT of a typical shape shape $\nu$ is  $\epsilon$-bad is exponentillay decreasing with $m_1$.   
By symmetry of LR coefficients, this holds for every $1\le i\le \ell$. 

{\bf Proof of Observation.} First, observe that the shape $\nu_1$ formed by the prefix of 
a SYT of a typical shape $n$ is also typical as long as $m_1=|\nu_1|$ tends to infinity. Now, 
the number of SYT of a typical shape $\nu_1$ is asymptotically  
equal to 
\[
\frac{\sqrt {m_1!}}{e^{\sqrt m_1}}.
\]
The number of SYT of shape $\nu_1$ that are $\epsilon$-bad is 
\[
\le \binom{m_1}{\epsilon m_1} \binom{\epsilon m_1}{\epsilon m_1/2} {\sqrt {(1-\epsilon)m_1!}\over e^{\sqrt{(1-\epsilon m_1}} }. 
\]
The quotient is exponentially decreasing with $m_1$. One concludes that the second sum\\ $\le  \cdot f^\nu \prod\limits_i e^{\sqrt m_i}\gamma^{-m_i}(1+m_i^{-1})$. The proof of lower bound is similar. 
}

\end{proof}

\begin{definition}
We will say that a sequence of diagrams $\la$ is {\em linear} if all rows of $\la$ grow linearly with $n$ and there are $o(n)$ of them.
\end{definition}

\begin{corollary}\label{cor:blocks}
A linear sequence of diagrams $\la$  is regular.
\end{corollary}

==========
\section{General shapes with linear growth}\label{sec:general}

\todo{NT: Put here everything that we want to say about Theorem~\ref{thm:blocks}.
}

In the following theorem we actually consider a sequence $\la^{(1)}, \la^{(2)}, \dots$ of diagrams, but we write only a single $\la$ for readability reasons.

\begin{theorem}\label{thm:blocks}
Fix $C, \delta, c_1, c_2, \ldots > 0$, and let $\la=(m_1^{j_1} \cdots m_p^{j_p}) \vdash n=\sum j_i m_i$, where $m_1 > \dots > m_p$. 
%Denote by $\ell$ the number of parts of $\la$, i.e., $\ell=j_1+\ldots+j_p$. 
As $n\to\infty$, assume that
\begin{itemize}
    \item[\rm(a)] 
    $j_i m_i \ge c_i n$ for each  $i$,

    \item[\rm(b)] 
    $j_i \le C m_i^{1 - \delta}$ for all $i$, and

    \item[\rm(c)] 
    for every fixed $M$, the total size of the subdiagram of $\la$ consisting of all rows of length at most $M$ is $o(n^{1/4})$.
    %{\rm(c)} $\ell = o(n^{1/4})$. 

\end{itemize}
Then this sequence of diagrams is regular:
$$
    \langle L_{\la},s_\nu \rangle \sim \frac{f^\nu}{z_\la} \qquad \text{as }n\to \infty
$$ 
for  typical  $\nu\vdash n$.
\end{theorem}

\magenta{
\begin{corollary}\label{cor:hooks}
The sequence of hook shapes $(n-k,1^k)$ with $k=o(n^{1/4})$ is regular.    
\end{corollary}
}

\subsection{Lemma on skew dimensions}

For the proof of Theorem~\ref{thm:blocks} we need a strengthening of the following result of Vershik and Kerov~\cite{VershikKerovChar} about dimensions of skew Young diagrams: if $k$ is fixed and $\nu$ is Plancherel-typical, then for every $\beta\vdash k$
\begin{equation}\label{eq:skewdim}
    \frac{f^{\nu/\beta}}{f^\nu} =\frac{f^\beta}{k!}(1+o(1))
\end{equation}
as $n = |\nu| \to \infty$. However, we need~\eqref{eq:skewdim} in the situation where $k$ is not fixed but grows not very fast and $\nu$ is not necessarily Plancherel-typical but typical in the sense of our Definition~\ref{def:typical}. 
There are many papers dealing with similar situations, and
here we closely follow the approach of Dousse and F\'eray~\cite{DousseFeray}. However, we cannot use their results directly, because they assume that both $\nu$ and $\beta$ are typical (in their terminology, ``balanced''), while in our case only $\nu$ is typical, but $\beta$ is arbitrary.


\begin{lemma}\label{lem:strong_Vershik_Kerov}
    If $k = o(n^{1/4})$, $\nu \vdash n$ is typical, and $\beta \vdash k$ is arbitrary, then~\eqref{eq:skewdim} holds
    uniformly on $\beta$ as $n \to \infty$.
\end{lemma}

\begin{proof}%[Proof of Lemma~\ref{t:strong_Vershik_Kerov}]
    Stanley~\cite[Theorem 3.1]{Stanley} proved the following formula:
    \[
        f^{\nu/\beta}
        = \frac{1}{k!} \sum_{g\in\S_k} \chi^\nu(g) \chi^\beta(g) 
        \qquad (\forall\, k \le n,\, \nu \vdash n,\, \beta \vdash k),
    \]
    where $\chi^\nu(g)$ is the value of $\chi^\nu$ at the permutation $g\in\S_k$ regarded as a permutation in $\S_n$ fixing $k+1, \dots, n$.
    Denote
    \begin{equation}\label{eq:Stanley}
        A_{\nu,\beta} :=        \frac{k!f^{\nu/\beta}}{f^\nu f^\beta}
        = \sum_{g\in\S_k} \frac{\chi^\nu(g)}{f^\nu} \frac{\chi^\beta(g)}{f^\beta} .
    \end{equation}
    We want to show that $A_{\nu,\beta}\to1$ as $n \to \infty$, uniformly on $\beta$. 
    Indeed, $|\frac{\chi^\beta(g)}{f^\beta}| \le 1$ for all $g \in \S_k$, and therefore
    \begin{equation}\label{eq:relaxed}
        \left| A_{\nu,\beta}-1 \right|
        \le \sum_{e \ne g \in \S_k} \left| \frac{\chi^\nu(g)}{f^\nu} \right|
        \qquad (\forall\, \beta \vdash k).
    \end{equation}
    To estimate this sum, we use the following result due to F\'eray and \'Sniady~\cite[Theorem 1]{FeraySniady}: 
    there exists a constant $a > 0$ such that for every $\nu\vdash n$ and every $g\in\S_n$
    \[
        \left| \frac{\chi^\nu(g)}{f^\nu} \right|
        \le \left[ a \cdot \max \left\{ \frac{\nu_1}{n},\frac{\nu'_1}{n},\frac{n-c_n(g)}{n}\right\} \right]^{n-c_n(g)},
    \]
    where $c_n(g)$ is the number of cycles of $g\in\S_n$.
    Note that in our case $\frac{\nu_1}{n},\frac{\nu'_1}n=O(\frac1{\sqrt n})$. Further, denoting by $c_k(g)$ the number of cycles of $g$ {\em as a permutation in $\S_k$}, we have $c_n(g) = c_k(g)+n-k$.
    It follows that 
    %$n-c_n(g)=k-c_k(g) \le k$, that is, 
    $\frac{n-c_n(g)}{n} = \frac{k - c_k(g)}{n} \le \frac kn=o(\frac1{\sqrt n})$, so that
    \[
        \left| \frac{\chi^\nu(g)}{f^\nu} \right|
        \le\left( \frac{b}{\sqrt n}\right)^{k-c_k(g)}
    \]
    for some other constant $b>0$. 
    For each $1 \le i \le k$, the number of $g\in\S_k$ with $c_k(g)=i$ is the unsigned Stirling numbers of the first kind $c(k,i)$. 
    Therefore,
    \[
        \sum_{e \ne g \in \S_k} \left| \frac{\chi^\nu(g)}{f^\nu} \right|
        \le \sum_{e \ne g \in \S_k} \left( \frac{b}{\sqrt n} \right)^{k-c_k(g)}
        = \sum_{i=1}^k c(k,i) \left( \frac{b}{\sqrt n} \right)^{k-i}-1.
    \]
    By \cite[Proposition~1.3.4]{Stanley_ECI}, a generating function for these Stirling numbers is  
    \[
        \sum_{i=1}^{k} c(k,i)x^i
        = x(x+1)\ldots(x+k-1)
    \]
    and therefore 
    \[
        \sum_{i=1}^{k}c(k,i)x^{k-i}
        = (1+x)\ldots(1+(k-1)x).
    \]
    Thus
    \[
        \sum_{e \ne g \in \S_k} \left| \frac{\chi^\nu(g)}{f^\nu} \right|
        \le \prod_{j=1}^{k-1} \left( 1+\frac{jb}{\sqrt n} \right) - 1 .
    \]
    Observe that 
    \[
       \sum_{j=1}^{k-1} \ln \left( 1+\frac{jb}{\sqrt n}\right)
        \le \sum_{j=1}^{k-1} \frac{jb}{\sqrt n}
        =\frac{bk(k-1)}{2\sqrt n} 
        \to 0
    \]
    as $n \to \infty$, since $k=o(n^{1/4})$. Then
    \[
        \sum_{e \ne g \in \S_k} \left| \frac{\chi^\nu(g)}{f^\nu} \right|
        %\le \prod_{j=1}^{k-1} \left( 1+\frac{jb}{\sqrt n} \right) - 1
        \le \exp \left( \sum_{j=1}^{k-1} \ln \left( 1+\frac{jb}{\sqrt n} \right) \right) - 1
        \sim \sum_{j=1}^{k-1} \ln \left( 1+\frac{jb}{\sqrt n}\right) 
        \to 0.
    \]
    In view of~\eqref{eq:Stanley} and~\eqref{eq:relaxed},
    this completes the proof.
\end{proof}


\subsection{Proof of Theorem~\ref{thm:blocks}}

%\begin{proof}[Proof of Theorem~\ref{thm:blocks}] 
The idea of the proof is to break the shape $\la$ into the disjoint union $\la=\mu\cup\tau$
of two shapes, $\mu$ consisting of long rows and $\tau$ consisting of short rows, and then use different techniques to deal with each of them. Namely, for the diagram $\mu$ we apply estimates from Theorem~\ref{thm:rectangular}, while for the diagram $\tau$ we proceed differently, relying on Lemma~\ref{lem:strong_Vershik_Kerov}.

\todo{
    More precisely, fix $C, c, \delta >0$. For any $\al > 0$ there exists $M > 0$, depending only on $\alpha$, $C$, $c$ and $\delta$,
    such that
    \begin{equation}\label{rectangularerror1}
        \langle L_{(m^k)},s_\nu \rangle
        = \frac{f^\nu}{m^kk!} (1+R_\nu)
        \quad \text{ with } \quad |R_\nu|\le m^{-\al}
    \end{equation}
    for any $m \ge M$, $k \le c m^{1 - \delta}$ and $\nu \vdash km$ satisfying $\nu_1, \nu'_1 \le C \sqrt{km}$.
}
Fix $\ve > 0$. By assumptions~(a) and~(b), each rectangular shape $(m_i^{j_i})$ satisfies the assumptions of
Theorem~\ref{thm:rectangular}, so for each $s>0$ there exists a positive constant $M$ such that if $m_i > M$ then, for (typical?) $\alpha_i \vdash m_i j_i$ satisfying $(\alpha_i)_1, (\alpha_i)'_1 \le D_i \sqrt{m_i j_i}$, 
\begin{equation}\label{eq:fromrectangular}
        \langle L_{(m_i^{j_i})},s_{\alpha_i} \rangle
        = \frac{f^{\alpha_i}}{m_i^{j_i}j_i!} (1+R_{m_i})
        \quad \text{ where } |R_{m_i}|\le m_i^{-s}.
\end{equation}
Fix $s$ (to be chosen later depending on $\varepsilon$), take the corresponding $M$, and 
let $\mu = (m_i^{j_i})$ with $m_i > M$ and
$\tau = (m_i^{j_i})$ with $m_i \le M$. Then $\la=\mu\cup\tau$ and  $L_\la=L_{\mu\cup\tau}=L_\mu L_\tau$; besides, $z_{\mu\cup\tau}=z_\mu z_\tau$. Denote $k=|\tau|$ and observe that $k = o(n^{1/4})$ by assumption~(c).

We will use the following formula: for arbitrary symmetric functions $f$ and $g$,
\begin{equation}\label{eq:product}
    \langle fg,s_\nu \rangle
    = \sum_{\alpha\subset\nu} \langle f,s_\alpha\rangle \langle s_{\nu/\alpha}, g \rangle 
    =\sum_{\al,\beta\subset\nu} c^\nu_{\al\beta} \langle f,s_\al \rangle \langle g,s_\beta \rangle.
\end{equation}
Then 
\begin{equation} \label{cup}
    \langle L_{\mu\cup\tau},s_\nu \rangle
    = \sum_{\al\vdash n-k,\,\beta\vdash k}c^\nu_{\al\beta}\langle L_\mu,s_\al\rangle\langle L_\tau,s_{\beta}\rangle.
\end{equation}

Let us first deal with $L_\mu$. 
The Littlewood--Richardson coefficient $c^\nu_{\al\beta}$ is nonzero only for $\al\subset\nu$. Since $\nu$ is typical,
the rows and columns of $\nu$ are $O(\sqrt{n})=O(\sqrt{n-k})$, hence the same is true for $\al\subset\nu$, i.e., $\al\vdash n-k$ is typical. Denote by $t$ the number of different blocks $(m_i^{j_i})$ in $L_\mu$. We have 
$$
L_\mu=\prod_{i=1}^t L_{(m_i^{j_i})}.
$$
By an obvious generalization of~\eqref{eq:product} to an arbitrary number of factors,
\begin{equation}\label{eq:decompmu}
    \langle L_\mu,s_\al \rangle
    = \langle \prod_{i=1}^t L_{(m_i^{j_i})}, s_\al \rangle
    = \sum_{\substack{\al_1,\ldots,\al_t\\ \al_i \vdash j_im_i}} c^\al_{\al_1,\ldots,\al_t} \prod_{i=1}^t \langle L_{(m_i^{j_i})},s_{\al_i} \rangle,
\end{equation}
where $c^\al_{\al_1,\ldots,\al_t}$ is a ``multi-Littlewood--Richardson coefficient'' equal to the coefficient of $s_\al$ in the product $s_{\al_1} \cdots s_{\al_t}$.
Observe that, again, $c^\al_{\al_1,\ldots,\al_t}$ is nonzero only if $\al_i\subset\al$.
We have seen that the rows and columns of $\al$ are $O(\sqrt{n})$, which is the same thing as $O(\sqrt{m_ij_i})$ in view of assumption~(a), and hence the same it true for $\al_i$. Thus, each $\al_i\vdash m_ij_i$ is typical.
Since $m_1,\ldots,m_t> M$,  we can apply~\eqref{eq:fromrectangular} to obtain
\[
    \prod_{i=1}^t\langle L_{(m_i^{j_i})},s_{\al_i}\rangle=\prod_{i=1}^t\frac{f^{\al_i}}{m_i^{j_i}j_i!}\prod_{i=1}^t(1+R_{m_i})
    =\frac1{z_\mu}\prod_{i=1}^tf^{\al_i}\prod_{i=1}^t(1+R_{m_i})
     \qquad\text{where}\qquad |R_{m_i}|\le m_i^{-s}.
\]
For $s>1$ we have 
\[
    \left|\ln\prod_{i=1}^t(1+R_{m_i})\right|
    = \left|\sum_{i=1}^t\ln(1+R_{m_i})\right|\le\sum_{i=1}^t|R_{m_i}|
    \le \sum_{i=1}^t\frac1{{m_i}^s}
    \le \sum_{j=1}^\infty\frac1{j^s}
    \le\int_1^\infty\frac1{x^s}dx
    =\frac{1}{s-1},
\]
hence
\[
    \left|\prod_{m_i>a}(1+R_{m_i})-1\right|
    \le e^{1/(s-1)}-1
    \le \varepsilon
\]
for sufficiently large $s$, which we now choose to satisfy this inequality. Then, in view of~\eqref{eq:decompmu},
\[
    \langle L_\mu,s_\al \rangle
    = (1+R)\frac1{z_\mu}\sum_{\substack{\al_1,\ldots,\al_t\\ \al_i \vdash j_im_i}} c^\al_{\al_1,\ldots,\al_t} \prod_{i=1}^t f^{\al_i}
    =(1+R)\frac{\langle h_1^{n-k},s_\al\rangle}{z_\mu}
    =(1+R)\frac{f^\al}{z_\mu} 
    \qquad\text{where}\qquad |R|<\varepsilon,
\]
since
\[
    \sum_{\substack{\al_1,\ldots,\al_t\\ \al_i \vdash j_im_i}} c^\al_{\al_1,\ldots,\al_t} \prod_{i=1}^t f^{\al_i}
    =\sum_{\substack{\al_1,\ldots,\al_t\\ \al_i \vdash j_im_i}} c^\al_{\al_1,\ldots,\al_t} \prod_{i=1}^t \langle h_1^{j_im_i}, s_{\al_i}\rangle
    =\langle h_1^{n-k},s_\al\rangle,
\]
again by the generalization of~\eqref{eq:product} to an arbitrary number of factors.
Thus, by \eqref{cup} and~\eqref{eq:product},
\begin{eqnarray}
    \langle L_{\mu\cup\tau},s_\nu \rangle
    &=& (1+R)z_\mu^{-1} 
    \sum_{\al\vdash n-k,\,\beta\vdash k}
    c^\nu_{\al\beta} \langle h_1^{n-k},s_\al \rangle \langle L_\tau,s_{\beta} \rangle
    =(1+R)z_\mu^{-1} \langle h_1^{n-k}L_\tau,s_\nu\rangle \nonumber
     \\
    &=& (1+R)z_\mu^{-1} \sum_{\beta\vdash k} \langle L_\tau,s_\beta \rangle \langle h_1^{n-k},s_{\nu/\beta} \rangle
    = (1+R)z_\mu^{-1}\sum_{\beta\vdash k} \langle L_\tau,s_\beta \rangle f^{\nu/\beta}.
    \label{eq:beforetau}
\end{eqnarray}

Recall that $z_\la=z_\mu z_\tau$, so it suffices to prove that the sum in the right-hand side of~\eqref{eq:beforetau} is equivalent to $f^\nu z_\tau^{-1}$ as $n\to\infty$.
Since $k = o(n^{1/4})$ by assumption~(c), Lemma~\ref{lem:strong_Vershik_Kerov} yields
\[
    f^{\nu/\beta} =\frac{f^\nu f^\beta}{k!}(1+o(1))
\]
as $n = |\nu| \to \infty$ uniformly on $\beta$, so we obtain
\begin{equation}\label{eq:tau}
    \sum_{\beta\vdash k}\langle L_\tau,s_\beta\rangle f^{\nu/\beta}
    = (1+o(1))\frac{f^\nu}{k!}\sum_{\beta\vdash k}\langle L_\tau,s_\beta\rangle f^\beta
    = (1+o(1)) \frac{f^\nu}{k!}\langle L_{\tau},h_1^k\rangle
    = \frac{f^\nu}{z_\tau}(1+o(1)),
\end{equation}
since 
\[
    \langle L_{\tau},h_1^k\rangle=\dim\psi^\tau=|C_\tau|=\frac{k!}{z_\tau}.
\]
Finally, combining~\eqref{eq:beforetau} and~\eqref{eq:tau}, we have
\[
    \langle L_{\la},s_\nu \rangle 
    =\langle L_{\mu\cup\tau},s_\nu \rangle 
    = (1+R)(1+o(1)) \frac{f^\nu}{z_\mu z_\tau}
    = \frac{f^\nu}{z_\la}(1+r),
\qquad\text{where}\qquad |r|<2\varepsilon
\]
for sufficienlty large $n$. This completes the proof of Theorem~\ref{thm:blocks}.
%\end{proof}

%%%%%%%%%%%%%%%%%%%%%%%%%%%%%%%%%%%%%%%%%%%%%%%%%%%%%%%%%%%%%%%%%%%%%%%%%%

\section{Old Section: Asymptotic equality of characters}
%%%%%%%%%%%%%%%%%%%%%%%%%%%%%%%%%%%%%%%%%%%

In this section we introduce and compare three different notions of asymptotic equality of characters of the symmetric group.

%\medskip

Consider two infinite sequences of characters,  
$\{\varphi_n\}_{n=1}^\infty$ and $\{\phi_n\}_{n=1}^\infty$,  
where for every positive integer $n$,  
$\varphi_n$ and $\phi_n$ are $\S_n$-characters. 

The first notion was considered in~\cite{AF}.
It uses the standard inner product of class functions,
\[
    \langle \varphi, \phi \rangle
    := \frac{1}{|\S_n|} \sum_{\pi \in \S_n} 
    \varphi(\pi) \overline{\phi(\pi)}
    = \sum_{\chi \in Irr(\S_n)} 
    \langle \varphi,\chi \rangle \overline{\langle \phi,\chi \rangle}.
\]

\begin{definition}
    $\varphi_n$ and $\phi_n$ are {\em asymptotically norm equal}  if
    \[
        \lim_{n \to \infty}
        \frac{\langle \varphi_n, \phi_n\rangle}{||\varphi_n||\cdot ||\phi_n||}
        = 1. 
    \]    
\end{definition}

Here is an example. 

%\subsection{preliminaries}
\begin{theorem}\cite{AF} 
    The characters of the conjugation representation and the regular representations of $\S_n$ are asymptotically norm equal. 
\end{theorem}

\todo{RA:
Note that if 
\[
    \lim_{n \to \infty} 
    \frac{||\varphi_n||}{||\phi_n||} 
    = 1
\]
and
\[
    \lim_{n \to \infty} 
    \frac{\langle \varphi_n,\phi_n \rangle}{||\varphi_n|| \cdot ||\phi_n||} 
    = 1
\]
then
\[
    \lim_{n \to \infty} 
    \frac{||\varphi_n - \phi_n||^2}{||\varphi_n|| \cdot ||\phi_n||} 
    = 0 .
\]
}

The second notion is based on convergence of the character multiplicity vector. 

\begin{definition}
    Two $\S_n$-characters $\varphi_n$ and $\phi_n$ are {\em asymptotically multiplicity proportional}  
    if, for all %generic (??) 
    irreducible characters $\chi \in Irr(\S_n)\setminus A$,
    \[
        \frac{\langle \varphi_n, \chi \rangle}{\varphi_n(id)} 
        = (1+o(1)) \cdot \frac{\langle \phi_n, \chi \rangle}{\phi_n(id)} ,  
    \] 
    while the rest satisfy 
    \[
    \lim_{n \to \infty}\frac{\sum\limits_{\chi \in A}\langle \varphi_n, \chi\rangle ^2}{\langle \varphi_n,\varphi_n\rangle}= \lim_{n \to \infty}\frac{\sum\limits_{\chi \in A}\langle \phi_n, \chi\rangle ^2}{\langle \phi_n,\phi_n\rangle}=0. 
    \]
\end{definition}

\todo{These two notions can be generalized to any sequence of finite groups.
}

An $\epsilon$-generic partition is a partition 
$\lambda=(\lambda_1,\dots,\lambda_t)\vdash n$ 
with 
%both width and height $\le (1-\epsilon)n$, namely
\[
    \max\{\lambda_1,t\}\le  (1-\epsilon)n.
\]



Here is an example. 

%\subsection{preliminaries}
\begin{theorem}\cite{R97} 
    For generic partitions $\lambda$, the multiplicity of  of the irreducible character $\chi^\lambda$ in the conjugation representation of $\S_n$ is $(1+o(1)) f^\lambda$. 
\end{theorem}

The third notion is based on the fact that the distribution of descent sets on a combinatorial basis of an $\S_n$-representation 
determines the character.

For an $\S_n$-character $\rho$, let $v^\rho$ be a vector indexed by subsets of $[n-1]$, ordered in anti-lexicographic order, such that the entries are defined as follows
\[
    v^\rho_J
    := \sum_{\lambda \vdash n} 
    \langle \rho,\chi^\lambda \rangle \cdot |\{T \in \SYT(\lambda) \,:\, \Des(T)=J\}|
    \qquad (\forall J \subseteq [n-1]).
\]

\begin{definition}
$\varphi_n$ and $\phi_n$ are {\em asymptotically descent equal}  if
\[
    \lim_{n \to \infty}
    \frac{\langle v^{\varphi_n}, v^{\phi_n} \rangle}{||v^{\varphi_n}|| \cdot ||v^{\phi_n}||} 
    = 1. 
\]      
\end{definition}

\begin{conjecture}
    If $\varphi_n$ and $\phi_n$ are asymptotic multiplicity proportional then they are  asymptotic norm equal.
\end{conjecture}


\begin{proof}
First, for any $\S_n$ character $\rho$, 
%$\rho=\sum\limits_{\lambda\vdash n} \langle \rho,\chi^\lambda \rangle \chi^\lambda$, thus, 
by the orthonormality of the irreducible characters, 
\[
    ||\rho||^2 
    = \langle \rho, \rho \rangle
    = \sum_{\lambda \vdash n} \langle \rho,\chi^\lambda \rangle^{2} . 
\]
Also
\[
\langle \varphi_n, \phi_n\rangle = \sum\limits_{\lambda\vdash n} \langle \varphi_n,\chi^\lambda \rangle \langle \phi_n,\chi^\lambda \rangle.
\]
Hence
    \[
    \frac{\langle \varphi_n, \phi_n\rangle}{||\varphi_n||\cdot ||\phi_n||}=  \frac{\sum\limits_{\lambda\vdash n}\langle \varphi_n,\chi^\lambda \rangle \langle \phi_n,\chi^\lambda \rangle}{\left(\sum\limits_{\lambda\vdash n}\langle \varphi_n,\chi^\lambda \rangle\right)^{1/2}\cdot \left(\sum\limits_{\lambda\vdash n}\langle \phi_n,\chi^\lambda \rangle\right)^{1/2}}.
    \]
Denote 
\[
m_{\lambda,\phi_n}:=\langle \varphi_n,\chi^\lambda\rangle.
\]
Then for every $\epsilon>0$ there exists $N$ such that for every 
$n>N$, and generic $\lambda\vdash n$
\[
|\langle \phi_n,\chi^\lambda\rangle - \langle \varphi_n,\chi^\lambda\rangle |< \epsilon \langle \varphi_n,\chi^\lambda\rangle.  
\]
We conclude that for $n>N$
\[  
\frac{(1-\epsilon)^{1/2}\sum\limits_{\lambda\vdash n}  m_{\lambda,\varphi_n} }
{(1+\epsilon)^{1/2}\left(\sum\limits_{\lambda\vdash n}m_{\lambda,\varphi_n}\right)} \le 
    \frac{\langle \varphi_n, \phi_n\rangle}{||\varphi_n||\cdot ||\phi_n||}\le  \frac{(1+\epsilon)^{1/2}\sum\limits_{\lambda\vdash n}  m_{\lambda,\varphi_n} }
{(1-\epsilon)^{1/2}\left(\sum\limits_{\lambda\vdash n}m_{\lambda,\varphi_n}\right)}.
\]
%Proof is completed. 
    
\todo{To be cleaned up, taking in account 
the normalization and  non-generic partitions. }
    
\end{proof}

\begin{conjecture}
 If
 \[
\frac{1}{2^{n/2}}||1-\frac{\langle v^{\varphi_n}, v^{\phi_n}\rangle}{||v^{\varphi_n}||\cdot ||v^{\phi_n}||}|| \longrightarrow 0
\]
then $\varphi_n$ and $\phi_n$ are asymptotically norm equal.
\end{conjecture}

\todo{Proof Idea: 
Combining~\cite[Observation 7.1]{AR15} with~\cite[Theorem 1]{GilAlon} implies asymptotic equality of the vector of character values under the assumption, which implies in turn asymptotic norm equality. 
}

\todo{Question: Does convergence of the vector of fiber sizes of the descent map imply a convergence of the the character multiplicities vector?}

%\todo{Question: Does convergence of the vector of fiber sizes of the descent map imply a convergence of the character multiplicities vector?}

\todo{To be extended, and applied in the paper}

=======================================================

\bigskip\bigskip
=================

However, if we take a sequence of independent uniformly distributed $\si_n\in\S_n$, then the cycle structures of $\si_n$ and $\si_{n+1}$  
may be very different, so there is little hope for a reasonable asymptotic behavior of $\psi^{\la(\si_n)}$. We need such a sequence of permutations to be coherent in the sense that $\la(\si_n)$  and $\la(\si_{n+1})$ should be close to each other. A convenient way to achieve this is to consider the so-called space of virtual permutations, introduced in~\cite{KOV} in the context of harmonic analysis on the infinite symmetric group.

Let $\pi_{n,n+1}:\S_n\to\S_{n+1}$ be the operator that, given a permutation $\si\in\S_{n+1}$, deletes the element $n+1$ from its cycle. The  {\it space of virtual permutations} $\X:=\varprojlim S_n$ is the inverse limit of finite symmetric groups $\S_n$ with respect to these operators. In other words, a virtual permutation, i.e., an element of $\X$, is a sequence $\omega=(\si_1,\si_2,\ldots)\in\S_1\times\S_2\times\ldots$ such that $\si_{n+1}$ is obtained from $\si_n$ by either inserting $n+1$ into one of the existing cycles, or adding $n+1$ as a fixed point (we will say that such a sequence of permutations is {\it coherent}).
Thus,  the cycle structures of $\si_n$ and $\si_{n+1}$ differ very little. It is easy to see that the sequence of uniform measures on $\S_n$ agrees with the projections $\pi_{n,n+1}$, so the inverse limit of these measures is a measure $M$ on $\X$, which is called the {\it Haar measure} on $\X$. By definition, the projection of $M$ to $\S_n$ coincides with the uniform measure.

\begin{theorem}\label{thm:main}
For almost every, with respect to the Haar measure, coherent sequence of permutations $\omega=(\si_1,\si_2,\ldots)\in\S_1\times\S_2\times\ldots$,  
\[
   \langle \psi^{\la(\si_n)},\chi^{\nu}\rangle\sim\frac{f^{\nu}}{z_{\la(\si_n)}} 
   \text{ as }n\to \infty
\]
for typical $\nu$.
In other words, a Haar-typical coherent sequence of diagrams is regular.
\end{theorem}

\begin{proof} %{\it (Sketch; more accurate estimates are needed.)}
Let $c(\si)$ be the number of cycles of a permutation $\si$. 
Also, denote by $l_i(\si)$ the length of the $i$th longest cycle of~$\si$ (with $l_i(\si)=0$ if $c(\si)<i$).
It is known (\cite[Proposition~2]{Tsilevich}, which is essentially an adaptation of Kingman's~\cite{Kingman} result for exchangeable random partitions) that the limits
\[
    x_i=\lim_{n\to\infty}\frac{l_i(\si_n)}{n},\qquad i=1,2,\ldots,
\]
exist for a.e.\ virtual permutation $\om=(\si_1,\si_2,\ldots)$.
%, and the random sequence $X(\om):=(X_1(\om),X_2(\om),\ldots)$ has the Poisson--Dirichlet
%distribution $\PD(1)$ on the infinite simplex $\Sigma:=\{x=(x_1,x_2,\ldots)\colon x_1\ge %x_2\ge\ldots\ge0,\,\sum x_i=1\}$. 
Thus, $l_i(\si_n)\sim x_in$, that is, all rows of $\si_n$ grow linearly with respect to $n$. 
%where, by the properties of $\PD(1)$, almost surely $x_i\ne x_j$. 
This means that $\la(\si_n)=(m_i^{j_i})$ satisfies condition~(a) of Theorem~\ref{thm:blocks}.
Further,   the total number of parts of $\la(\si_n)$  is equal to $c(\si_n)$, and $c(\si_n)\sim\ln n$ as $n\to\infty$ (\cite{Tsilevich}), which ensures that assumptions~(b) and~(c) of the same theorem are also satisfied. Thus, applying Theorem~\ref{thm:blocks} completes the proof.
\end{proof}

Applying Remark~\ref{rem:Planch}, we obtain the following corollary.

\begin{corollary}
For almost every, with respect to the Haar measure, coherent sequence of permutations $\omega=(\si_1,\si_2,\ldots)\in\S_1\times\S_2\times\ldots$ and for almost every, with respect to the Plancherel measure, growing sequence of Young diagrams $\nu=(\nu_1,\nu_2,\ldots)$,  
\[
    \langle \psi^{\la(\si_n)},\chi^{\nu_n}\rangle\sim\frac{f^{\nu_n}}{z_{\la(\si_n)}} \text{ as }n\to \infty.
\] 
In other words, a typical coherent sequence of diagrams is regular.
\end{corollary}

\todo{RA: It seems that this work has several main results:
\begin{enumerate}

    \item 
    The elegant result about virtual permutations, Theorem~\ref{thm:main}, relying on Theorem~\ref{thm:blocks}.
    
    \item 
    The result about rectangular shapes, Theorem~\ref{thm:rectangular}. 

    \item 
    The result about generalized conjugacy characters, namely hLc of hook shapes, Corollary~\ref{cor:hooks}.
    
\end{enumerate}
